\documentclass[10pt, reqno]{amsart}

% --  --  --  --  --  --  --  --  --  --  --  --  --  --  --  --  --  --  --  --  --  --  -- -
% Packages and Encoding
% --  --  --  --  --  --  --  --  --  --  --  --  --  --  --  --  --  --  --  --  --  --  -- -
\usepackage[utf8]{inputenc}
\usepackage[margin=1.2in]{geometry}

\usepackage[T1]{fontenc}
\usepackage[english]{babel}
\usepackage{amssymb, amsmath, amsthm, mathrsfs}
\usepackage[dvipsnames]{xcolor}
\usepackage{amsfonts}%
%\usepackage{natbib}% removed -- amsalpha is not compatible with natbib author-year mode; only plain \cite is used
%\usepackage{datetime}% removed -- it redefined \\today to include the weekday
\usepackage{enumitem}
\usepackage{mathtools}
\usepackage{stmaryrd}
\usepackage{thmtools}
\usepackage{bm}
\usepackage{comment}
\usepackage{hyperref}
\usepackage{amscd}
\usepackage{tikz-cd}
\usetikzlibrary{arrows.meta,positioning,calc}
\usepackage{caption}

\newcommand{\dashbullet}{\raisebox{0.2ex}{\rule{0.8em}{0.4pt}}}

\newlist{subtlelist}{itemize}{1}
\setlist[subtlelist,1]{
  label=\dashbullet,
  labelsep=0.6em,
  leftmargin=2em,
  itemsep=0.25\baselineskip,
  topsep=0.5\baselineskip
}

% Hyperref configuration
\hypersetup{
    colorlinks=true,
    linkcolor=NavyBlue,
    citecolor=NavyBlue,
    urlcolor=NavyBlue
}

% --  --  --  --  --  --  --  --  --  --  --  --  --  --  --  --  --  --  --  --  --  --  -- -
% Typography and Layout
% --  --  --  --  --  --  --  --  --  --  --  --  --  --  --  --  --  --  --  --  --  --  -- -
\linespread{1.15}
\setlength{\parskip}{0.1em}

% --  --  --  --  --  --  --  --  --  --  --  --  --  --  --  --  --  --  --  --  --  --  -- -
% Theorem Environments (AMS Style)
% --  --  --  --  --  --  --  --  --  --  --  --  --  --  --  --  --  --  --  --  --  --  -- -
\theoremstyle{plain}
\newtheorem{theorem}{Theorem}[section]
\newtheorem{proposition}[theorem]{Proposition}
\newtheorem{lemma}[theorem]{Lemma}
\newtheorem{corollary}[theorem]{Corollary}
\newtheorem{claim}[theorem]{Claim}
\newtheorem{example}[theorem]{Example}

\theoremstyle{definition}
\newtheorem{definition}[theorem]{Definition}
\newtheorem{remark}[theorem]{Remark}
\newtheorem{axiom}[theorem]{Structural axiom}
\newtheorem{observation}[theorem]{Observation}
\newtheorem{protocol}[theorem]{Protocol}
\newtheorem{inquiry}[theorem]{Structural inquiry}

\theoremstyle{plain}
\newtheorem{hypothesis}[theorem]{Hypothesis}
\newtheorem{conjecture}[theorem]{Conjecture}

% --  --  --  --  --  --  --  --  --  --  --  --  --  --  --  --  --  --  --  --  --  --  -- -
% Macros and Notation
% --  --  --  --  --  --  --  --  --  --  --  --  --  --  --  --  --  --  --  --  --  --  -- -
\newcommand{\Q}{\mathbb{Q}}
\newcommand{\Z}{\mathbb{Z}}
\newcommand{\R}{\mathbb{R}}
\newcommand{\C}{\mathbb{C}}
\newcommand{\BSD}{\mathsf{BSD}}
\newcommand{\ZF}{\mathsf{ZF}}
\newcommand{\ZFC}{\mathsf{ZFC}}
\newcommand{\Spec}{\operatorname{Spec}}
\newcommand{\Isom}{\operatorname{Isom}}
\newcommand{\Ho}{\operatorname{Ho}}
\newcommand{\Reg}{\operatorname{Reg}}
\newcommand{\MW}{\mathrm{MW}}
\newcommand{\coloneqq}{\mathrel{\mathop:}=}
\newcommand{\defect}{\delta_{\mathrm{mod}}}

% Tate--Shafarevich symbol (reserved for the outlook on Birch--Swinnerton-Dyer)
\DeclareFontFamily{U}{wncy}{}
\DeclareFontShape{U}{wncy}{m}{n}{<->wncyr10}{}
\DeclareSymbolFont{mcy}{U}{wncy}{m}{n}
\DeclareMathSymbol{\Sha}{\mathord}{mcy}{"58}
\newcommand{\Sel}{\operatorname{Sel}}
\DeclareMathOperator{\ord}{ord}

\newcommand{\parahead}[1]{%
  \par\medskip
  \noindent\textit{#1.}\enspace
}

\begin{document}

\title[Intensional semantics for comparison problems]{Intensional semantics for comparison problems in arithmetic geometry}
\date{\today}
\author{R. Laniewski}

\subjclass[2020]{Primary 18N10, 03F35; Secondary 18D20, 03B38, 11G05, 11G40}
\keywords{Lawvere enriched categories, quantitative categories, $(2,1)$-categories, C-systems, identity types, intensional semantics, transport cost, Szpiro conjecture, $abc$ conjecture, Birch and Swinnerton-Dyer conjecture, Tate--Shafarevich group}

\begin{abstract}
This work proposes a semantic environment for arithmetic geometry in which the passage between two presentations of one object carries a measurable weight, rather than dissolving into a transparent identification. Following Lawvere, a category equipped with a cost on its morphisms is enriched over the monoidal poset $([0,\infty],\ge,+)$, so that a height becomes a map controlled by a generalized distance. Following Voevodsky, the identity type of a C-system is a locus of structure open to transport. A labeled quantitative $(2,1)$-category reconciles the two readings, with invertible $2$-morphisms witnessing that two transports differ only by a coherent change of presentation. A comparison between two labeled realizations of one object becomes a transport inequality. If it is asked to descend functorially to an inequality between extensional invariants, then it collapses to a tautology, unless a non-zero load is assigned to the transport itself, at which point the inequality regains its content.

We formulate an intensional Szpiro inequality for elliptic packages, where the discriminant height transports along a morphism with a controlled model defect, and relate a bound of this shape to the $abc$ conjecture.

We then treat the Birch and Swinnerton-Dyer conjecture in its pre-modularity form, so that the rank and refined forms descend to $\mathrm{ACA}_{0}$ without analytic continuation. The refined constant reads as a transport cost in which the regulator, the period, the local Tamagawa data, and the order of the Tate--Shafarevich group $\mathrm{Sha}(E/K)$ are held apart by the descent obstruction. We reduce the finiteness of $\mathrm{Sha}(E/K)$ to one hypothesis.
\end{abstract}

\maketitle

\tableofcontents

\section{Introduction}\label{sec:intro}

A convention governs much of contemporary arithmetic geometry, so familiar that it is rarely stated. Once two number fields, two schemes, or two elliptic curves are known to be isomorphic, and once an isomorphism between them has been chosen, one treats them as the same object of discourse. Heights and volumes are then required to be insensitive to the route by which the identification was reached. This discipline is convenient, and for a great many purposes it is the correct one. It is nonetheless a choice, and it carries a cost. It discards the history of a construction, and with that history a quantity that, in certain arithmetic questions, one would wish to measure.

The present work takes that quantity as its subject. We develop a semantic environment in which the passage from one presentation of an arithmetic object to another carries an explicit weight, and in which a height is permitted to change along such a passage by an amount bounded by that weight. The proposal does not depart from established mathematics. It draws together three lines of thought, each of long standing, and observes that they describe one structure.

The first line is the theory of enriched categories in the form given to it by Lawvere~\cite{Lawvere1973}. A generalized metric space, in his reading, is nothing other than a small category enriched over the monoidal poset $([0,\infty],\ge,+)$, with the triangle inequality appearing as the compositional law and with no requirement of symmetry. A category whose morphisms carry a non-negative cost, subadditive under composition and vanishing on identities, is precisely such an enriched structure. The distance from one object to another becomes the least cost of a transport between them, and a height becomes a non-expansive map controlled by this metric. This is the language in which the weight of a transport is expressed.

The second line is the theory of C-systems and univalent foundations developed by Voevodsky~\cite{voevodsky2023martin}, together with the $\infty$-topos semantics of Riehl~\cite{Riehl2024}. In a C-system the identity type $\mathrm{Id}_A(a,b)$ between two terms is itself a type rather than a bare proposition, its inhabitants witnessing particular paths of identification, and its J-structure regulating how an identification may be transported or contracted. When the system is non-trivial in the homotopical sense, two terms may be identified by paths that carry distinct content. This is the language in which we register that an identification is a structure rather than a redundancy.

The third line is the resource-sensitive type theory of Atkey~\cite{Atkey2018}, whose quantitative categories with families enrich the standard semantics of Martin-L\"of type theory by a semiring of usage annotations. The role of the semiring is played here by the Lawvere quantale $([0,\infty],\ge,+)$, and the usage of a transport is its cost. This is the language in which the metric and the identity type are reconciled, so that transport along a path of zero usage preserves a height exactly, while transport along a path of positive usage incurs a bounded penalty.

These three lines meet inside a single object, which we term a labeled quantitative $(2,1)$-category. Its objects are arithmetic data carrying labels and construction histories. Its $1$-morphisms are transports endowed with a cost. Its $2$-morphisms are invertible identifications between transports, present so that two transports differing only by a coherent change of auxiliary choice are never forced into literal equality, yet silent in the numerical account, since the cost is required to depend on a transport only up to such a $2$-isomorphism. Over this structure sits a forgetful functor to the extensional category of ordinary arithmetic geometry, and the fiber over a fixed object -- its vertical groupoid -- is the place where pure intensionality lives. In the extensional world this groupoid is contractible and every transport is free. In the intensional world it retains positive distances, and these distances are what the framework measures.

The contribution of this device to the study of arithmetic comparison statements is the following. A comparison between two presentations of one underlying object is a transport inequality in the labeled quantitative $(2,1)$-category. We isolate, in a form requiring no special geometry, the exact condition under which such an inequality descends to a non-trivial inequality between extensional invariants. If one demands that the descent be functorial with respect to extensional isomorphisms and transparent to the labels at the end, then the two sides of the comparison are forced to be evaluations of one invariant on one object, and the inequality collapses to the harmless form $H \le H + \Delta$. The content of a comparison statement, when there is any, resides exactly in the refusal to contract the vertical groupoid -- that is, in the assignment of a positive load to the transport. We regard this observation as the categorical reinterpretation of comparison problems named in the title.

The remainder of the work carries this reading to a Diophantine question and then grounds it. We introduce elliptic packages, in which an elliptic curve is presented together with a chosen model, local data, and a label. A morphism of packages transports the discriminant height along a change of Weierstrass coordinates, and the cost of the transport -- the model defect -- bounds the increase of the height. Within this quantitative category we state an intensional Szpiro inequality, in which the retained defect appears as a necessary term, and we describe how a bound of this shape relates to the $abc$ conjecture through the classical reduction of Szpiro~\cite{Szpiro1990} and Masser--Oesterl\'e~\cite{MasserOesterle1985}, once the defect is held small. We then descend the entire account to a pre-set-theoretic arithmetic universe formalizable in the subsystem $\mathsf{ACA}_0$ of second-order arithmetic~\cite{simpson2009subsystems}, where the inequalities become comparisons between finite tuples of integers. A small formalization in the Coq system~\cite{Bertot2004} exhibits the grammar of the framework inside a kernel, keeping the $2$-categorical level explicit while the scalar bounds attach to $1$-morphisms alone.

It remains to delimit what is, and what is not, claimed. We do not assert that the intensional Szpiro inequality is a theorem, and the central Diophantine statement is offered as a hypothesis whose conditional consequences we trace. What we offer is a semantic environment, with a precise categorical shape and a conservative arithmetic shadow, in which comparison statements of a certain kind acquire stable and measurable content. The environment is, we believe, of independent interest, and the same vocabulary is turned, later in this paper, toward the Birch and Swinnerton-Dyer conjecture, where regulators, Tate--Shafarevich data, and periods invite a reading as transport costs over a labeled base. That application occupies Sections~\ref{sec:BSD} and~\ref{sec:TateShafarevich}.

\parahead{Orientation} Throughout, two environments recur, and it is convenient to name them at the outset. Universe~H is the Hilbertian and extensional background, where a single model of $\ZF$ is fixed, isomorphic objects are identified once an isomorphism has been chosen, and heights depend only on isomorphism classes. Universe~A is the intensional and quantitative background, where objects carry labels and construction histories, where identifications carry a cost, and where the logical distance between two isomorphic copies may register as a datum. A third environment, Universe~P, appears later as the pre-set-theoretic arithmetic recipient of the framework, formalizable in $\mathsf{ACA}_0$. The letters serve only as signposts. The mathematics resides in the structures they label.


\section{Two disciplines of identification}\label{sec:identification}

Before any cost can be measured, one must agree on when two presentations are to count as the same object of discourse. Two disciplines are available, and the distance between them is the subject of this work.

\begin{definition}[Static and dynamic identification]\label{def:StaticDynamic}
Let $\mathcal{C}$ be a category of arithmetic objects such as number fields or schemes.

\smallskip
\noindent \textit{Static identification.} Under \emph{static identification}, one works, for Diophantine purposes, with isomorphism classes in $\mathcal{C}$. Once two objects $X$ and $Y$ of $\mathcal{C}$ are known to be isomorphic, they may be interchanged without comment and treated as the same object of discourse.

\smallskip
\noindent \textit{Dynamic identification.} Under \emph{dynamic identification}, the atomic object of discourse comprises both the underlying entity $X$ of $\mathcal{C}$ and a label encoding its contextual origin. Two labeled objects $(X,\ell_X)$ and $(Y,\ell_Y)$ keep their semantic separation within formulas regardless of any structural isomorphism $X \cong Y$ in $\mathcal{C}$, so that intensional distinctions persist as mathematical content.

Static identification reflects the prevailing standpoint of arithmetic geometry. Dynamic identification models a discipline in which distinct labeled copies of arithmetic data are kept apart even where the underlying objects coincide up to isomorphism.
\end{definition}

\begin{remark}[Structural rather than arbitrary identifications]\label{rem:ArbitraryStructural}
The phrase ``removing labels'' is to be understood in a purely structural sense. The background environment $\mathcal{S}$ is a fixed model of set-theoretic semantics for arithmetic geometry. Within $\mathcal{S}$, static identification applies only along isomorphisms preserving the relevant structure. When two number fields, schemes, or elliptic curves are isomorphic, and once an isomorphism has been fixed, one recovers the familiar convention by which these objects are treated as identical for Diophantine purposes.

Outside this scope, no identifications are imposed. In particular, the arguments do not rely on collapsing arbitrary distinct elements of a ring, nor on fusing nested universes in ways that would violate the axioms of $\ZF$. When the text speaks of forgetting labels, it refers to discarding auxiliary names and indices once a structural or protocol-determined identification has been fixed in $\mathcal{S}$.
\end{remark}

\begin{definition}[Standard and labeled signatures]\label{def:Signatures}
Within the background environment $\mathcal{S}$, the \emph{standard signature} refers to the regime in which one works with static identification together with ring coherence as in Axiom~\ref{ax:RingCoherence}. In this regime, isomorphic copies of a number field or scheme are not distinguished once an identification has been fixed, and quantifiers range over underlying objects of a $\ZF$-style universe.

The \emph{labeled signature} refers to the regime in which dynamic identification is maintained. Arithmetic objects carry explicit context labels, and not every isomorphism is available as an identification inside formulas. Ring coherence may then be suspended locally and partly reconstituted by auxiliary devices.

The two signatures share the same informal arithmetic vocabulary. What changes is the discipline that governs when two presentations may be regarded as the same object of discourse.
\end{definition}

\begin{definition}[Protocol-determined data]\label{def:ProtocolDetermined}
Fix a semantic environment $\mathcal{S}$. By a \emph{protocol} we mean an explicit rule of construction, written as finite mathematical instructions, that takes specified input data and outputs specified objects or morphisms.

A datum, for instance a family of local identifications or a gluing map, will be called \emph{protocol-determined in $\mathcal{S}$} when it is produced by a declared protocol from declared input data, and when its standing as the output of the protocol does not rely on an ambient uniqueness statement in the sense of static identification.

In this work, the phrase ``protocol-determined'' is used in places where the word ``canonical'' would invite an extensional reading that is not intended. Inside the extensional environment $\mathcal{S}$, an expression such as ``canonical'' often carries an implicit promise of uniqueness up to a unique isomorphism, together with insensitivity to the route by which an identification is reached. In the labeled signature the standing of a construction comes instead from a declared derivation protocol and from the construction history of the objects being transported, so that the output is singled out by the fact that the protocol produces it, and not by an ambient uniqueness principle that erases the history of transport.
\end{definition}

\begin{remark}[The divergence of the two disciplines]\label{rem:divergence}
A comparison statement, read in the standard signature, identifies its two sides as soon as their underlying objects are isomorphic, so that any apparent discrepancy between them is absorbed into the chosen isomorphism. The same statement, read in the labeled signature, keeps the two sides apart, and any discrepancy between them may survive as a measured quantity. The whole of this work lives in the space between these two readings, and its purpose is to give the second reading a precise categorical form together with a conservative arithmetic shadow.
\end{remark}


\section{Ring coherence and suspended evaluation}\label{sec:ring-coherence}

The informal assumption that underlies much of conventional arithmetic geometry may be described, for present purposes, as a principle of \emph{ring coherence}. One starts with a single ambient ring and expects the arithmetic quantities under discussion to derive from this ring in a compatible way. Heights, log-volumes, and similar expressions are then regarded as maps of ring elements that respect both the additive and the multiplicative structures. This coherence is rarely stated, yet it governs many habitual identifications.

\begin{axiom}[Ring coherence]\label{ax:RingCoherence}
Let $R$ be a commutative ring with identity. Arithmetic invariants that depend simultaneously on the additive and multiplicative structures of $R$ are taken to be defined on the underlying set of $R$ in a manner compatible with both structures. Once an element $x \in R$ has been fixed, all such invariants take values that depend only on $x$ as a ring element, and not on any auxiliary presentation or labeling of $x$.
\end{axiom}

The content of Axiom~\ref{ax:RingCoherence} may be unfolded by separating the symbolic level of arithmetic expressions from the level of their evaluation in a fixed ring, and then imposing a compatibility requirement between the two levels.

\begin{definition}[Arithmetic term category]\label{def:ArithmeticTermCategory}
Let $\mathcal{T}_\mathrm{arith}$ be the small category defined as follows.

\begin{enumerate}
\item[(i)] Objects are finite lists of formal variables
  \[
    X = (x_1,\dots,x_n)
  \]
with $n \geq 0$. The empty list is allowed and represents a nullary context.

\item[(ii)] A morphism
  \[
    \tau \colon (x_1,\dots,x_n) \longrightarrow (y_1,\dots,y_m)
  \]
is an $m$-tuple of formal expressions in the variables $x_1,\dots,x_n$ built from the symbols $0,1$, the binary operations $+$ and $\cdot$, and the unary operation $-$, using finitely many applications of these operations. Composition is defined by substitution of expressions.
\end{enumerate}
\end{definition}

An object of $\mathcal{T}_\mathrm{arith}$ represents a choice of formal parameters, while a morphism represents a vector of arithmetic terms in these parameters. To connect this with a fixed ring, one introduces an evaluation functor.

\begin{definition}[Evaluation functor]\label{def:EvaluationFunctor}
Let $R$ be a commutative ring with identity. The evaluation functor
\[
  E_R \colon \mathcal{T}_\mathrm{arith} \longrightarrow \mathbf{Set}
\]
is defined as follows.

\begin{enumerate}
\item[(i)] On objects, $E_R(x_1,\dots,x_n) = R^n$, with the understanding that $E_R(\varnothing)$ is a singleton set.

\item[(ii)] A morphism $\tau \colon (x_1,\dots,x_n) \to (y_1,\dots,y_m)$, corresponding to expressions $(t_1,\dots,t_m)$, induces
  \[
    E_R(\tau) \colon R^n \longrightarrow R^m,
    \qquad
    (a_1,\dots,a_n) \longmapsto \bigl(t_1(a),\dots,t_m(a)\bigr)
  \]
where the right-hand side is computed using the ring operations of $R$.
\end{enumerate}
\end{definition}

The symbolic calculus of $\mathcal{T}_\mathrm{arith}$ is thereby interpreted inside a fixed ring through $E_R$. The principle of ring coherence may then be expressed as the requirement that arithmetic quantities respect this interpretation.

\begin{definition}[Arithmetic quantities]\label{def:MultiAssignment}
Let $R$ be a commutative ring. An \emph{arithmetic quantity} on $R$ in $n$ variables is a map $Q \colon R^n \to V$, where $V$ is an abelian group or a real vector space of values. A \emph{multi-assignment of arithmetic quantities} on $R$ consists of a choice, for each object $X = (x_1,\dots,x_n)$ of $\mathcal{T}_\mathrm{arith}$, of a family $\mathcal{Q}_X$ of arithmetic quantities $Q \colon E_R(X) = R^n \to V_Q$ with various value groups $V_Q$.
\end{definition}

In practice $Q$ may be a height, a log-volume, or an expression that aggregates local contributions attached to a tuple of elements of $R$. The coherence condition expresses the expectation that equivalent symbolic presentations lead to equal values of such quantities.

\begin{definition}[Categorical ring coherence]\label{def:RingCoherenceCategorical}
A multi-assignment of arithmetic quantities on $R$ satisfies \emph{categorical ring coherence} when for every morphism $\tau \colon X \to Y$ of $\mathcal{T}_\mathrm{arith}$ and every quantity $Q \in \mathcal{Q}_Y$, the composite $Q \circ E_R(\tau)$ belongs to the family $\mathcal{Q}_X$. That is, the value of an arithmetic quantity attached to a tuple of terms depends only on the elements of $R$ that these terms denote, and not on the symbolic route by which they were assembled.
\end{definition}

\begin{remark}[Suspended evaluation]\label{rem:SuspendedEvaluation}
The labeled signature of Definition~\ref{def:Signatures} suspends categorical ring coherence locally. Two symbolic routes that denote the same element of $R$ are no longer obliged to yield the same value, because the route itself -- the label, the construction history -- has been admitted as part of the datum. The framework of the following sections may be read as a disciplined way of keeping such suspended evaluations, by attaching to each symbolic route a cost, so that two routes denoting the same element differ by an amount controlled by the distance between their histories.
\end{remark}


\section{Foundational motivations for quantitative identity}\label{sec:FoundationalIdentity}

The proposal that the logical distance between two isomorphic copies of an arithmetic object need not be vacuous stands at variance with the formalism that governs current practice. It is therefore fitting to recall that alternative attitudes toward identity have a long and well-established history, and that the quantitative discipline pursued here is less unfamiliar than it may at first appear.

\subsection{Hilbertian formalism and the exclusion of history}

Within the classical Hilbertian perspective~\cite{vonHeijenoort1967Hilbert}, a formal theory is constituted by a stock of symbols, a grammar of formation, and a set of derivation rules, while the semantic interpretation is a supplementary apparatus added once the formal system has been fixed. Practice that aligns with this attitude treats a structure $\mathcal{X}$ strictly as its isomorphism class $[\mathcal{X}]$. Once two models $M$ and $M'$ are proved isomorphic, no room remains for a distinction between them, and the isomorphism $\psi \colon M \xrightarrow{\sim} M'$ behaves as a transparent path.

In such an environment, any height or volume $H$ on the category of arithmetic objects must satisfy the extensionality principle
\[
  X \cong Y \implies H(X) = H(Y)
\]
A distance between $X$ and $Y$ is forbidden, regarded either as identically zero or as a residue of presentation that must be purged once the formalism has been expressed in invariant terms. This is a coherent and powerful stance. It is also, as we observe in Section~\ref{sec:Comparison}, exactly the stance under which a comparison between two presentations collapses to a tautology.

\subsection{Pre-Hilbertian concerns: sense, reference, and construction history}
\label{subsec:PreHilbert}

The distinction between extensional sameness and intensional identity has a persistent ancestry, running from the \textit{Sinn} and \textit{Bedeutung} distinction of Frege~\cite{Frege1892}, through the constructive epistemology of Brouwer~\cite{Brouwer1908}, to the intensional type theory of Martin-L\"of~\cite{MartinLof1975}. Within this lineage, the membership of a mathematical object in an extensional set does not exhaust its identity, which is constituted also by the manner in which the object was introduced and by the history of the constructions leading to it. An element $a$ of a type $A$ remains distinct from an element $b$ of $A$ unless they share the same introductory form, even where they eventually yield the same observable value.

This constructive lineage finds a contemporary realization in the theory of C-systems and univalent foundations~\cite{voevodsky2023martin}. There the interpretation of identity is governed by the J-structure, a mechanism that treats equality as a structure open to transport. The formal distance between syntactic C-systems, which carry the rich data of derivation, and semantic C-systems, which reside in categories such as sets or sheaves, measures the price one pays when attempting to collapse intensional history into extensional semantics.

\begin{remark}[Identity types as paths]\label{rem:IntensionalRealignment}
Within the homotopy-theoretic perspective, identity types behave as paths rather than as predicates. The effort required to construct non-degenerate J-structures on semantic categories, as documented by Voevodsky~\cite{voevodsky2023martin} for the passage from syntactic to semantic C-systems, shows that intensional identity must be installed in a model through explicit architectural choices. It does not follow inferentially from the mere fixation of an underlying universe.

The work of Riehl~\cite{Riehl2024} on the $\infty$-topos semantics of homotopy type theory affords a further confirmation. There the univalence axiom identifies the identity type of a universe with the type of equivalences, so that distinct paths between isomorphic objects keep their individuality up to higher coherence. A space in which all such paths have been contracted to a single point corresponds, in the vocabulary of this work, to Universe~H -- the extensional environment where isomorphic number fields are identified without qualification. Standard arithmetic practice thus appears as the adoption of a particular extensional discipline, a choice to work in a truncated space where the paths between isomorphic points have been rendered contractible. The refusal of this contraction produces a richer ambient in which the logical distance between two isomorphic copies may be witnessed by explicit transport data.
\end{remark}

From this vantage, a label attached to a copy of a number field, or the construction history of a configuration, serves as a concise trace of sense or of derivation rather than as a typographical decoration. Two copies $K_1$ and $K_2$ of a field $K$ that are isomorphic as rings need not be interchangeable in a context where their labels encode distinct paths of construction. The statement $K_1 \cong K_2$ then carries equality of algebraic form, while leaving open a logical space in which the senses of $K_1$ and $K_2$ differ.

\begin{remark}[An intensional cover of an extensional category]\label{rem:IntensionalCover}
The role of labels may be expressed by viewing them as part of an intensional cover of an extensional arithmetic category. Let $\mathcal{C}_{\mathrm{ext}}$ be a category such as that of number fields and ring homomorphisms, read in Universe~H. Consider a second category $\mathcal{C}_{\mathrm{int}}$ whose objects are pairs $(X,\lambda)$, where $X$ is an object of $\mathcal{C}_{\mathrm{ext}}$ and $\lambda$ is a label or construction history, and whose morphisms respect these labels. There is a forgetful functor
\[
  U \colon \mathcal{C}_{\mathrm{int}} \longrightarrow \mathcal{C}_{\mathrm{ext}}
\]
which erases labels and remembers only the underlying extensional object. Two labeled copies $(X,\lambda_1)$ and $(X,\lambda_2)$ share the same extensional image, yet they need not be identified in $\mathcal{C}_{\mathrm{int}}$ in the absence of a label-preserving isomorphism between them. The passage from a sentence about labeled objects to its image under $U$ may collapse distinctions that were significant above. The whole apparatus of the following sections refines this forgetful functor into a $2$-categorical and quantitative form.
\end{remark}

\subsection{Quantitative identity as controlled intensionality}

The quantitative categories introduced in Section~\ref{sec:QuantitativeCategories} formalize this nuanced stance. Within such a category, a morphism $f \colon X \to Y$ represents an isomorphism able to carry a non-zero load. The cost $c(f)$ registers a semantic distortion or informational effort that is invisible at the structural level of static identification. The height $H$ then exchanges strict invariance for subadditivity, satisfying
\[
  H(Y) \leq H(X) + c(f)
\]
for every morphism $f \colon X \to Y$. An isomorphism with $c(f) > 0$ thus articulates an equivalence that preserves algebraic structure while registering the dissipation of information about the construction history of the object. The cost $c(f)$ measures the decoding effort associated with the derivation, so that the history of formation and the position in a network of links acquire a numerical weight, comparable to the description length in the theory of Kolmogorov complexity~\cite{Kolmogorov1965}.

\subsection{Exigencies for a distance-based semantics}

The passage from a qualitative category to a quantitative one imposes new architectural requirements, more exacting than those it replaces. It is requisite, first, that the costs $c(f)$ be stable under composition and factorization, respecting the triangle inequality of a Lawvere metric space in the sense of Section~\ref{sec:QuantitativeCategories}, for otherwise the distances would become arbitrary decorations without structural import.

Simultaneously, this quantitative regime must form a conservative extension of the classical view. When restricted to a subcategory of protocol-determined morphisms, such as identity maps or standard changes of coordinates, the costs must vanish, reproducing the invariance of heights and volumes. The theory must contain the standard signature as its zero-cost limit, a requirement made precise as Proposition~\ref{prop:ConservativeExtension}.

For these distances to serve as substantive bounds in Diophantine geometry, the cost must finally be effective. The value $c(f)$ cannot remain abstract. It must be computable from the arithmetic data and must relate transparently to known invariants such as the discriminant and the conductor. The elliptic packages of Section~\ref{sec:EllipticPackages} are introduced precisely to meet this last exigency.


\section{Quantitative categories and Lawvere enrichment}\label{sec:QuantitativeCategories}

We now describe the structure that carries the weight of a transport. The reading through enriched category theory is placed first, since it furnishes the language in which every later cost will be measured.

\subsection{The Lawvere quantale}

\begin{definition}[The cost quantale]\label{def:CostQuantale}
Let $\mathcal{V}$ denote the structure
\[
  \mathcal{V} = \bigl([0,\infty],\ \ge,\ +,\ 0\bigr)
\]
where $[0,\infty]$ is the extended non-negative reals, the order is the \emph{reverse} of the natural order (a morphism $a \to b$ exists exactly when $a \ge b$, so that $0$, the costless identification, is the terminal object), the tensor is addition with the convention $a + \infty = \infty$, and the unit is $0$. As a poset $[0,\infty]$ is complete, addition preserves the order in each variable and distributes over arbitrary infima, and $0$ is the greatest element for $\ge$. Thus $\mathcal{V}$ is a small complete monoidal poset -- a commutative quantale in the sense suited to enrichment.
\end{definition}

The reversal of the order is the point on which the whole reading turns. A smaller cost is a stronger relation, the unit $0$ is the cost of doing nothing, and the tensor adds costs along a composite. The triangle inequality and the vanishing of self-distance, which one would otherwise impose by hand, become the two axioms of an enriched category.

\begin{definition}[$\mathcal{V}$-category]\label{def:VCategory}
A \emph{$\mathcal{V}$-category} $\mathsf{X}$ consists of a set $\mathrm{Ob}(\mathsf{X})$ together with, for every ordered pair of objects $X,Y$, an element
\[
  \mathsf{X}(X,Y) \in [0,\infty]
\]
subject to the two requirements
\[
  \mathsf{X}(X,X) = 0
  \qquad\text{and}\qquad
  \mathsf{X}(X,Z) \le \mathsf{X}(X,Y) + \mathsf{X}(Y,Z)
\]
for all objects $X,Y,Z$. A \emph{$\mathcal{V}$-functor} $F \colon \mathsf{X} \to \mathsf{Y}$ is a map of objects with $\mathsf{Y}(FX,FY) \le \mathsf{X}(X,Y)$ for all $X,Y$.
\end{definition}

Following Lawvere~\cite{Lawvere1973}, a $\mathcal{V}$-category is exactly a generalized metric space, with the asymmetry of $\mathsf{X}(X,Y)$ and $\mathsf{X}(Y,X)$ permitted, with infinite distances permitted, and with a $\mathcal{V}$-functor playing the role of a short, that is, distance-decreasing, map. The framework that follows produces, for each fixed extensional object, a $\mathcal{V}$-category whose distances are the least costs of transports, and the height is a $\mathcal{V}$-functor into the real line read as a generalized metric space.

\subsection{Quantitative categories of isomorphisms}

\begin{definition}[Quantitative category of isomorphisms]\label{def:QuantitativeCategory}
A \emph{quantitative category of isomorphisms} is a category $\mathcal{Q}$ endowed with the following data. Invertible and directed morphisms may reside side by side, and the cost applies to all arrows.

\begin{itemize}
\item[(i)] Each morphism $f \colon X \to Y$ carries a non-negative real number
  \[
    c(f) \in \R_{\ge 0}
  \]
termed its \emph{distortion cost}. The cost attaches to the morphism rather than to the endpoints, so that parallel morphisms $f,g \colon X \to Y$ may carry distinct costs. One requires
  \[
    c(\mathrm{id}_X) = 0
    \qquad\text{and}\qquad
    c(g \circ f) \le c(f) + c(g)
  \]
for every object $X$ and every composable pair.

\item[(ii)] A morphism $f \colon X \to Y$ is an \emph{isomorphism} of $\mathcal{Q}$ when it admits a two-sided inverse $g \colon Y \to X$ in the categorical sense. There is no requirement that $c(f)$ or $c(g)$ vanish. One may have $c(f) > 0$ and $c(g) > 0$ even for isomorphisms.

\item[(iii)] A functor $H \colon \mathcal{Q} \to \mathbf{R}$ to the category of real numbers and order-preserving maps, playing the role of a height or volume, satisfying the transport inequality
  \[
    H(Y) \le H(X) + c(f)
  \]
for every morphism $f \colon X \to Y$.
\end{itemize}
\end{definition}

\begin{remark}[Lawvere metric spaces and the enriched reading]\label{rem:LawvereMetric}
The axioms imposed on the cost $c$ -- namely $c(\mathrm{id}_X) = 0$ and the subadditivity $c(g \circ f) \le c(f) + c(g)$ -- reproduce the structure of a category enriched over the quantale $\mathcal{V}$ of Definition~\ref{def:CostQuantale}. The distance from $X$ to $Y$ is the infimum of costs over all morphisms between them, and the triangle inequality follows from the subadditivity of $c$ under composition. This generalized distance need not be symmetric, since the cost of transporting data from $X$ to $Y$ may differ from that of the reverse passage. The asymmetry is exactly what the framework requires, for the cost of a transport and that of its inverse are not constrained to coincide.

Universe~A thus views the category of arithmetic objects as a generalized metric space in which the distance between two configurations measures the effort, computational or geometric, of decoding one construction history into another. The height functor of part~(iii) appears in this light as a map controlled by the enrichment, satisfying a short-map condition with respect to the Lawvere distance. This anchors the quantitative categories within an established categorical lineage~\cite{Lawvere1973, Kelly1982}, while the specifically arithmetic substance of the cost -- its dependence on discriminants, conductors, and Galois-theoretic data -- remains proper to the present inquiry.
\end{remark}

\begin{definition}[Transport distance]\label{def:TransportDistance}
Let $\mathcal{Q}$ be a quantitative category of isomorphisms. For objects $X,Y$ of $\mathcal{Q}$ set
\[
  d_{\mathcal{Q}}(X,Y) \coloneqq \inf\bigl\{ c(f) : f \in \mathcal{Q}(X,Y) \bigr\}
\]
with the convention $\inf\varnothing = +\infty$. The assignment $d_{\mathcal{Q}}$ endows $\mathrm{Ob}(\mathcal{Q})$ with the structure of a $\mathcal{V}$-category in the sense of Definition~\ref{def:VCategory}, termed the \emph{transport distance} of $\mathcal{Q}$.
\end{definition}

\begin{proposition}[The height is short for the transport distance]
\label{prop:HeightShort}
Let $\mathcal{Q}$ be a quantitative category of isomorphisms with height functor $H \colon \mathcal{Q} \to \mathbf{R}$, and equip $\R$ with the $\mathcal{V}$-category structure $\R(a,b) = \max\{b - a, 0\}$. Then $H$ is a $\mathcal{V}$-functor
\[
  \bigl(\mathrm{Ob}(\mathcal{Q}), d_{\mathcal{Q}}\bigr)
  \longrightarrow
  (\R, \R(\,\cdot\,,\,\cdot\,))
\]
that is to say, $\max\{H(Y) - H(X),\, 0\} \le d_{\mathcal{Q}}(X,Y)$ for all objects $X,Y$.
\end{proposition}

\begin{proof}
Fix $X,Y$. For every morphism $f \colon X \to Y$ the transport inequality of Definition~\ref{def:QuantitativeCategory}(iii) gives $H(Y) - H(X) \le c(f)$, hence $\max\{H(Y) - H(X),0\} \le c(f)$ since the left side is non-negative. Taking the infimum over all such $f$ yields $\max\{H(Y) - H(X),0\} \le d_{\mathcal{Q}}(X,Y)$. If there is no morphism $X \to Y$, then $d_{\mathcal{Q}}(X,Y) = +\infty$ and the inequality holds trivially.
\end{proof}

Proposition~\ref{prop:HeightShort} is the smallest statement that expresses the intended reading. A height is a short map for a generalized distance rather than an invariant of isomorphism classes, and the distance is the cost of transport. In the extensional limit, where every transport between isomorphic objects is free, the distance between isomorphic objects is zero, and the short-map condition recovers the extensionality principle $X \cong Y \implies H(X) = H(Y)$.

\begin{remark}[Where the cost is arithmetic]\label{rem:CostArithmetic}
The quantale $\mathcal{V}$ and the enriched reading carry no arithmetic of their own. They organize the account. The arithmetic enters when one populates $\mathcal{Q}$ with copies of number fields and elliptic curves carrying Galois and model data, at which point the identity of an object is determined chiefly by its position within a web of correlations rather than by the extensional membership of its underlying set. Two objects may be isomorphic as bare rings while occupying distinct positions in this web, and a morphism then realigns the positions at a cost that the later sections compute explicitly.
\end{remark}


\section{The labeled quantitative \texorpdfstring{$(2,1)$}{(2,1)}-category}\label{sec:UniverseAFormal}

A quantitative category of isomorphisms assigns a cost to transports. It does not yet express the fact that two transports may agree only up to a coherent change of presentation. To keep this last distinction without collapsing it into equality, one passes to a $(2,1)$-category, where invertible $2$-morphisms witness such agreements while the scalar cost remains attached to the $1$-morphisms.

\begin{definition}[$(2,1)$-category]\label{def:TwoOneCategory}
A \emph{$(2,1)$-category} is a $2$-category in which every $2$-morphism is invertible. Equivalently, for any two objects $A$ and $B$, the hom-category $\mathfrak{A}(A,B)$ is a groupoid.
\end{definition}

\begin{remark}[A point on strictness]\label{rem:strictness}
In what follows the word ``$2$-category'' is used in a pragmatic sense. Should one prefer a weak $2$-category, that is a bicategory, then the same discussion may be reproduced by replacing $2$-functors with pseudofunctors and inserting the requisite coherence isomorphisms. Since the present inquiry treats the arithmetic role of transport rather than a full higher-categorical development, the strict language is retained.

Strictness in this sense concerns the associativity and unit laws of composition, together with the on-the-nose preservation of composition by $2$-functors, rather than the presence of $2$-morphisms. The invertible $2$-cells between parallel $1$-morphisms, upon which the invariance of the cost and the pseudonaturality of Definition~\ref{def:PseudoNatural} rest, are retained in full. The point at which the strict language reaches its limit, namely the transport of geometric data in place of scalars, is taken up in Remark~\ref{rem:StrictPrototypeLimit}.

This choice has no bearing on the arithmetic descent of Section~\ref{sec:UniverseP}. A unitor or associator cell, strict or merely coherent, acts as the identity upon the scalar data -- the costs and heights -- that Universe~P ever reads, since a transport is registered there by its cost alone and not by the $1$-morphism or the coherence cell that produced it. The finite integer tuples of a code are built from these scalars, so no coherence cell of either the strict or the weak reading can alter an entry of such a code. Where the descent does pass through a coherence cell, as it does for the intensional Szpiro inequality of Section~\ref{sec:IntensionalSzpiro}, the point is verified in detail for the toy Coq interface in Remark~\ref{rem:StrictPrototypeLimit}, and the same argument governs every later descent to a finite integer code.
\end{remark}

\begin{definition}[Homotopy $1$-category]\label{def:HoOneCategory}
Let $\mathfrak{A}$ be a $(2,1)$-category. The \emph{homotopy $1$-category} $\Ho(\mathfrak{A})$ is defined as follows. Its objects coincide with the objects of $\mathfrak{A}$. For objects $A$ and $B$, the morphisms in $\Ho(\mathfrak{A})(A,B)$ are the $2$-isomorphism classes of $1$-morphisms $A \to B$. Composition is induced by the composition of $1$-morphisms, which remains well-defined by virtue of the invertibility of all $2$-morphisms. There is a canonical functor
\[
  \pi \colon \mathfrak{A} \longrightarrow \Ho(\mathfrak{A})
\]
which acts as the identity on objects and sends a $1$-morphism to its $2$-isomorphism class.
\end{definition}

The passage $\mathfrak{A} \mapsto \Ho(\mathfrak{A})$ renders visible a distinction often lived informally. One might present a transport alongside its auxiliary choices, and two such presentations may prove equivalent without being equal. Working within $\Ho(\mathfrak{A})$ allows one to speak of transports only up to such equivalence.

\begin{definition}[Labeled quantitative $(2,1)$-category]\label{def:UniverseAFormal}
A \emph{Universe~A structure} consists of the following data:

\begin{enumerate}
\item[(i)] A $(2,1)$-category $\mathfrak{A}$ in the sense of Definition~\ref{def:TwoOneCategory}, whose
  \begin{itemize}
\item objects are triples $(X,\lambda,\mathcal{H})$, where $X$ is an underlying arithmetic object in the sense of Universe~H, such as a number field, a scheme, or an elliptic curve, where $\lambda$ is a label or construction index, and where $\mathcal{H}$ is a package of auxiliary structures such as Galois actions, model data, or other decorations
\item $1$-morphisms $f \colon (X,\lambda,\mathcal{H}) \to (Y,\lambda',\mathcal{H}')$ are structure-preserving transports, in the manner of the examples of Sections~\ref{sec:Anabelian} and~\ref{sec:EllipticPackages}
\item $2$-morphisms are invertible identifications between $1$-morphisms, encoding that two transports may differ by auxiliary choices while keeping a formal distinctness.
  \end{itemize}

\item[(ii)] A forgetful $2$-functor
  \[
    U \colon \mathfrak{A} \longrightarrow \mathcal{C}_{\mathrm{ext}}
  \]
where $\mathcal{C}_{\mathrm{ext}}$ is the extensional category of arithmetic objects of Universe~H, such as number fields with ring homomorphisms or elliptic curves with isogenies, viewed as a $2$-category with only identity $2$-morphisms. On objects, $U(X,\lambda,\mathcal{H}) = X$. On $1$-morphisms, $U$ forgets labels and auxiliary structures and retains the underlying extensional morphism. The fibers of $U$ over a fixed $X$ are the family of intensional realizations of $X$ with their labels and histories.

\item[(iii)] A quantitative structure on $1$-morphisms, given by a map
  \[
    c \colon \operatorname{Mor}_1(\mathfrak{A}) \longrightarrow \R_{\ge 0}
  \]
termed the \emph{distortion cost}, satisfying $c(\mathrm{id}_A) = 0$ and $c(g \circ f) \le c(f) + c(g)$ for every object $A$ and every composable pair. The cost is an invariant of $2$-isomorphism. If there exists an invertible $2$-morphism $\alpha \colon f \Rightarrow g$ between parallel $1$-morphisms, then $c(f) = c(g)$. Equivalently, $c$ factors through the homotopy $1$-category $\Ho(\mathfrak{A})$ of Definition~\ref{def:HoOneCategory}.

\item[(iv)] A height assignment $H \colon \operatorname{Ob}(\mathfrak{A}) \to \R$ such that for every $1$-morphism $f \colon A \to B$ the transport inequality
  \[
    H(B) \le H(A) + c(f)
  \]
holds. Since $H$ is defined on objects and $c$ respects $2$-isomorphisms, the inequality depends only on the $2$-isomorphism class of the chosen transport.

\item[(v)] A contextual structure in the sense of C-systems~\cite{voevodsky2023martin}, realized as a fibration over a base of types. For each object $A = (X,\lambda,\mathcal{H})$ there is a small category $\mathrm{Ctx}(A)$ of \emph{contexts above $A$}, whose objects are refinements or extensions of the auxiliary data $\mathcal{H}$. A $1$-morphism $f \colon A \to B$ induces a pullback functor
  \[
    f^\ast \colon \mathrm{Ctx}(B) \longrightarrow \mathrm{Ctx}(A)
  \]
compatible with composition up to coherent invertible $2$-morphisms. This structure describes how the identity rules, and in particular J-elimination, behave when one transports auxiliary data along a morphism of $\mathfrak{A}$. The pullback functors carry no cost of their own. The whole weight of a transport is borne once, by the scalar $c(f)$ of item (iii), and the substitution $f^\ast$ is free, in the sense that no quantity $c(f^\ast)$ is levied over and above $c(f)$.
\end{enumerate}
A Universe~A structure in this sense is termed a \emph{labeled quantitative $(2,1)$-category with contextual flavor}.
\end{definition}

\begin{remark}[Forgetting labels through $\Ho(\mathfrak{A})$]\label{rem:ForgetThroughHo}
Since $\mathcal{C}_{\mathrm{ext}}$ carries only identity $2$-morphisms, the forgetful $2$-functor $U$ factors canonically through $\Ho(\mathfrak{A})$, yielding a $1$-functor
\[
  \Ho(U) \colon \Ho(\mathfrak{A}) \longrightarrow \mathcal{C}_{\mathrm{ext}}
\]
which discards labels and auxiliary structures once transports have been reduced to their $2$-isomorphism classes. Within this language, the inequalities governed by $c$ and $H$ read as statements attached to morphisms in $\Ho(\mathfrak{A})$.
\end{remark}

\begin{definition}[Vertical groupoid over an extensional object]\label{def:VerticalGroupoid}
Let $(\mathfrak{A},U)$ be a Universe~A structure and let $X$ be an object of $\mathcal{C}_{\mathrm{ext}}$. The \emph{vertical groupoid over $X$}, denoted $\mathfrak{A}_X^{\mathrm{vert}}$, has as objects those $A \in \operatorname{Ob}(\mathfrak{A})$ with $U(A) = X$, and as morphisms from $A$ to $B$ the $2$-isomorphism classes of $1$-morphisms $f \colon A \to B$ such that $U(f) = \mathrm{id}_X$ and $f$ is an equivalence in $\mathfrak{A}$. Composition is induced from $\mathfrak{A}$. Because equivalences are stable under composition and admit quasi-inverses, $\mathfrak{A}_X^{\mathrm{vert}}$ is a groupoid.
\end{definition}

\begin{remark}[The locus of pure intensionality]\label{rem:locus_pure_intensionality}
In Universe~H the vertical groupoid $\mathfrak{A}_X^{\mathrm{vert}}$ is contractible, since all vertical morphisms carry zero cost. In Universe~A, morphisms in $\mathfrak{A}_X^{\mathrm{vert}}$ may carry a strictly positive distortion cost, reflecting the non-trivial effort of realigning labels or auxiliary data above a fixed extensional object. The whole substance of the framework lives in this fiber, and a comparison statement that retains content is one whose two sides sit at positive vertical distance.
\end{remark}

\begin{definition}[Normalization datum]\label{def:UniverseANormalization}
Let $(\mathfrak{A},U,c,H)$ be a Universe~A structure. A \emph{normalization datum} on $\mathfrak{A}$ comprises, for every object $X$ of $\mathcal{C}_{\mathrm{ext}}$, a representative $\mathrm{NF}(X) \in \operatorname{Ob}(\mathfrak{A})$ with $U(\mathrm{NF}(X)) = X$, and for every object $A$, a distinguished $1$-morphism
\[
  \nu_A \colon \mathrm{NF}(U(A)) \longrightarrow A
\]
whose image under $U$ is $\mathrm{id}_{U(A)}$. One defines the \emph{normalization defect} of $A$ as $\delta_{\mathfrak{A}}(A) \coloneqq c(\nu_A)$.
\end{definition}

\begin{remark}[On reference points]\label{rem:reference_points}
Such a datum demands neither the uniqueness of $\mathrm{NF}(X)$ nor that $\nu_A$ be of minimal cost. It serves to fix a disciplined selection of a reference point within each extensional fiber, together with a distinguished path required to reach a given object. This mirrors the normalization procedures invoked for elliptic packages in Definition~\ref{def:AccumulatedDefect}, extending that logic to the ambient categorical environment.
\end{remark}


\section{The J-structure and controlled elimination}\label{sec:Jstructure}

The vertical groupoid of Definition~\ref{def:VerticalGroupoid} admits a reading through the optics of C-systems, and it is this reading that ties the arithmetic cost to the identity type of a dependent theory.

\begin{remark}[Vertical groupoids and identity types]\label{rem:VerticalGroupoidCsystems}
In a C-system~\cite{voevodsky2023martin}, the identity type $\mathrm{Id}_A(a,a)$ over a fixed term $a$ of type $A$ is itself a type whose inhabitants witness paths from $a$ to $a$. When the C-system is non-trivial in the homotopical sense, this type may contain inhabitants other than reflexivity, each carrying its own computational or geometric content. In the present framework, the role of $\mathrm{Id}_A(a,a)$ is played by the vertical groupoid $\mathfrak{A}_X^{\mathrm{vert}}$. A morphism $f \colon A \to B$ in this groupoid, with $U(f) = \mathrm{id}_X$, represents a non-trivial self-identification of the extensional object $X$ -- a path that permutes labels or auxiliary data while leaving the underlying arithmetic structure fixed. In Universe~H, where the J-structure has been collapsed and all self-identifications are forced to be reflexivity, the vertical groupoid is contractible. In Universe~A it retains the richer structure of a groupoid whose morphisms may carry positive distortion cost.

This parallel reinforces the thesis that the intensional environment of Universe~A arises of its own accord, wherever one declines to contract the identity type to a mere proposition. Such a refusal has received independent mathematical life in homotopy type theory and its $\infty$-topos semantics~\cite{Riehl2024}.
\end{remark}

\begin{definition}[Lawvere--Voevodsky interpretation of the vertical cost]
\label{def:LawvereVoevodskyInterpretation}
Let $(\mathfrak{A},U,c,H)$ be a Universe~A structure and let $X$ be an object of $\mathcal{C}_{\mathrm{ext}}$. The \emph{Lawvere--Voevodsky interpretation} of the vertical groupoid $\mathfrak{A}_X^{\mathrm{vert}}$ comprises the following data:

\begin{enumerate}
\item[(i)] \textit{Vertical distance.} For objects $A,B \in \operatorname{Ob}(\mathfrak{A}_X^{\mathrm{vert}})$ set
\[
  d_X^{\mathrm{vert}}(A,B)
  \coloneqq
  \inf\bigl\{ c(f) : f \in \mathfrak{A}_X^{\mathrm{vert}}(A,B) \bigr\}
\]
with $\inf\varnothing = +\infty$. By the cost axioms of Definition~\ref{def:UniverseAFormal}(iii), the assignment $d_X^{\mathrm{vert}}$ satisfies $d_X^{\mathrm{vert}}(A,A) = 0$ and the triangle inequality, and so endows $\operatorname{Ob}(\mathfrak{A}_X^{\mathrm{vert}})$ with the structure of a Lawvere generalized metric space, a $\mathcal{V}$-category in the sense of Definition~\ref{def:VCategory}.

\item[(ii)] \textit{Path-space reading.} Under the correspondence
\[
  \mathrm{Id}_{\mathcal{C}_{\mathrm{ext}}}(X,X)
  \;\longleftrightarrow\;
  \mathfrak{A}_X^{\mathrm{vert}}
\]
an inhabitant of the identity type corresponds to a vertical morphism $f \colon A \to B$ with $U(f) = \mathrm{id}_X$, and the cost $c(f)$ measures the length of that path in the Lawvere metric.

\item[(iii)] \textit{Reflexivity and zero cost.} The reflexivity term $\mathrm{refl}_A$ corresponds to $\mathrm{id}_A$ in $\mathfrak{A}_X^{\mathrm{vert}}$, and the cost axiom $c(\mathrm{id}_A) = 0$ expresses that reflexivity carries zero length.
\end{enumerate}
\end{definition}

\begin{proposition}[Controlled J-elimination from the transport inequality]
\label{prop:ControlledJElimination}
Let $(\mathfrak{A},U,c,H)$ be a Universe~A structure and let $d_X^{\mathrm{vert}}$ be the vertical distance of Definition~\ref{def:LawvereVoevodskyInterpretation}.

\begin{enumerate}
\item[\textup{(i)}] \textit{Quantitative J-rule.} For every vertical morphism $f \colon A \to B$ in $\mathfrak{A}_X^{\mathrm{vert}}$ and every height $H$ satisfying the transport inequality of Definition~\ref{def:UniverseAFormal}(iv), one has
\[
  |H(B) - H(A)| \le c(f) + c(f^{-1})
\]
where $f^{-1}$ is the inverse in the groupoid. If $c(f) = c(f^{-1}) = 0$, then $H(B) = H(A)$, recovering the classical J-elimination principle that transport along reflexivity preserves all structure.

\item[\textup{(ii)}] \textit{Metric control on the height.} The height $H$, restricted to $\operatorname{Ob}(\mathfrak{A}_X^{\mathrm{vert}})$, is Lipschitz with respect to the symmetrized vertical distance
\[
  \bar{d}_X^{\mathrm{vert}}(A,B)
  \coloneqq
  \inf_{f \in \mathfrak{A}_X^{\mathrm{vert}}(A,B)}
  \bigl( c(f) + c(f^{-1}) \bigr)
\]
in the precise sense that $|H(B) - H(A)| \le \bar{d}_X^{\mathrm{vert}}(A,B)$ for all objects $A,B$.

\item[\textup{(iii)}] \textit{Truncation.} In Universe~H the vertical groupoid is essentially trivial, every vertical morphism carrying zero cost. The metric $d_X^{\mathrm{vert}}$ vanishes identically, and the J-elimination rule degenerates to the extensional principle that isomorphic objects share the same height.
\end{enumerate}
\end{proposition}

\begin{proof}
For (i), apply the transport inequality to $f$ and to $f^{-1}$ separately, then combine. For (ii), the infimum over all $f \in \mathfrak{A}_X^{\mathrm{vert}}(A,B)$ of $c(f) + c(f^{-1})$ bounds $|H(B) - H(A)|$ by (i). For (iii), in Universe~H the vertical groupoid is contractible by Remark~\ref{rem:locus_pure_intensionality}, so every vertical morphism $f$ satisfies $c(f) = 0$.
\end{proof}

\begin{remark}[Quantitative categories with families]\label{rem:QuantitativeCwF}
The structure of Definition~\ref{def:LawvereVoevodskyInterpretation} parallels the quantitative categories with families introduced by Atkey~\cite{Atkey2018} as the categorical semantics of quantitative type theory. There the standard categories-with-families semantics for Martin-L\"of type theory is enriched by a semiring of usage annotations, and substitution is required to respect these annotations. In the present context the role of the semiring is played by the Lawvere quantale $([0,\infty],\ge,+)$, and the usage of a morphism is its distortion cost $c(f)$. The J-elimination rule of Proposition~\ref{prop:ControlledJElimination}(i) reads as the statement that transport along a zero-usage path preserves height exactly, while transport along a path of positive usage incurs a bounded penalty. This is the sense in which Universe~A implements a resource-sensitive discipline of identity, where the effort of decoding one intensional presentation into another is held as a quantitative datum rather than erased by a truncation.

A word is owed on substitution, since Definition~\ref{def:UniverseAFormal}(v) introduces the pullback functors $f^\ast \colon \mathrm{Ctx}(B) \to \mathrm{Ctx}(A)$ without a cost of their own. The choice made there, that context substitution is free, is the choice of Atkey's semantics as well. In a quantitative category with families, substitution respects the usage annotations and re-indexes them along the morphism, and it consumes no resource in doing so, the whole expenditure being carried by the terms and, in the present reading, by the $1$-morphism $f$. The cost $c$ is accordingly defined upon $\operatorname{Mor}_1(\mathfrak{A})$ alone, and a context of the form $f^\ast\Gamma$ is charged only through $c(f)$ wherever it enters a transport inequality. One could entertain a finer discipline in which a quantity $c(f^\ast)$ were levied upon the substitution itself, subadditive along the coherence isomorphisms $(g \circ f)^\ast \cong f^\ast \circ g^\ast$ and bounded above by $c(f)$, and such a discipline would sit closer still to the annotated substitutions of~\cite{Atkey2018}. Nothing below requires that refinement, and every construction of this paper reads the weight of a transport once, at the level of the $1$-morphism.

The $\infty$-topos semantics of Riehl~\cite{Riehl2024} affords the geometric realization of this picture. In an $\infty$-topos the identity type of a universe is the type of equivalences, and distinct paths between equivalent objects keep their individuality up to higher coherence. Universe~H corresponds to the $0$-truncation of this space, where all paths have been contracted. Universe~A corresponds to an intermediate regime in which paths are retained and measured, yet the measurement obeys the compositional discipline of a Lawvere metric.
\end{remark}

\subsection{Universe H as the zero-cost limit}

The conservativity demanded in Section~\ref{sec:FoundationalIdentity} can now be stated as a proposition. Universe~H is recovered from any Universe~A structure by contracting every vertical fiber, and the recovery is compatible with the height.

\begin{definition}[The extensional collapse]\label{def:ExtensionalCollapse}
Let $(\mathfrak{A},U,c,H)$ be a Universe~A structure. The \emph{extensional collapse} is the functor
\[
  \kappa \colon \Ho(\mathfrak{A}) \longrightarrow \mathcal{C}_{\mathrm{ext}}
\]
of Remark~\ref{rem:ForgetThroughHo}, together with the equivalence relation on objects that identifies $A$ and $B$ whenever $U(A) = U(B)$ and there exists a vertical equivalence $A \to B$. We write $\mathcal{C}_{\mathrm{ext}}^{\,\kappa}$ for the quotient.
\end{definition}

\begin{proposition}[Conservative extension]\label{prop:ConservativeExtension}
Let $(\mathfrak{A},U,c,H)$ be a Universe~A structure, and suppose that the height $H$ is \emph{extensionally coherent}, meaning that $H(A) = H(B)$ whenever $U(A) = U(B)$ and the vertical distance $\bar{d}_{U(A)}^{\mathrm{vert}}(A,B)$ vanishes. Then there is a well-defined extensional height
\[
  H^{\mathrm{ext}} \colon \operatorname{Ob}(\mathcal{C}_{\mathrm{ext}}^{\,\kappa})
  \longrightarrow \R,
  \qquad
  H^{\mathrm{ext}}\bigl([X]\bigr) \coloneqq H(A) \text{ for any } A \text{ with } U(A) = X
\]
and $H^{\mathrm{ext}}$ is an invariant of isomorphism classes in $\mathcal{C}_{\mathrm{ext}}$. Conversely, every extensional height on $\mathcal{C}_{\mathrm{ext}}$ lifts to an extensionally coherent height on $\mathfrak{A}$ by composition with $U$, and the cost of every vertical morphism under this lift is zero.
\end{proposition}

\begin{proof}
If $U(A) = U(B) = X$ and $A,B$ lie at zero symmetrized vertical distance, then $|H(A) - H(B)| \le \bar{d}_X^{\mathrm{vert}}(A,B) = 0$ by Proposition~\ref{prop:ControlledJElimination}(ii), so extensional coherence makes $H^{\mathrm{ext}}([X])$ independent of the chosen representative. Isomorphic objects of $\mathcal{C}_{\mathrm{ext}}$ have equal extensional images under any representative, so $H^{\mathrm{ext}}$ is an invariant of isomorphism classes. For the converse, given an extensional height $h$ on $\mathcal{C}_{\mathrm{ext}}$, set $H \coloneqq h \circ U$. Then for any $1$-morphism $f \colon A \to B$ one has $H(B) - H(A) = h(U(B)) - h(U(A))$, which equals $0$ when $U(f) = \mathrm{id}$ since $h$ is invariant. Thus the transport inequality holds with cost zero on every vertical morphism.
\end{proof}

\begin{remark}[The reading of the proposition]\label{rem:readingConservative}
Proposition~\ref{prop:ConservativeExtension} states the precise sense in which the framework loses nothing. Every theorem of the extensional discipline is a theorem of the intensional one, read at zero cost, while the intensional discipline carries in addition the positive distances of the vertical fibers. A statement of Universe~H is the shadow under $\kappa$ of a statement of Universe~A whose vertical content has been contracted. The next section shows that a comparison statement is exactly a statement whose vertical content refuses to be contracted, and that this refusal is the source of whatever Diophantine substance the comparison may carry.
\end{remark}


\section{Comparison problems and the descent obstruction}\label{sec:Comparison}

We come to the categorical reinterpretation named in the title. A comparison problem, in the present reading, is the data of two intensional realizations of one extensional object together with a transport between them, and the question attached to it is whether the resulting transport inequality descends to a non-trivial inequality between extensional invariants. The answer is governed entirely by the vertical groupoid, and the obstruction to descent is exact.

\subsection{Comparison problems}

\begin{definition}[Comparison problem]\label{def:ComparisonProblem}
Let $(\mathfrak{A},U,c,H)$ be a Universe~A structure and let $X$ be an object of $\mathcal{C}_{\mathrm{ext}}$. A \emph{comparison problem over $X$} is a pair of objects $A_{\mathrm{src}}, A_{\mathrm{tgt}}$ of $\mathfrak{A}$ with $U(A_{\mathrm{src}}) = U(A_{\mathrm{tgt}}) = X$, together with a vertical transport
\[
  f \colon A_{\mathrm{src}} \longrightarrow A_{\mathrm{tgt}}
\]
in $\mathfrak{A}_X^{\mathrm{vert}}$. Its \emph{comparison inequality} is the transport inequality
\[
  H(A_{\mathrm{tgt}}) \le H(A_{\mathrm{src}}) + c(f)
\]
and its \emph{content} is the quantity $c(f) \in \R_{\ge 0}$, the vertical cost of the transport.
\end{definition}

\begin{remark}[Why this captures the informal notion]\label{rem:WhyComparison}
In arithmetic practice a comparison statement places two realizations of one object side by side -- two models of an elliptic curve, two reconstructions of a field from local data, two presentations of a fundamental group -- and asserts an inequality between invariants attached to them. The two realizations share an underlying extensional object $X$, and the transport $f$ is the explicit passage from one realization to the other. The content of the comparison is the cost of that passage. When the passage is free, the comparison says nothing beyond $H(X) \le H(X)$. When the passage carries a positive cost, the comparison registers that cost. The next two results make this dichotomy precise.
\end{remark}

\subsection{A two-sided descent lemma}

We first describe, in a form requiring no special geometry, the mechanism by which a comparison collapses when one insists on reading it in the extensional signature. The lemma is internal to Universe~H.

\begin{definition}[Toy labeled and standard data]\label{def:ToyAliasedStandard}
Let $X$ be an object of $\mathcal{C}_{\mathrm{ext}}$. A \emph{toy labeled configuration over $X$} is a tuple
\[
  \mathcal{C} = \bigl(A_{\mathrm{src}}, A_{\mathrm{tgt}}, f, V^{\mathrm{src}}, V^{\mathrm{tgt}}, \delta\bigr)
\]
consisting of two objects $A_{\mathrm{src}}, A_{\mathrm{tgt}}$ above $X$, a vertical transport $f \colon A_{\mathrm{src}} \to A_{\mathrm{tgt}}$, two real numbers $V^{\mathrm{src}}, V^{\mathrm{tgt}}$ playing the role of intensional volumes attached to the two avatars, and a non-negative real number $\delta$ with
\[
  V^{\mathrm{tgt}} \le V^{\mathrm{src}} + \delta
\]
There is a natural involution $\sigma$ on such data that exchanges source and target, sending $\mathcal{C}$ to
\[
  \sigma(\mathcal{C}) = \bigl(A_{\mathrm{tgt}}, A_{\mathrm{src}}, f^{-1}, V^{\mathrm{tgt}}, V^{\mathrm{src}}, \delta'\bigr)
\]
where $\delta'$ is the non-negative real number appearing in the reverse inequality $V^{\mathrm{src}} \le V^{\mathrm{tgt}} + \delta'$.
\end{definition}

\begin{lemma}[Two-sided descent lemma]\label{lemma:ToyNoBridge}
Suppose there is a rule
\[
  \mathcal{B}^{\mathrm{toy}}
  \colon
  \{\text{toy labeled configurations over } X\}
  \longrightarrow
  \R^2 \times \R_{\ge 0},
  \qquad
  \mathcal{C} \longmapsto
  \bigl(H^{\mathrm{std}}_{\mathrm{src}}(\mathcal{C}), H^{\mathrm{std}}_{\mathrm{tgt}}(\mathcal{C}), \Delta(\mathcal{C})\bigr)
\]
to be read as a standard inequality in Universe~H,
\begin{equation}\label{eq:ToyStdIneq}
  H^{\mathrm{std}}_{\mathrm{tgt}}(\mathcal{C})
  \le
  H^{\mathrm{std}}_{\mathrm{src}}(\mathcal{C}) + \Delta(\mathcal{C}).
\end{equation}
Assume the following:
\begin{itemize}
\item[(i)] \emph{Equivariance under exchange.} If $\mathcal{C}$ and $\mathcal{C}'$ differ only by replacing the objects above $X$ by isomorphic copies and by exchanging the roles of source and target, then $\mathcal{B}^{\mathrm{toy}}(\mathcal{C})$ and $\mathcal{B}^{\mathrm{toy}}(\mathcal{C}')$ agree up to the same exchange. In particular $\Delta(\mathcal{C}) = \Delta(\sigma(\mathcal{C}))$ and $H^{\mathrm{std}}_{\mathrm{src}}(\mathcal{C}) = H^{\mathrm{std}}_{\mathrm{tgt}}(\sigma(\mathcal{C}))$.

\item[(ii)] \emph{Label transparency.} The quantities $H^{\mathrm{std}}_{\mathrm{src}}(\mathcal{C})$ and $H^{\mathrm{std}}_{\mathrm{tgt}}(\mathcal{C})$ are values of a single ring-coherent invariant in the sense of Axiom~\ref{ax:RingCoherence}, evaluated on the underlying object $X$ after static identification. Once the labels have been erased, there is no remaining standard datum that distinguishes the source evaluation from the target evaluation.

\item[(iii)] \emph{Intended survival of the discrepancy.} There is at least one configuration $\mathcal{C}$ with $V^{\mathrm{tgt}} \ne V^{\mathrm{src}}$ for which~\eqref{eq:ToyStdIneq} is intended to keep a non-trivial trace of the discrepancy, so that the rule does not force $H^{\mathrm{std}}_{\mathrm{src}}(\mathcal{C}) = H^{\mathrm{std}}_{\mathrm{tgt}}(\mathcal{C})$ by stipulation.
\end{itemize}
Then no such rule exists. Equivalently, under (i) and (ii) any rule producing a standard inequality of the form~\eqref{eq:ToyStdIneq} forces $H^{\mathrm{std}}_{\mathrm{src}}(\mathcal{C}) = H^{\mathrm{std}}_{\mathrm{tgt}}(\mathcal{C})$, and the comparison collapses in Universe~H to a harmless inequality.
\end{lemma}

\begin{proof}
Fix a configuration $\mathcal{C}$. By equivariance (i), the output triples of $\mathcal{C}$ and $\sigma(\mathcal{C})$ agree up to exchange, so that
\[
  H^{\mathrm{std}}_{\mathrm{src}}(\mathcal{C}) = H^{\mathrm{std}}_{\mathrm{tgt}}(\sigma(\mathcal{C}))
  \qquad\text{and}\qquad
  H^{\mathrm{std}}_{\mathrm{tgt}}(\mathcal{C}) = H^{\mathrm{std}}_{\mathrm{src}}(\sigma(\mathcal{C})).
\]
By label transparency (ii), once one passes to Universe~H and performs static identification, there is no standard datum distinguishing the source evaluation from the target evaluation, so the two outputs coincide,
\[
  H^{\mathrm{std}}_{\mathrm{src}}(\mathcal{C}) = H^{\mathrm{std}}_{\mathrm{tgt}}(\mathcal{C}).
\]
Then~\eqref{eq:ToyStdIneq} reduces to $H \le H + \Delta(\mathcal{C})$ with $\Delta(\mathcal{C}) \ge 0$, which holds independently of any discrepancy between $V^{\mathrm{src}}$ and $V^{\mathrm{tgt}}$. This contradicts the intended survival condition (iii).
\end{proof}

\subsection{The general descent obstruction}

The two-sided lemma is the elementary form of a single statement about the framework. A comparison inequality descends to the extensional quotient as a non-trivial inequality only when the transport is not a vertical equivalence of zero symmetrized cost. The proof is immediate from Proposition~\ref{prop:ControlledJElimination}, and the transparency of the statement is the principal benefit of the categorical reading.

\begin{theorem}[Descent obstruction]\label{thm:DescentObstruction}
Let $(\mathfrak{A},U,c,H)$ be a Universe~A structure with an extensionally coherent height in the sense of Proposition~\ref{prop:ConservativeExtension}, and let $(A_{\mathrm{src}}, A_{\mathrm{tgt}}, f)$ be a comparison problem over an object $X$. Let $H^{\mathrm{ext}}$ be the extensional height induced on $\mathcal{C}_{\mathrm{ext}}^{\,\kappa}$ by Proposition~\ref{prop:ConservativeExtension}. Then the images of $A_{\mathrm{src}}$ and $A_{\mathrm{tgt}}$ under the extensional collapse $\kappa$ satisfy
\[
  H^{\mathrm{ext}}\bigl(\kappa(A_{\mathrm{src}})\bigr)
  =
  H^{\mathrm{ext}}\bigl(\kappa(A_{\mathrm{tgt}})\bigr)
\]
whenever the symmetrized vertical cost $c(f) + c(f^{-1})$ vanishes. Consequently the comparison inequality descends to $H^{\mathrm{ext}} \le H^{\mathrm{ext}}$, a tautology, exactly when the transport is free. The comparison transmits content to the extensional level only through the residual term $c(f) + c(f^{-1})$, and that term is invisible to $H^{\mathrm{ext}}$.
\end{theorem}

\begin{proof}
Both objects lie above $X$, so $\kappa(A_{\mathrm{src}})$ and $\kappa(A_{\mathrm{tgt}})$ are the same class $[X]$ in $\mathcal{C}_{\mathrm{ext}}^{\,\kappa}$ as soon as $\bar{d}_X^{\mathrm{vert}}(A_{\mathrm{src}}, A_{\mathrm{tgt}}) = c(f) + c(f^{-1}) = 0$, by Definition~\ref{def:ExtensionalCollapse}. Extensional coherence then gives $H^{\mathrm{ext}}(\kappa(A_{\mathrm{src}})) = H^{\mathrm{ext}}(\kappa(A_{\mathrm{tgt}}))$, and the comparison inequality $H(A_{\mathrm{tgt}}) \le H(A_{\mathrm{src}}) + c(f)$ descends to the equality of the two extensional heights together with the slack $c(f)$. By Proposition~\ref{prop:ControlledJElimination}(ii) the height cannot separate two objects at zero symmetrized vertical distance, so no part of the content survives in $H^{\mathrm{ext}}$.
\end{proof}

\begin{remark}[The import of the obstruction]\label{rem:MoralObstruction}
Theorem~\ref{thm:DescentObstruction} and Lemma~\ref{lemma:ToyNoBridge} express the same fact from two directions. A comparison statement that one wishes to read, at the end, as an inequality between ordinary invariants of a single object will collapse to a tautology, because the extensional height is blind to the vertical fiber where the comparison lives. The difficulty lies in the tension between the dynamic identification of the labeled signature and the static identification of Universe~H, rather than in any deficiency of analytic structure on either side of the comparison. If one wishes a comparison to retain Diophantine content, then one must accept the positive vertical cost as part of the statement, and read the inequality in Universe~A rather than after the collapse. The Diophantine sections below adopt exactly this discipline, keeping the model defect visible as a term of the inequality.
\end{remark}


\section{Anabelian packages: a quantitative category from practice}\label{sec:Anabelian}

Before turning to the Diophantine application, it is worth exhibiting that quantitative categories of isomorphisms are not an artifice imposed from outside. They arise within established arithmetic practice, and anabelian geometry furnishes a transparent instance, where the bare profinite group of a curve carries less information than the group endowed with its Galois action and decomposition data, so that an isomorphism of bare groups incurs a cost from the start when read at the enriched level.

Let $k$ be a finitely generated field over $\Q$, and let $X$ be a smooth hyperbolic curve over $k$, with arithmetic fundamental group
\[
  \pi_1(X) \coloneqq \pi_1(X_{\overline{k}}) \rtimes G_k
\]
where $G_k$ is the absolute Galois group of $k$ and $\pi_1(X_{\overline{k}})$ is the geometric fundamental group. The action of $G_k$ on $\pi_1(X_{\overline{k}})$ is regarded up to inner automorphism, through the outer representation
\[
  \rho_X \colon G_k \longrightarrow \mathrm{Out}\bigl(\pi_1(X_{\overline{k}})\bigr).
\]
Anabelian results, originating in Grothendieck's vision and developed by Mochizuki~\cite{mochizuki1999local, MochizukiAbAnIII}, show that under suitable hypotheses one reconstructs the isomorphism class of the pair $(X,k)$ from the profinite group $\pi_1(X_{\overline{k}})$ equipped with this outer action and with boundary decorations such as inertia and decomposition subgroups.

In Universe~H one tends to read $\pi_1(X_{\overline{k}})$ as a profinite group up to isomorphism. If two curves have isomorphic geometric fundamental groups as abstract profinite groups, then under this extensional outlook they cannot be separated by any invariant depending only on the group structure, and the outer action is extraneous unless explicitly encoded. Universe~A introduces a different viewpoint.

\begin{definition}[Anabelian package]\label{def:AnabelianPackage}
Let $k$ be a finitely generated field over $\Q$. An \emph{anabelian package over $k$} is a tuple $\mathcal{G} = (\Pi, G_k, \rho, \mathcal{D})$, where $\Pi$ is a profinite group, $\rho \colon G_k \to \mathrm{Out}(\Pi)$ is a continuous homomorphism encoding the outer Galois action, and $\mathcal{D} = \{D_v\}_{v \in S}$ is a finite collection of closed subgroups $D_v \subset \Pi \rtimes G_k$, indexed by a finite set $S$ of places or marked points, playing the role of decomposition or inertia groups. A \emph{morphism} $f \colon \mathcal{G} \to \mathcal{G}'$ over the same base field $k$ is an isomorphism of profinite groups $\varphi \colon \Pi \xrightarrow{\sim} \Pi'$ such that the induced map on outer automorphism groups intertwines $\rho$ and $\rho'$, and such that for each $v \in S$ the image $\varphi(D_v)$ is conjugate to some $D'_{v'} \in \mathcal{D}'$. The category of anabelian packages over $k$ is denoted $\mathbf{AnPkg}_k$.
\end{definition}

\begin{definition}[Transport defect]\label{def:TransportDefect}
Let $f \colon \mathcal{G} \to \mathcal{G}'$ be a morphism in $\mathbf{AnPkg}_k$. Its \emph{transport defect} is
\[
  \delta(f) \coloneqq \#\bigl\{ v \in S : \varphi(D_v) \ne D'_{v'} \text{ while } \varphi(D_v) \sim D'_{v'} \bigr\}
\]
where $\sim$ denotes conjugacy. That is, $\delta(f)$ counts the decomposition groups preserved only up to conjugacy rather than exactly. The identity morphism satisfies $\delta(\mathrm{id}_{\mathcal{G}}) = 0$, and for composable morphisms $\delta(g \circ f) \le \delta(f) + \delta(g)$.
\end{definition}

With the transport defect $\delta$ in the role of the distortion cost, the category $\mathbf{AnPkg}_k$ becomes a quantitative category of isomorphisms in the sense of Definition~\ref{def:QuantitativeCategory}. An isomorphism $f \colon \mathcal{G} \xrightarrow{\sim} \mathcal{G}'$ may have $\delta(f) > 0$ when the underlying group isomorphism fails to preserve the decomposition data exactly. The defect measures the positional information -- the placement of distinguished subgroups within the ambient group -- that is shifted when passing from $\mathcal{G}$ to $\mathcal{G}'$.

\begin{proposition}[A statement with stable content in Universe~A]\label{prop:OrdinaryInA}
Let $X$ and $Y$ be smooth hyperbolic curves over $k$ with anabelian packages $\mathcal{G}_X$ and $\mathcal{G}_Y$. Suppose there is an isomorphism of abstract profinite groups
\[
  \varphi \colon \pi_1(X_{\overline{k}}) \xrightarrow{\sim} \pi_1(Y_{\overline{k}})
\]
which does not respect the outer Galois actions $\rho_X$ and $\rho_Y$. Then the bare profinite groups are identified in Universe~H, becoming isomorphic once one keeps only the abstract group structure, while the enriched objects $\mathcal{G}_X$ and $\mathcal{G}_Y$ remain non-isomorphic in $\mathbf{AnPkg}_k$. In particular the statement
\begin{equation}\label{eq:AnabelianDistinction}
  \text{``}\,\mathcal{G}_X \not\cong \mathcal{G}_Y\,\text{''}
\end{equation}
has a determinate and true value in Universe~A, whereas the corresponding statement about the bare groups evaluates to false in Universe~H.
\end{proposition}

\begin{proof}
By hypothesis $\varphi$ is an isomorphism of abstract profinite groups, so in Universe~H, where objects are identified up to isomorphism of the underlying group, $X$ and $Y$ become indistinguishable at the level of their geometric fundamental groups. In $\mathbf{AnPkg}_k$ a morphism from $\mathcal{G}_X$ to $\mathcal{G}_Y$ requires $\varphi$ to intertwine $\rho_X$ and $\rho_Y$. Since by hypothesis no such intertwining exists, there is no morphism, and the packages are non-isomorphic. Thus~\eqref{eq:AnabelianDistinction} is true in Universe~A and expresses the arithmetically meaningful fact that $X$ and $Y$, while sharing isomorphic geometric fundamental groups, differ in their Galois-theoretic structure.
\end{proof}

For Diophantine purposes one requires a real-valued invariant playing the role of a height. The following definition furnishes a schematic one.

\begin{definition}[Anabelian height]\label{def:AnabelianHeight}
Let $\mathcal{G} = (\Pi, G_k, \rho, \mathcal{D})$ arise from a smooth hyperbolic curve $X$ over $k$, with $\Pi = \pi_1(X_{\overline{k}})$ and $\mathcal{D}$ encoding decomposition groups at a finite set of rational points or points at infinity. The \emph{anabelian height} is
\[
  h_{\mathrm{ab}}(\mathcal{G}) \coloneqq h(X) + \sum_{v \in S} \log\, [G_k : \mathrm{Stab}_{G_k}(D_v)]
\]
where $h(X)$ is a classical height of the curve, for instance the Faltings height or a projective height, and the sum runs over the logarithmic indices of the stabilizers of the decomposition groups under the Galois action.
\end{definition}

\begin{lemma}[Transport inequality for anabelian height]\label{lem:TransportAnabelian}
Let $f \colon \mathcal{G} \to \mathcal{G}'$ be a morphism in $\mathbf{AnPkg}_k$ with transport defect $\delta(f)$. Then
\[
  h_{\mathrm{ab}}(\mathcal{G}') \le h_{\mathrm{ab}}(\mathcal{G}) + C_k \cdot \delta(f)
\]
for a constant $C_k$ depending only on $k$.
\end{lemma}

\begin{proof}
When $f$ preserves a decomposition group exactly, the corresponding term in the anabelian height is unchanged. When $f$ preserves a decomposition group only up to conjugacy, the stabilizer index may shift by a bounded amount depending on the structure of $G_k$. The transport defect $\delta(f)$ counts the number of such shifts, and the constant $C_k$ bounds the contribution of each.
\end{proof}

\begin{remark}[A faithful fragment of Universe~A]\label{rem:AnabelianFragment}
The category $\mathbf{AnPkg}_k$ affords a compact yet precise fragment of Universe~A. In Universe~H, where one works with the curve up to isomorphism, only $h(X)$ is visible. In Universe~A the full anabelian height becomes the natural invariant, and the transport inequality of Lemma~\ref{lem:TransportAnabelian} shows it transforms in a controlled way under morphisms, with the transport defect entering as an explicit correction. The anabelian reconstruction theorems imply that, under suitable hypotheses, the package $\mathcal{G}_X$ determines the isomorphism class of $X$ over $k$, so that the intensional data carry the full Diophantine content of the curve. Universe~A, by keeping this data as part of the object, preserves the capacity to make fine arithmetic distinctions that Universe~H effaces. It is not claimed that this fragment suffices to organize any large theory. It shows that categories with quantitative structure on morphisms occur in arithmetic geometry, and that the distinction between the two universes carries tangible mathematical content.
\end{remark}

%==================================================================
\section{Elliptic packages and the discriminant height}\label{sec:EllipticPackages}
%==================================================================

The preceding sections have stayed at the level of formalism. We now descend to arithmetic geometry and exhibit a quantitative category whose objects are presentations of elliptic curves and whose morphisms carry an explicit transport cost. The construction is elementary, and it rests upon the transformation law of the discriminant under a change of Weierstrass coordinates. What the intensional reading contributes is the decision to keep the presentation as part of the object and to measure, rather than to suppress, the price of passing from one presentation to another.

\subsection{An extensional anchor}\label{subsec:ConductorComplexity}

Before introducing the intensional data, we isolate the one invariant that belongs to the curve itself and not to any of its presentations.

\begin{definition}[Conductor complexity]\label{def:ConductorComplexity}
Let $\mathcal{E}$ be an elliptic package over a number field $K$ in the sense of Definition~\ref{def:EllipticPackage} below, with underlying curve $E$. The \emph{conductor complexity} of $\mathcal{E}$ is
\[
n(\mathcal{E}) \coloneqq \frac{1}{[K:\Q]} \log N_E
\]
where $N_E$ is the absolute norm of the conductor ideal of $E$ over $K$. This quantity depends solely upon the underlying curve $E$, so that it descends to Universe~H and serves as the extensional anchor against which the intensional heights are later compared.
\end{definition}

\subsection{Packages and their morphisms}\label{subsec:EllPkgDef}

\begin{definition}[Elliptic package]\label{def:EllipticPackage}
Let $K$ be a number field. An \emph{elliptic package over $K$} consists of a tuple
\[
\mathcal{E} = \bigl(E, \mathcal{W}, \{\sigma_v\}_{v \in S}, \lambda\bigr)
\]
where the four components are the following:
\begin{enumerate}
\item[(i)] $E$ is an elliptic curve over $K$.
\item[(ii)] $\mathcal{W}$ is a Weierstrass model for $E$ over $\mathcal{O}_K$ or over a localization thereof.
\item[(iii)] $\{\sigma_v\}_{v \in S}$ is a finite collection of local sections or level structures at the places $v$ within a finite set $S$ that contains the places of bad reduction of $E$.
\item[(iv)] $\lambda$ is a label or construction history, keeping the manner in which the data $(E,\mathcal{W},\{\sigma_v\})$ arose from the surrounding arithmetic context -- for instance a derivation from a modular or Shimura curve, a base change along a morphism of number fields, or a pointer to the ambient presentation from which the model was extracted.
\end{enumerate}

A \emph{morphism of elliptic packages} $f \colon \mathcal{E} \to \mathcal{E}'$ over the same base field $K$ is an isomorphism of elliptic curves
\[
  \varphi_f \colon E \xrightarrow{\sim} E'
\]
accompanied by the following compatible data:
\begin{enumerate}
\item[(i)] For each place $v$ of $K$, a change of Weierstrass coordinates relating the local equation $\mathcal{W}_v$ of $\mathcal{E}$ to the local equation $\mathcal{W}'_v$ of $\mathcal{E}'$. After extension of scalars to $K_v$, this change assumes the standard form
  \begin{equation}\label{eq:local-change-of-variables}
    x = u_v(f)^2\,x' + r_v(f)
    \qquad
    y = u_v(f)^3\,y' + u_v(f)^2 s_v(f)\,x' + t_v(f)
  \end{equation}
with $u_v(f) \in K_v^\times$ and $r_v(f), s_v(f), t_v(f) \in K_v$, identifying $\mathcal{W}_v$ with the pullback of $\mathcal{W}'_v$ along $\varphi_f$. For all but finitely many finite places $v$, the coefficients lie in $\mathcal{O}_v$ and $u_v(f)$ is a unit in $\mathcal{O}_v$. The element $u_v(f)$ is termed the \emph{local dilation factor} of $f$ at $v$, and the associated factor
  \[
    c_v(f) \coloneqq u_v(f)^{-12}
  \]
renders the discriminant relation in the simple form
  \[
    \Delta_{\mathcal{W}',v} = c_v(f)\,\Delta_{\mathcal{W},v}.
  \]
\item[(ii)] A correspondence between the local data $\{\sigma_v\}$ and $\{\sigma'_v\}$ keeping the transformation of sections or level structures under $\varphi_f$ and the coordinate changes in~\eqref{eq:local-change-of-variables}.
\end{enumerate}

The labels $\lambda$ and $\lambda'$ are not constrained to match. In applications they often refer to distinct construction histories. They play a passive role in the inequalities below, and yet they remain faithful to the principle that one compares presentations rather than collapsing them. The category of elliptic packages over $K$ is denoted $\mathbf{EllPkg}_K$.
\end{definition}

\begin{remark}[Variable base fields]\label{rem:EllPkgVarCrossRef}
Definition~\ref{def:EllipticPackage} fixes the base field $K$, and the category $\mathbf{EllPkg}_K$ keeps morphisms over $K$. When one wishes to speak of Kummerian moves, such as adjoining an $\ell$-th root $K \hookrightarrow K(\sqrt[\ell]{q})$, then it is convenient to pass to the variable-base category $\mathbf{EllPkg}_{\mathrm{var}}$ of Definition~\ref{def:EllPkgVar} and to separate the model defect $\mu$ from the Kummer defect of Definition~\ref{def:KummerDefect}.
\end{remark}

\subsection{The model defect}\label{subsec:ModelDefect}

\begin{definition}[Model defect]\label{def:ModelDefect}
Let $f \colon \mathcal{E} \to \mathcal{E}'$ be a morphism in $\mathbf{EllPkg}_K$. The \emph{model defect} of $f$ is
\[
\mu(f) \coloneqq \frac{1}{[K:\Q]} \sum_{v} \log^{+}\!\bigl(\lvert c_v(f)\rvert_v^{-1}\bigr)
\]
where $c_v(f) \in K_v^\times$ is the factor of Definition~\ref{def:EllipticPackage}, the absolute values are normalized, the sum runs over all places $v$ of $K$, and $\log^{+}(x) \coloneqq \max\{\log x,0\}$ for $x>0$. The sum is finite since for all but finitely many finite places $v$ one finds $\lvert c_v(f)\rvert_v = 1$, so that the corresponding summand vanishes. For the identity morphism one has $\mu(\mathrm{id}_{\mathcal{E}}) = 0$, and for composable morphisms $f \colon \mathcal{E}_1 \to \mathcal{E}_2$ and $g \colon \mathcal{E}_2 \to \mathcal{E}_3$ one observes the sub-additivity
\[
\mu(g \circ f) \le \mu(f) + \mu(g).
\]
\end{definition}

\begin{proof}[Proof of sub-additivity in Definition~\ref{def:ModelDefect}]
For each place $v$ one has $c_v(g \circ f) = c_v(g)\,c_v(f)$. For positive real numbers $a$ and $b$ one has the elementary inequality $\log^{+}(ab) \le \log^{+}(a) + \log^{+}(b)$. Applying this with $a = \lvert c_v(f)\rvert_v^{-1}$ and $b = \lvert c_v(g)\rvert_v^{-1}$ and then summing over $v$ yields the stated inequality.
\end{proof}

\begin{remark}[Equivalent expressions for the model defect]\label{rem:ModelDefectEquivalent}
Let $f \colon \mathcal{E}\to\mathcal{E}'$ be a morphism in $\mathbf{EllPkg}_K$, with local dilation factors $u_v(f)$ and $c_v(f)=u_v(f)^{-12}$ as in Definition~\ref{def:EllipticPackage}. Then
\[
\mu(f)
=
\frac{1}{[K:\Q]} \sum_v \log^{+}\!\bigl(|u_v(f)|_v^{12}\bigr)
=
\frac{12}{[K:\Q]} \sum_v \log^{+}\!\bigl(|u_v(f)|_v\bigr).
\]
In particular, if $|u_v(f)|_v\le 1$ for all places $v$, then $\mu(f)=0$.
\end{remark}

\begin{remark}[Elliptic packages as a quantitative category]\label{rem:EllPkgQuantitative}
With the model defect $\mu$ in the role of a transport cost, the category $\mathbf{EllPkg}_K$ acquires the structure of a quantitative category of isomorphisms in the sense of Definition~\ref{def:QuantitativeCategory}. Two elliptic packages may be isomorphic in $\mathbf{EllPkg}_K$, and yet be linked by an isomorphism $f$ that carries a positive defect $\mu(f)$. This defect preserves the visibility of the arithmetic price paid to alter the presentation, which is the whole purpose of the intensional reading.
\end{remark}

\subsection{The discriminant height and its transport law}\label{subsec:DiscriminantHeight}

\begin{lemma}[Local change of discriminant]\label{lem:LocalDiscriminantChange}
Let $K_v$ be a local field, and let $\mathcal{W}_v$ and $\mathcal{W}'_v$ be two Weierstrass equations of an elliptic curve over $K_v$. Suppose that $\mathcal{W}'_v$ is obtained from $\mathcal{W}_v$ by a change of variables of the form
\[
  x = u^2 x' + r
  \qquad
  y = u^3 y' + u^2 s x' + t
\]
with $u \in K_v^{\times}$ and $r,s,t \in K_v$. Then their discriminants satisfy
\[
  \Delta(\mathcal{W}'_v) = u^{-12}\,\Delta(\mathcal{W}_v)
\]
in $K_v^\times$.
\end{lemma}

\begin{proof}
This is the customary transformation law for the discriminant under a change of Weierstrass coordinates. One may consult Silverman \cite[III, Proposition~1.1]{SilvermanAEC}.
\end{proof}

We fix the normalization of absolute values once, so that the definition below may be stated without further comment. For each place $v$ of $K$ we take $\lvert\,\cdot\,\rvert_v$ to be the normalized absolute value of the completion $K_v$, with the local degree carried inside the absolute value rather than left as a separate weight. At a finite place $v$ above the rational prime $p$, with normalized valuation $\mathrm{ord}_v$ taking the value $1$ at a uniformizer, this reads $\lvert x\rvert_v = q_v^{-\mathrm{ord}_v(x)}$, where $q_v = \#\kappa(v) = p^{\,f_v}$ is the residue-field cardinality and $f_v = [\kappa(v):\mathbb{F}_p]$ the residue degree. At a real place $\lvert x\rvert_v = \lvert\sigma_v(x)\rvert$ and at a complex place $\lvert x\rvert_v = \lvert\sigma_v(x)\rvert^{2}$, with $\sigma_v$ the corresponding embedding. The local degrees $n_v = [K_v:\Q_p]$ are thereby absorbed into the absolute values, the product formula $\prod_v \lvert x\rvert_v = 1$ holds for every $x\in K^\times$, and the factor $1/[K:\Q]$ in the height below renders the result an absolute height, independent of the field over which the package is presented.

\begin{definition}[Discriminant height on elliptic packages]\label{def:DiscriminantHeight}
Let $\mathcal{E} = (E, \mathcal{W}, \{\sigma_v\}, \lambda)$ be an elliptic package over $K$. The \emph{discriminant height} of $\mathcal{E}$ is
\[
h_\Delta(\mathcal{E}) \coloneqq \frac{1}{[K:\Q]} \sum_{v} \log \max\bigl(1, \lvert \Delta_{\mathcal{W},v} \rvert_v^{-1}\bigr)
\]
where $\Delta_{\mathcal{W},v}$ is the discriminant of the local model $\mathcal{W}_v$, the sum runs over all places $v$ of $K$, and the absolute values are normalized as above.
\end{definition}

In particular, the contribution of a finite place to the sum is
\[
  \log \max\bigl(1, \lvert \Delta_{\mathcal{W},v}\rvert_v^{-1}\bigr)
  = f_v \, \max\bigl(0,\, \mathrm{ord}_v(\Delta_{\mathcal{W}})\bigr)\,\log p
\]
so that the residue degree $f_v$ appears as the explicit local weight from the outset, in agreement with Remark~\ref{rem:HeightsAsLogIntData}. Unlike the conductor complexity of Definition~\ref{def:ConductorComplexity}, this quantity reads the chosen model $\mathcal{W}$ and therefore lives in Universe~A.

\begin{lemma}[Transport inequality for the discriminant height]\label{lem:TransportDiscriminant}
Let $K$ be a number field and let $f \colon \mathcal{E} \to \mathcal{E}'$ be a morphism in $\mathbf{EllPkg}_K$. Suppose that for each place $v$ of $K$ the discriminants satisfy $\Delta_{\mathcal{W}',v} = c_v(f)\,\Delta_{\mathcal{W},v}$. Then
\[
  h_\Delta(\mathcal{E}')
  \le
  h_\Delta(\mathcal{E}) + \mu(f).
\]
\end{lemma}

\begin{proof}
Fix a place $v$ and set $A_v \coloneqq |\Delta_{\mathcal{W},v}|_v^{-1}$ and $\gamma_v \coloneqq \log |c_v(f)|_v^{-1}$, so that $|\Delta_{\mathcal{W}',v}|_v^{-1} = e^{\gamma_v} A_v$. Writing $\alpha_v \coloneqq \log A_v$, one has
\[
\log \max(1, A_v) = \max\{0,\alpha_v\}
\qquad
\log \max(1, e^{\gamma_v}A_v) = \max\{0,\alpha_v+\gamma_v\}.
\]
The elementary inequality $\max\{0,x+y\} \le \max\{0,x\} + \max\{0,y\}$ yields, at each place,
\[
\log \max\bigl(1,\ |\Delta_{\mathcal{W}',v}|_v^{-1}\bigr)
\le
\log \max\bigl(1,\ |\Delta_{\mathcal{W},v}|_v^{-1}\bigr)
+
\log^{+}\!\bigl(|c_v(f)|_v^{-1}\bigr).
\]
Summing over $v$ and dividing by $[K:\Q]$ gives $h_\Delta(\mathcal{E}') \le h_\Delta(\mathcal{E}) + \mu(f)$, following Definition~\ref{def:ModelDefect}.
\end{proof}

\parahead{Non-arbitrariness of the transport cost} The model defect $\mu(f)$ enters as a cost attached to a chosen transport $f$. It is helpful to note that $\mu$ cannot be treated as a free parameter once the morphism $f$ is fixed, since any increase of the discriminant height forces a corresponding payment.

\begin{lemma}[Height increase forces defect]\label{lem:HeightIncreaseForcesDefect}
Let $f \colon \mathcal{E} \to \mathcal{E}'$ be a morphism in $\mathbf{EllPkg}_K$. If $h_\Delta(\mathcal{E}') \ge h_\Delta(\mathcal{E})$, then
\[
h_\Delta(\mathcal{E}') - h_\Delta(\mathcal{E}) \le \mu(f).
\]
In particular, if $h_\Delta(\mathcal{E}') - h_\Delta(\mathcal{E}) = D$ with $D>0$, then $\mu(f)\ge D$.
\end{lemma}

\begin{proof}
Lemma~\ref{lem:TransportDiscriminant} gives $h_\Delta(\mathcal{E}') \le h_\Delta(\mathcal{E}) + \mu(f)$, so the claim follows by rearranging the terms.
\end{proof}

\begin{corollary}[Two-sided control by a transport and its inverse]\label{cor:TwoSidedControlTransport}
Let $f \colon \mathcal{E}\to\mathcal{E}'$ be an isomorphism in $\mathbf{EllPkg}_K$. Then
\[
\bigl|h_\Delta(\mathcal{E}') - h_\Delta(\mathcal{E})\bigr|
\le
\mu(f) + \mu(f^{-1}).
\]
In particular, if $\mu(f)=\mu(f^{-1})=0$, then $h_\Delta(\mathcal{E}') = h_\Delta(\mathcal{E})$.
\end{corollary}

\begin{proof}
Apply Lemma~\ref{lem:TransportDiscriminant} to $f$ and to $f^{-1}$, then combine the two inequalities.
\end{proof}

\begin{remark}[A small interpretive comment]\label{rem:ZeroCostInterpretation}
When the morphism $f$ is fixed, the condition $\mu(f)=0$ expresses that the local coordinate changes carry no dilation cost, and Corollary~\ref{cor:TwoSidedControlTransport} then shows that the two presentations are indistinguishable as far as $h_\Delta$ is concerned. This is the arithmetic shadow of the extensional collapse of Definition~\ref{def:ExtensionalCollapse} -- a costless isomorphism is precisely one that survives the passage to Universe~H.
\end{remark}

\subsection{Keeping the coherence level visible}\label{subsec:EllPkgTwoOne}

In the parts of the theory where functoriality or the commutativity of squares is invoked, coherence data is not ancillary, and we work in a quantitative $(2,1)$-categorical enhancement of the $1$-category $\mathbf{EllPkg}_K$.

\begin{definition}[Quantitative $(2,1)$-category of elliptic packages]\label{def:EllPkgTwoOne}
Fix a number field $K$. A \emph{quantitative $(2,1)$-category of elliptic packages over $K$} is a $(2,1)$-category $\mathfrak{EllPkg}_K$ with the following features:
\begin{enumerate}
\item[\textup{(i)}] Objects are elliptic packages over $K$ in the sense of Definition~\ref{def:EllipticPackage}.
\item[\textup{(ii)}] $1$-morphisms are morphisms of elliptic packages, with the same underlying data as in $\mathbf{EllPkg}_K$.
\item[\textup{(iii)}] All $2$-morphisms are invertible. We write $f \cong g$ to mean that there exists an invertible $2$-morphism $f \Rightarrow g$.
\item[\textup{(iv)}] The defect $\mu$ of Definition~\ref{def:ModelDefect} extends to $1$-morphisms of $\mathfrak{EllPkg}_K$ and is invariant under $2$-isomorphism, so that $f \cong g$ implies $\mu(f)=\mu(g)$.
\end{enumerate}
The homotopy $1$-category $\Ho(\mathfrak{EllPkg}_K)$ may be identified with $\mathbf{EllPkg}_K$ by passing to $2$-isomorphism classes of $1$-morphisms.
\end{definition}

\begin{definition}[Commutativity up to an invertible $2$-morphism]\label{def:CommUpToDefect}
Let $\mathfrak{A}$ be a $(2,1)$-category. Given $1$-morphisms
\[
  A \xrightarrow{u} B \xrightarrow{v} C
  \qquad
  A \xrightarrow{u'} B' \xrightarrow{v'} C
\]
a \emph{comparison} is an invertible $2$-morphism $\alpha \colon v\circ u \Rightarrow v'\circ u'$. In this situation we say that the square \emph{commutes via $\alpha$}.
\end{definition}

\begin{definition}[Pseudonatural transformation in a $(2,1)$-category]\label{def:PseudoNatural}
Let $\mathfrak{A}$ be a $(2,1)$-category and let $F,G\colon\mathfrak{A}\to\mathfrak{A}$ be $2$-functors. A \emph{pseudonatural transformation} $\eta\colon F \Rightarrow G$ consists of the following data:
\begin{enumerate}
\item[\textup{(i)}] for every object $A$, a $1$-morphism $\eta_A\colon F(A)\to G(A)$,
\item[\textup{(ii)}] for every $1$-morphism $f\colon A\to B$, an invertible $2$-morphism $\eta_f \colon \eta_B\circ F(f)\Rightarrow G(f)\circ \eta_A$,
\end{enumerate}
subject to the unit and composition coherence axioms.
\end{definition}

\begin{remark}[Why pseudonaturality stays at the $(2,1)$-level]\label{rem:PseudoNaturalNotBicategorical}
Definitions~\ref{def:CommUpToDefect} and~\ref{def:PseudoNatural} are kept in their native $(2,1)$-categorical form, with coherence registered by explicit invertible $2$-morphisms. The quantitative structure sits on $1$-morphisms, by way of $\mu$, and Definition~\ref{def:EllPkgTwoOne}\textup{(iv)} requires $\mu$ to be invariant under $2$-isomorphism. In this manner the visibility of the $2$-morphisms prevents accidental identifications while the numerical accounting stays attached to the transports themselves.

It is worth saying where these $2$-cells exert force and where they fall silent, since the two are easily confused. They exert force at the level of derivation. The invariance of $\mu$ under $2$-isomorphism is what turns the cost into a quantity attached to the transport rather than to its presentation, and so it is what permits the cost to descend to the homotopy $1$-category by Proposition~\ref{prop:CoherenceStability}. Pseudonaturality plays the parallel role for transport along a morphism, since the invertible $2$-cell $(\tau_\varepsilon)_f$ of the pseudonatural comparison in Hypothesis~\ref{hyp:ReducedNormalizedProfile} renders the two routes around a normalization square equal in cost, so that the family of bounds coheres across $\mathbf{EllPkg}_K$ rather than being chosen object by object with no relation between the choices.

They fall silent at the level of the final statement. The normalized inequality is objectwise, and its derivation uses none of the pseudonaturality data, as Lemma~\ref{lem:ObjectwiseInsensitive} states. When the construction descends to Universe~P in Proposition~\ref{prop:SoundnessDescentNormalization}, the $2$-cells leave no residue. This silence does not mark the coherence level as idle. The descent is sound for the very reason that the cost was first shown invariant under the $2$-cells, so that the coherence data is the scaffolding by which the scalar is proven independent of its presentation. Once that independence is in hand, the scaffolding may be set aside for the arithmetical statement, in the conservative spirit of Proposition~\ref{prop:ConservativeExtension}.
\end{remark}

\begin{proposition}[Coherence stability of the cost]\label{prop:CoherenceStability}
Let $\mathfrak{A}$ be a $(2,1)$-category equipped with a cost $c \colon \mathrm{Mor}_1(\mathfrak{A}) \to \R_{\ge 0}$ satisfying the axioms of Definition~\ref{def:UniverseAFormal}\textup{(iii)}. Then the following hold:
\begin{enumerate}
\item[\textup{(i)}] \textit{Descent to the homotopy $1$-category.} The assignment $[f] \mapsto c(f)$ defines a map $\bar{c} \colon \mathrm{Mor}(\Ho(\mathfrak{A})) \to \R_{\ge 0}$ on the homotopy $1$-category of Definition~\ref{def:HoOneCategory}, satisfying $\bar{c}([\mathrm{id}_A]) = 0$ and the sub-additivity $\bar{c}([g]\circ[f]) \le \bar{c}([f]) + \bar{c}([g])$.
\item[\textup{(ii)}] \textit{Cost equality for $2$-isomorphic composites.} If $f, g \colon A \to B$ are $1$-morphisms and there exists an invertible $2$-morphism $\alpha \colon f \Rightarrow g$, then $c(f) = c(g)$.
\end{enumerate}
In consequence, no cost accrues from the $2$-morphisms that witness the commutativity of faces. Every accumulation of cost arises from the sub-additivity of $c$ applied to the $1$-morphisms that appear as factors of a composite.
\end{proposition}

\begin{proof}
For \textup{(i)}, the invariance of $c$ under $2$-isomorphism, which is part of Definition~\ref{def:EllPkgTwoOne}\textup{(iv)}, shows that $c$ is constant on $2$-isomorphism classes, so that $\bar{c}$ is well defined. The normalization and sub-additivity descend from the corresponding properties of $c$. Part \textup{(ii)} is the invariance under $2$-isomorphism restated for a single pair of parallel $1$-morphisms.
\end{proof}

\subsection{Normalization and the accumulated defect}\label{subsec:Normalization}

To compare an arbitrary presentation with the curve itself, one needs a canonical representative within each isomorphism class and a distinguished transport toward it. We package these choices into a normalization procedure.

\begin{definition}[Normalization defect]\label{def:AccumulatedDefect}
Let $K$ be a number field. A \emph{normalization procedure} for elliptic packages over $K$ comprises the following choices:
\begin{enumerate}
\item[(i)]\label{item:NF_rule} a rule $\mathrm{NF}_K$ assigning to every elliptic package $\mathcal{E}$ over $K$ a normalized elliptic package $\mathrm{NF}_K(\mathcal{E})$ defined over the same field $K$,
\item[(ii)]\label{item:NF_transport} for every $\mathcal{E}$, a distinguished morphism $\nu_{\mathcal{E}} \colon \mathrm{NF}_K(\mathcal{E}) \to \mathcal{E}$ in $\mathbf{EllPkg}_K$ whose local coordinate changes are produced by the normalization rule,
\item[(iii)]\label{item:NF_finitary} a finitary discipline, where the outputs $\mathrm{NF}_K(\mathcal{E})$ and $\nu_{\mathcal{E}}$ depend only upon the finitely presented data of $\mathcal{E}$, so that one has in mind a deterministic procedure that clears denominators, performs local minimization at the finite places, and resolves the remaining ambiguities by a fixed tie-breaker,
\item[(iv)]\label{item:NF_idempotent} an idempotence discipline, where the normalized output is a fixed point of the procedure and its distinguished transport is the strict identity, so that
\[
  \mathrm{NF}_K\bigl(\mathrm{NF}_K(\mathcal{E})\bigr) = \mathrm{NF}_K(\mathcal{E})
  \qquad\text{and}\qquad
  \nu_{\mathrm{NF}_K(\mathcal{E})} = \mathrm{id}_{\mathrm{NF}_K(\mathcal{E})},
\]
\item[(v)]\label{item:NF_curve_invariance} a curve invariance discipline, where for every $\mathcal{E} = (E,\mathcal{W},\{\sigma_v\}_{v\in S},\lambda)$ the normalized output $\mathrm{NF}_K(\mathcal{E})$ has the same underlying curve $E$ and the distinguished morphism $\nu_{\mathcal{E}}$ has underlying curve isomorphism $\varphi_{\nu_{\mathcal{E}}} = \mathrm{id}_E$.
\end{enumerate}
Given such a choice, the \emph{normalization defect} of $\mathcal{E}$ is
\[
\delta_{\mathrm{mod}}(\mathcal{E}) \coloneqq \mu(\nu_{\mathcal{E}})
\]
and one also writes $\operatorname{defect}(\mathcal{E}) \coloneqq \delta_{\mathrm{mod}}(\mathcal{E})$. In particular, the idempotence discipline~\ref{item:NF_idempotent} implies $\defect(\mathrm{NF}_K(\mathcal{E})) = 0$.
\end{definition}

\begin{corollary}[A lower bound for the normalization defect]\label{cor:DefectLowerBoundNormalization}
Let $\mathrm{NF}_K$ be a normalization procedure in the sense of Definition~\ref{def:AccumulatedDefect}. For every elliptic package $\mathcal{E}$ over $K$ one has
\[
\defect(\mathcal{E})
\ge
\max\Bigl\{0,\ h_\Delta(\mathcal{E}) - h_\Delta\bigl(\mathrm{NF}_K(\mathcal{E})\bigr)\Bigr\}.
\]
\end{corollary}

\begin{proof}
By Lemma~\ref{lem:TransportDiscriminant} applied to $\nu_{\mathcal{E}} \colon \mathrm{NF}_K(\mathcal{E}) \to \mathcal{E}$ one has $h_\Delta(\mathcal{E}) \le h_\Delta(\mathrm{NF}_K(\mathcal{E})) + \mu(\nu_{\mathcal{E}})$, and Definition~\ref{def:AccumulatedDefect} identifies $\mu(\nu_{\mathcal{E}})$ with $\defect(\mathcal{E})$.
\end{proof}

\begin{remark}[The deterministic nature of the normalization]\label{rem:DeterministicNormalizationDiscipline}
The finitary discipline of Definition~\ref{def:AccumulatedDefect}\textup{\ref{item:NF_finitary}} requires the normalization procedure to be deterministic in the sense of Universe~P. At each finite place $v$, Tate's algorithm proceeds through a fixed sequence of conditional branches, each applying a coordinate change with parameters $(u_v,r_v,s_v,t_v)$ drawn from $\mathcal{O}_{K_v}$. When the algorithm permits more than one choice of local coordinate, then a tie-breaking rule must be specified. One fixes, once and for all, the following convention. Every local parameter $\alpha \in \mathcal{O}_{K_v}$ is expressed as its unique tuple of integers $(z_1, \dots, z_d) \in \Z^d$ in the \textit{certified integral basis} $(\omega_1, \dots, \omega_d)$ of the field code (Definition~\ref{def:CertifiedFieldCode}), and among all coordinate parameters yielding a minimal model at $v$, the algorithm selects the one whose combined integer vector is lexicographically smallest. This renders the output a typographically unique string of integers determined by the input data alone. The global normalization $\mathrm{NF}_K$ applies the local minimization at each prime dividing the discriminant in increasing order and retains the archimedean model unchanged, so that the idempotence discipline of Definition~\ref{def:AccumulatedDefect}\textup{\ref{item:NF_idempotent}} follows from the lexicographic rule selecting the same representative on a fixed point.
\end{remark}

%------------------------------------------------------------------
\subsection{A dilation tower at the bad prime}\label{subsec:DilationTower}
%------------------------------------------------------------------

Before the discriminant height becomes the term of an inequality, it is instructive to see how it moves along an isomorphism of packages. The two-sided control of Corollary~\ref{cor:TwoSidedControlTransport} permits such movement, and the tower below shows that the movement is unbounded. This is the arithmetic face of the distinction between the static and the dynamic disciplines of Section~\ref{sec:identification}. It is also what determines, in Section~\ref{sec:IntensionalSzpiro}, that a scalar comparable to the discriminant height is read as a quantity transforming under transport rather than as an invariant of an isomorphism class.

\begin{definition}[Dilation tower at the bad prime]\label{def:DilationTowerAtBadPrime}
Let
\[
  E \colon y^2+y=x^3-x
\]
over $\Q$, and let $\mathcal{E}_0$ be the elliptic package whose Weierstrass model is
\[
  \mathcal{W}_0 \colon y^2+y=x^3-x.
\]
For each integer $m\ge 0$, let $\mathcal{W}_m$ be the Weierstrass model obtained from $\mathcal{W}_0$ by the change of variables
\[
  x=37^{-2m}x'
  \qquad
  y=37^{-3m}y'.
\]
Equivalently,
\[
  \mathcal{W}_m \colon y'^2+37^{3m}y'=x'^3-37^{4m}x'.
\]
Let $\mathcal{E}_m$ be the elliptic package obtained from $\mathcal{E}_0$ by replacing the model with $\mathcal{W}_m$ and by transporting the remaining local data along the same coordinate change. We write
\[
  \rho_m \colon \mathcal{E}_0 \longrightarrow \mathcal{E}_m
\]
for the induced morphism of elliptic packages.
\end{definition}

\begin{lemma}[Growth along the dilation tower]\label{lem:DilationTowerGrowth}
For every integer $m\ge 0$, the morphism
\[
  \rho_m \colon \mathcal{E}_0 \longrightarrow \mathcal{E}_m
\]
is an isomorphism in $\mathbf{EllPkg}_{\Q}$, and one has
\[
  h_\Delta(\mathcal{E}_m)=(12m+1)\log 37.
\]
Moreover,
\[
  \mu(\rho_m)=12m\log 37.
\]
\end{lemma}

\begin{proof}
The coordinate change in Definition~\ref{def:DilationTowerAtBadPrime} has local dilation factor
\[
  u=37^{-m}
\]
at the prime $37$. Hence the discriminant transformation law of Lemma~\ref{lem:LocalDiscriminantChange} gives
\[
  \Delta_{\mathcal{W}_m}=u^{-12}\Delta_{\mathcal{W}_0}
  =
  37^{12m}\cdot 37
  =
  37^{12m+1}.
\]
There are no new finite prime divisors. Thus the finite contribution to $h_\Delta(\mathcal{E}_m)$ is
\[
  \log\max\bigl(1,\lvert 37^{12m+1}\rvert_{37}^{-1}\bigr)
  =
  (12m+1)\log 37.
\]
The archimedean contribution vanishes for this height, since $\lvert 37^{12m+1}\rvert_\infty^{-1}<1$. This proves the displayed formula for $h_\Delta(\mathcal{E}_m)$.

The morphism $\rho_m$ is invertible, since the coordinate change is invertible over $\Q$. Its model defect is computed from
\[
  c_{37}(\rho_m)=u^{-12}=37^{12m}.
\]
Therefore
\[
  \mu(\rho_m)
  =
  \log^{+}\bigl(\lvert 37^{12m}\rvert_{37}^{-1}\bigr)
  =
  12m\log 37.
\]
\end{proof}

\begin{proposition}[No isomorphism-invariant scalar comparable to the discriminant height]
\label{prop:NoInvariantLogVolume}
There is no assignment
\[
  V \colon \operatorname{Ob}(\mathbf{EllPkg}_{\Q})\longrightarrow \R_{\ge 0}
\]
satisfying both of the following conditions:
\begin{enumerate}
\item[\textup{(i)}] If two packages are isomorphic in $\mathbf{EllPkg}_{\Q}$, then their values under $V$ are equal.
\item[\textup{(ii)}] There exist constants $A,B>0$ such that for every elliptic package $\mathcal{E}$ over $\Q$ one has
\[
  A^{-1}h_\Delta(\mathcal{E})-B
  \le
  V(\mathcal{E})
  \le
  Ah_\Delta(\mathcal{E})+B.
\]
\end{enumerate}
\end{proposition}

\begin{proof}
Apply the two conditions to the packages $\mathcal{E}_m$ of Definition~\ref{def:DilationTowerAtBadPrime}. By Lemma~\ref{lem:DilationTowerGrowth}, each $\mathcal{E}_m$ is isomorphic to $\mathcal{E}_0$. Hence condition \textup{(i)} gives $V(\mathcal{E}_m)=V(\mathcal{E}_0)$ for all $m\ge 0$. The lower bound in condition \textup{(ii)} gives
\[
  V(\mathcal{E}_0)
  =
  V(\mathcal{E}_m)
  \ge
  A^{-1}(12m+1)\log 37-B
\]
for all $m\ge 0$. The right member tends to infinity with $m$, while the left member is fixed. This is impossible.
\end{proof}

\begin{remark}[The meaning of the obstruction]\label{rem:MeaningOfLogVolumeObstruction}
The obstruction marks the expected cost of keeping positive-cost isomorphisms, and one reads it as a feature of the discipline of Universe~A. The morphisms $\rho_m$ of Definition~\ref{def:DilationTowerAtBadPrime} are precisely the kind of transports that the universe permits. They are isomorphisms of packages, and yet they alter the discriminant height by a visible amount. A scalar comparable to $h_\Delta$ must see this alteration, and so it cannot at the same time be invariant under all such isomorphisms.
\end{remark}

%==================================================================
\section{The Kummerian cost protocol}\label{sec:Kummer}
%==================================================================

Definition~\ref{def:EllipticPackage} fixes the base field, and yet many arithmetic operations move between fields. Adjoining a root, passing to a splitting field, or extending scalars to trivialize a torsion structure all replace $K$ by a larger field $L$. In Universe~H such a passage is often treated as free once a base has been fixed. In Universe~A it carries a cost of its own, and the present section isolates that cost and calibrates it against the ramification of the extension.

\subsection{A variable-base category}\label{subsec:EllPkgVar}

\begin{definition}[A variable-base category of elliptic packages]\label{def:EllPkgVar}
Let $\mathbf{EllPkg}_{K}$ be as in Definition~\ref{def:EllipticPackage}. Define a category $\mathbf{EllPkg}_{\mathrm{var}}$ as follows. An object of $\mathbf{EllPkg}_{\mathrm{var}}$ is a pair $(K,\mathcal{E})$ where $K$ is a number field and $\mathcal{E}$ is an elliptic package over $K$. A morphism
\[
  f \colon (K,\mathcal{E}) \longrightarrow (L,\mathcal{E}')
\]
consists of an inclusion of number fields $\iota \colon K \hookrightarrow L$ together with a morphism of elliptic packages over $L$
\[
  f_L \colon \iota_\ast(\mathcal{E}) \longrightarrow \mathcal{E}'
\]
where $\iota_\ast(\mathcal{E})$ is the package obtained from $\mathcal{E}$ by base change along $\iota$, with the same underlying Weierstrass equation viewed over $\mathcal{O}_L$ and with the same label transported as a symbol.
\end{definition}

\begin{remark}[Why the variable base is introduced]\label{rem:WhyEllPkgVar}
The category $\mathbf{EllPkg}_{\mathrm{var}}$ is not a replacement for $\mathbf{EllPkg}_K$. Its role is to provide a place where one may speak of a morphism that witnesses a passage $K \leadsto L$. In Universe~A the morphism $f\colon (K,\mathcal{E}) \to (L,\mathcal{E}_L)$ may carry a strictly positive distortion cost $\mu(f) > 0$, reflecting the arithmetic price of the field extension. This term enters the intensional Szpiro inequality of Section~\ref{sec:IntensionalSzpiro} through the detachment defect introduced below.
\end{remark}

\subsection{Radical transports and the Kummer defect}\label{subsec:KummerDefect}

We isolate the morphisms that correspond to root extraction.

\begin{definition}[Radical transport morphism]\label{def:RadicalTransport}
Let $(K,\mathcal{E})$ be an object of $\mathbf{EllPkg}_{\mathrm{var}}$. Fix an integer $\ell\ge 2$ and an element $q\in K^\times$, and let $L = K(\sqrt[\ell]{q})$. A \emph{radical transport morphism} of degree $\ell$ is a morphism
\[
  f_{\mathrm{rad}} \colon (K,\mathcal{E}) \longrightarrow (L,\mathcal{E}_L)
\]
in $\mathbf{EllPkg}_{\mathrm{var}}$ for which the underlying package $\mathcal{E}_L$ is the base change of $\mathcal{E}$ to $L$, and for which the Weierstrass coordinate changes in the package part are trivial at every place. In particular, the model defect contribution coming from the local dilation factors vanishes.
\end{definition}

The classical part of the cost sees nothing through such a morphism, so that one needs a separate term to keep the detachment visible.

\begin{definition}[Kummer degree and Kummer defect]\label{def:KummerDefect}
Let $f \colon (K,\mathcal{E}) \to (L,\mathcal{E}')$ be a morphism in $\mathbf{EllPkg}_{\mathrm{var}}$ with underlying field inclusion $\iota \colon K \hookrightarrow L$. The \emph{Kummer degree} of $f$ is
\[
  \deg_{\mathrm{kum}}(f) \coloneqq [L:K]
\]
and, having fixed a constant $C_{\mathrm{kum}}\ge 0$ depending only upon the chosen protocol, the \emph{Kummer defect} of $f$ is
\[
  \mu_{\mathrm{kum}}(f) \coloneqq C_{\mathrm{kum}} \cdot \log\bigl(\deg_{\mathrm{kum}}(f)\bigr).
\]
\end{definition}

\begin{remark}[Total transport cost]\label{rem:TotalTransportCost}
For morphisms internal to $\mathbf{EllPkg}_{K}$, the model defect $\mu$ of Definition~\ref{def:ModelDefect} is the operative cost. For morphisms in $\mathbf{EllPkg}_{\mathrm{var}}$ it is convenient to keep two layers visible, and one writes
\[
  \mu_{\mathrm{tot}}(f) \coloneqq \mu_{\mathrm{mod}}(f) + \mu_{\mathrm{kum}}(f)
\]
where $\mu_{\mathrm{mod}}(f)$ is the model defect computed from the local dilation factors in the package component of $f$, and $\mu_{\mathrm{kum}}(f)$ is the detachment term of Definition~\ref{def:KummerDefect}. A radical transport in the sense of Definition~\ref{def:RadicalTransport} satisfies $\mu_{\mathrm{mod}}(f_{\mathrm{rad}})=0$, while $\mu_{\mathrm{kum}}(f_{\mathrm{rad}})$ may be non-zero.
\end{remark}

\begin{lemma}[Cost of Kummer detachment]\label{lem:KummerDetachmentCost}
Let $f_{\mathrm{rad}}$ be a radical transport morphism of degree $\ell$ in the sense of Definition~\ref{def:RadicalTransport}. Then
\[
  \mu_{\mathrm{kum}}(f_{\mathrm{rad}})
  \le
  C_{\mathrm{kum}} \cdot \log(\ell).
\]
\end{lemma}

\begin{proof}
Let $L=K(\sqrt[\ell]{q})$, so that $[L:K]\le \ell$. By Definition~\ref{def:KummerDefect} one then has $\mu_{\mathrm{kum}}(f_{\mathrm{rad}}) = C_{\mathrm{kum}} \cdot \log([L:K]) \le C_{\mathrm{kum}} \cdot \log(\ell)$.
\end{proof}

\begin{remark}[Arithmetic calibration of the logarithmic scaling]\label{rem:KummerDefectArithmeticCalibration}
The choice of a logarithmic dependence on $\ell$ in Definition~\ref{def:KummerDefect} may appear axiomatic at first reading, and yet it admits a natural justification from the arithmetic of ramification in Kummer extensions, which it is useful to indicate. Let $L = K(\sqrt[\ell]{q})$ with $[L:K] \leq \ell$. The extension $L/K$ is abelian of exponent dividing $\ell$, and its ramification is confined to two families of places of $K$: those dividing $q$ and those dividing the rational prime $\ell$. At all other places, $L/K$ is unramified.

It is the relative different $\mathfrak{D}_{L/K}$ that registers this ramification place by place, and the relative discriminant is its norm, $\mathfrak{d}_{L/K}=N_{L/K}(\mathfrak{D}_{L/K})$, so that the exponent of $\mathfrak{D}_{L/K}$ at a place $w$ of $L$ above $v$ is the quantity that fixes the ceiling. For a place $w \mid v$ with ramification index $e_{w/v}$, the exponent $d_{w/v}$ of $\mathfrak{D}_{L/K}$ obeys the bounds of Serre~\cite{Serre1979}. When $v \nmid \ell$ the ramification is tame and the exponent is exactly $d_{w/v}=e_{w/v}-1\le \ell-1$. When $v \mid \ell$ the ramification is wild and the exponent rises, subject to
\[
  e_{w/v} \le d_{w/v} \le e_{w/v}-1+v_w(e_{w/v})
\]
where $v_w$ is the normalized valuation of $L$ at $w$. For the cyclic extension of prime degree $\ell$ the ramified places $w \mid \ell$ are totally ramified, so that $e_{w/v}=\ell$ and $v_w(\ell)=\ell\, e_{v/\ell}$, whence $d_{w/v}\le (\ell-1)+\ell\, e_{v/\ell}$, a quantity linear in $\ell$. It is this linear growth of the wild different exponent, weighted by the residual contribution $f_{w/v}\log\mathrm{N}v$ and then normalized by the degree $[L:\Q]$, that provokes a logarithmic ceiling rather than a linear one. The place-by-place reckoning that follows realizes these bounds, the tame places first and then the wild.

At a place $v$ of $K$ dividing $q$ but not $\ell$, the extension is at most tamely ramified. The local discriminant exponent at such a place is bounded by $e_{w/v} - 1 \leq \ell - 1$, where $e_{w/v}$ is the ramification index of a place $w$ of $L$ above $v$. After normalization by $[L:\Q] \geq [K:\Q]$, the total contribution from the tame places remains bounded independently of $\ell$, since the factor $(\ell - 1) / [L:K]$ does not exceed $1$ and the number and residual degrees of the places dividing $q$ depend only upon $K$.

At a place $v \mid \ell$, wild ramification may occur. The conductor-discriminant formula for abelian extensions expresses the discriminant exponent at $v$ in terms of the Artin conductor exponents of the non-trivial characters of $\mathrm{Gal}(L/K)$. For a cyclic extension of prime degree $\ell$, these $\ell - 1$ characters share a common conductor exponent $\mathfrak{f}_v$ at $v$, and the discriminant exponent is $(\ell - 1) \cdot \mathfrak{f}_v$. The quantity $\mathfrak{f}_v$ is bounded above in terms of the local ramification data of $K$ at $v$ -- that is, in terms of the ramification index $e_{v/\ell}$ and the residue degree $f_{v/\ell}$ of $v$ over the rational prime $\ell$ (see Neukirch~\cite{Neukirch1999} or Serre~\cite{Serre1979} for the relevant local and global theory). Since $\sum_{v \mid \ell} e_{v/\ell} \, f_{v/\ell} = [K:\Q]$ and each $\log N_{K/\Q}(v) = f_{v/\ell} \cdot \log \ell$, the total wild contribution to the logarithmic relative discriminant is of order $\ell \cdot \log \ell$ before normalization. After dividing by $[L:\Q]$, one obtains
\begin{equation}\label{eq:KummerRelativeDiscriminant}
  \frac{1}{[L:\Q]}\,\log N_{K/\Q}(\mathfrak{d}_{L/K})
  \leq
  C_{K,q} \cdot \log \ell
\end{equation}
for a constant $C_{K,q}$ depending only upon $K$ and upon the places of $K$ dividing $q$.

This bound grounds the logarithmic scaling of $\mu_{\mathrm{kum}}$ in Definition~\ref{def:KummerDefect}. When one descends a Szpiro-type bound from $L$ to $K$ -- as in the controlled descent of Hypothesis~\ref{hyp:ControlledDescentWithSupport} below -- then the price of the descent is governed by the relative discriminant of $L/K$ through the standard height comparison machinery. The wild contribution from the places above $\ell$ generates the $\log \ell$ term, while the tame contribution from the places dividing $q$ enters the full descent cost through a separate support term, whose necessity is shown in Section~\ref{subsec:ControlledDescentUniformity}. The scaling at issue here concerns the detachment cost $\mu_{\mathrm{kum}}$ alone, which sees only the degree $\ell$. A linear dependence $\mu_{\mathrm{kum}} \sim \ell$ would overestimate this cost by a factor of $\ell / \log \ell$, while a sub-logarithmic scaling would leave the wild ramification at the places above $\ell$ inadequately accounted for. The logarithmic regime is the point of equilibrium between the growth of the discriminant exponent and the normalization by the degree of the extension.
\end{remark}

\subsection{Descent back to the base field}\label{subsec:ControlledDescentBridge}

To compare a bound proved over a Kummer extension $L/K$ with a statement over the original field $K$, one needs a descent step. In the variable-base category we isolate this as a transport whose cost carries two terms of distinct provenance, one logarithmic in the Kummer degree and one sensitive to the tame support of the radicand. The necessity of the second term, and the explicit form of both, are taken up in Section~\ref{subsec:ControlledDescentUniformity}.

\begin{definition}[Package descent morphism]\label{def:DescentMorphism}
Let $L/K$ be a finite extension of number fields and let $\mathcal{E}_L$ be an elliptic package over $L$. The \emph{descent morphism} $f_{\mathrm{desc}} \colon \mathcal{E}_L \to \mathcal{E}_{L\mid K}$ is the morphism in $\mathbf{EllPkg}_{\mathrm{var}}$ defined by two operations. First, by scalar restriction, where the coefficients of the Weierstrass model are viewed over the maximal subfield $K$. Second, by re-minimization, where a normalization $\mathrm{NF}_K$ is applied to the restricted coefficients at each place $v$ of $K$ following Tate's algorithm. The dilation factors $u_v$ of $f_{\mathrm{desc}}$ keep the ratio between the $L$-minimal discriminant and the re-normalized $K$-minimal discriminant.
\end{definition}

\begin{definition}[Controlled descent datum]\label{def:ControlledDescentDatum}
Let $K$ be a number field, let $L/K$ be a finite extension, let $\mathcal{E}$ be an elliptic package over $K$, and write $\mathcal{E}_L$ for its base change to $L$. A \emph{controlled descent datum} for $(\mathcal{E},L/K)$ is a morphism in $\mathbf{EllPkg}_{\mathrm{var}}$
\begin{equation}\label{eq:ControlledDescentMorphism}
  \delta_{\mathcal{E},L/K} \colon (L,\mathrm{NF}_L(\mathcal{E}_L))
  \longrightarrow (K,\mathrm{NF}_K(\mathcal{E})).
\end{equation}
Its \emph{descent cost} is $\mu_{\mathrm{tot}}(\delta_{\mathcal{E},L/K})$.
\end{definition}

\begin{definition}[Tame Kummer support]\label{def:TameKummerSupport}
Let $K$ be a number field and let
\[
  L=K(\sqrt[\ell]{q})
\]
with $\ell\ge 2$ and $q\in K^{\times}$. The \emph{tame Kummer support} of $(q,\ell)$ is
\[
  \operatorname{Supp}^{\mathrm{tame}}_{K}(q,\ell)
  =
  \bigl\{
    v<\infty
    \mid
    v\nmid \ell
    \text{ and }
    v(q)\ne 0
  \bigr\}.
\]
Its logarithmic size is
\[
  R_{K}(q,\ell)
  \coloneqq
  \frac{1}{[K:\Q]}
  \sum_{v\in \operatorname{Supp}^{\mathrm{tame}}_{K}(q,\ell)}
  \log \mathrm{N}v
\]
where $\mathrm{N}v$ denotes the absolute norm of the prime ideal $v$.
\end{definition}

\begin{remark}[A harmless overcount]\label{rem:TameSupportOvercount}
The set of Definition~\ref{def:TameKummerSupport} may overcount the ramified tame places, since a valuation $v(q)$ divisible by $\ell$ can disappear after the removal of an $\ell$-th power. This overcount is useful. It is computable directly from the finite code of $q$, and it gives an upper bound for the tame ramification places of $L/K$ outside the primes above $\ell$.
\end{remark}

\begin{hypothesis}[Controlled descent with tame support]\label{hyp:ControlledDescentWithSupport}
Fix a number field $K$. There exist constants
\[
  C_{\mathrm{desc}}>0
  \qquad
  A_{\mathrm{desc}}\ge 0
  \qquad
  B_{\mathrm{desc},K}\ge 0
\]
such that the following holds. For every elliptic package $\mathcal{E}$ over $K$, and for every Kummer extension $L=K(\sqrt[\ell]{q})$, there exists a controlled descent datum $\delta_{\mathcal{E},L/K}$ in the sense of Definition~\ref{def:ControlledDescentDatum} satisfying
\begin{equation}\label{eq:ControlledDescentWithSupport}
  \mu_{\mathrm{tot}}(\delta_{\mathcal{E},L/K})
  \le
  C_{\mathrm{desc}}\log \ell
  +
  A_{\mathrm{desc}}\,R_{K}(q,\ell)
  +
  B_{\mathrm{desc},K}.
\end{equation}
\end{hypothesis}

\begin{remark}[The descent cost belongs to finite-Galois arithmetic]\label{rem:DescentCostClassicalNature}
Hypothesis~\ref{hyp:ControlledDescentWithSupport} occupies a distinguished position among the assumptions of this paper. The controlled descent datum $\delta_{\mathcal{E},L/K}$ of Definition~\ref{def:ControlledDescentDatum} involves three operations from the standard repertoire of algebraic number theory: scalar restriction of Weierstrass coefficients, re-minimization by way of Tate's algorithm at each place of $K$, and the computation of the dilation factors $u_v$ arising from the re-minimization. Each of these belongs to finite-Galois arithmetic. The two coefficients of the bound draw upon distinct sources of ramification. The term $C_{\mathrm{desc}}\log \ell$ collects the wild contribution at the places above $\ell$, read through the local theory of different exponents (see Serre~\cite{Serre1979} and Neukirch~\cite{Neukirch1999}), together with the Kummer detachment cost $C_{\mathrm{kum}}$ of Lemma~\ref{lem:KummerDetachmentCost}. The term $A_{\mathrm{desc}}R_{K}(q,\ell)$ collects the tame contribution at the places in the support of the radicand, of Definition~\ref{def:TameKummerSupport}. In this sense the descent cost is not an exotic quantity, and the bound is a calibrated restatement of a familiar ramification estimate within the language of transport costs. That both of its terms are forced, and that both may be made explicit away from the primes above $2$ and $3$, is shown in Section~\ref{subsec:ControlledDescentUniformity}.
\end{remark}

\begin{lemma}[Bridge inequality for the discriminant height]\label{lem:BridgeDescentHeight}
Let $K$ be a number field and assume Hypothesis~\ref{hyp:ControlledDescentWithSupport}. Let $\mathcal{E}$ be an elliptic package over $K$ and let $L=K(\sqrt[\ell]{q})$ be a Kummer extension, with $\mathcal{E}_L$ the base change of $\mathcal{E}$ to $L$. Then
\[
  h_\Delta\bigl(\mathrm{NF}_K(\mathcal{E})\bigr)
  \le
  h_\Delta\bigl(\mathrm{NF}_L(\mathcal{E}_L)\bigr)
  + C_{\mathrm{desc}}\log \ell
  + A_{\mathrm{desc}}\,R_{K}(q,\ell)
  + B_{\mathrm{desc},K}.
\]
\end{lemma}

\begin{proof}
By Hypothesis~\ref{hyp:ControlledDescentWithSupport} choose a controlled descent datum $\delta_{\mathcal{E},L/K} \colon (L,\mathrm{NF}_L(\mathcal{E}_L)) \to (K,\mathrm{NF}_K(\mathcal{E}))$ with $\mu_{\mathrm{tot}}(\delta_{\mathcal{E},L/K}) \le C_{\mathrm{desc}}\log \ell + A_{\mathrm{desc}}R_{K}(q,\ell) + B_{\mathrm{desc},K}$. The transport inequality of Lemma~\ref{lem:TransportDiscriminant}, applied to the package component of $\delta_{\mathcal{E},L/K}$, gives
\[
  h_\Delta\bigl(\mathrm{NF}_K(\mathcal{E})\bigr)
  \le
  h_\Delta\bigl(\mathrm{NF}_L(\mathcal{E}_L)\bigr)
  + \mu_{\mathrm{mod}}(\delta_{\mathcal{E},L/K})
  \le
  h_\Delta\bigl(\mathrm{NF}_L(\mathcal{E}_L)\bigr)
  + \mu_{\mathrm{tot}}(\delta_{\mathcal{E},L/K})
\]
since $\mu_{\mathrm{mod}} \le \mu_{\mathrm{tot}}$ by Remark~\ref{rem:TotalTransportCost}. Substituting the bound on the descent cost yields the claim.
\end{proof}

\subsection{The necessity of the tame support term and explicit local bounds}\label{subsec:ControlledDescentUniformity}

The descent estimate of Hypothesis~\ref{hyp:ControlledDescentWithSupport} carries two terms, and it is worth seeing why each is present. This subsection establishes that the tame support term is forced -- a bound logarithmic in the Kummer degree alone cannot hold uniformly as the radicand varies -- and it then proves the estimate explicitly away from the residue characteristics $2$ and $3$, where the bound holds with $C_{\mathrm{desc}}=12+C_{\mathrm{kum}}$ and $A_{\mathrm{desc}}=12$. The part that follows from local minimization is thereby separated from the part that must see the tame support of $q$.

There is also a small directional matter. The arrows of Definition~\ref{def:EllPkgVar} follow inclusions of fields from the smaller field to the larger one, while the descent arrow of Definition~\ref{def:ControlledDescentDatum} points in the reverse direction. For the estimates below, this reverse arrow is read as a finite descent code in a descent completion of $\mathbf{EllPkg}_{\mathrm{var}}$. This completion has the same objects as $\mathbf{EllPkg}_{\mathrm{var}}$. In addition to the base-change arrows, it carries formal reverse descent arrows whose data are finite coordinate changes produced by scalar restriction and re-minimization.

\begin{definition}[Height-admissible descent datum]\label{def:HeightAdmissibleDescentDatum}
Let $L/K$ be a finite extension of number fields, and let $\mathcal{E}$ be an elliptic package over $K$. A \emph{height-admissible descent datum} for $(\mathcal{E},L/K)$ is a controlled descent datum
\[
  \delta_{\mathcal{E},L/K}
  \colon
  (L,\mathrm{NF}_{L}(\mathcal{E}_{L}))
  \longrightarrow
  (K,\mathrm{NF}_{K}(\mathcal{E}))
\]
in the sense of Definition~\ref{def:ControlledDescentDatum}, read as an arrow of the descent completion described above, such that the transport inequality
\[
  h_{\Delta}\bigl(\mathrm{NF}_{K}(\mathcal{E})\bigr)
  \le
  h_{\Delta}\bigl(\mathrm{NF}_{L}(\mathcal{E}_{L})\bigr)
  +
  \mu_{\mathrm{tot}}(\delta_{\mathcal{E},L/K})
\]
holds.
\end{definition}

\begin{remark}[Finite nature of descent codes]\label{rem:FiniteDescentCodes}
A height-admissible descent datum is a finite object. It consists of a certified extension code for $L/K$, local coordinate changes at the finitely many places where the base-changed normalized model ceases to be minimal, and the finite list of local dilation exponents that enter the model defect. Thus the descent completion does not enlarge the foundational strength of the framework. It gives a typed place for the reverse arrows that are computed by the finite procedures of Definition~\ref{def:EllPkgCode} and Remark~\ref{rem:PkgCodeComputability}.
\end{remark}

\subsubsection{The tame support term is forced}\label{subsubsec:QuadraticTwistObstruction}

The first result shows that the coefficient $A_{\mathrm{desc}}$ of Hypothesis~\ref{hyp:ControlledDescentWithSupport} cannot be taken to be zero. A quadratic twist family lets a single tame prime grow without bound while the Kummer degree stays fixed at two.

\begin{definition}[Quadratic twist descent family]\label{def:QuadraticTwistDescentFamily}
For each rational prime $p\ge 5$, let
\[
  E_{p}\colon y^{2}=x^{3}-p^{2}x
\]
over $\Q$. Let $\mathcal{E}_{p}$ be the elliptic package over $\Q$ whose Weierstrass model is this equation, with the remaining package data chosen by the fixed deterministic conventions of Definition~\ref{def:AccumulatedDefect}. Let
\[
  L_{p}=\Q(\sqrt p)
\]
and write $t_{p}=\sqrt p$. Over $L_{p}$, the change of variables
\[
  x=t_{p}^{2}X
  \qquad
  y=t_{p}^{3}Y
\]
identifies $E_{p}$ with the curve
\[
  E_{0}\colon Y^{2}=X^{3}-X.
\]
\end{definition}

\begin{lemma}[Height gap in the twist family]\label{lem:QuadraticTwistHeightGap}
For the family of Definition~\ref{def:QuadraticTwistDescentFamily}, one has
\[
  h_{\Delta}\bigl(\mathrm{NF}_{\Q}(\mathcal{E}_{p})\bigr)
  -
  h_{\Delta}\bigl(\mathrm{NF}_{L_{p}}((\mathcal{E}_{p})_{L_{p}})\bigr)
  \ge
  6\log p-6\log 2.
\]
\end{lemma}

\begin{proof}
The discriminant of the equation $y^{2}=x^{3}-p^{2}x$ is $\Delta(E_{p})=64p^{6}$. At the prime $p\ge 5$, the displayed equation is minimal. Indeed, for a short Weierstrass equation in residue characteristic at least $5$, the condition $v_{p}(-p^{2})=2<4$ prevents a scaling by $p^{-1}$ from preserving integrality. Thus the $p$-part of $h_{\Delta}(\mathrm{NF}_{\Q}(\mathcal{E}_{p}))$ is exactly $6\log p$, and all other finite contributions are non-negative. Hence
\[
  h_{\Delta}\bigl(\mathrm{NF}_{\Q}(\mathcal{E}_{p})\bigr)\ge 6\log p.
\]

Over $L_{p}$, Definition~\ref{def:QuadraticTwistDescentFamily} identifies $(E_{p})_{L_{p}}$ with $E_{0}$. The model $Y^{2}=X^{3}-X$ has discriminant $64$. Therefore its normalized finite discriminant contribution over $L_{p}$ is at most $6\log 2$. Since normalization can only decrease the discriminant height, one has
\[
  h_{\Delta}\bigl(\mathrm{NF}_{L_{p}}((\mathcal{E}_{p})_{L_{p}})\bigr)
  \le
  6\log 2.
\]
The two estimates give the claim.
\end{proof}

\begin{theorem}[The support term cannot be dropped]\label{thm:NoUniformControlledDescent}
There are no constants $C>0$ and $B\ge 0$ with the following property. For every prime $p\ge 5$, there exists a height-admissible descent datum
\[
  \delta_{p}
  \colon
  (L_{p},\mathrm{NF}_{L_{p}}((\mathcal{E}_{p})_{L_{p}}))
  \longrightarrow
  (\Q,\mathrm{NF}_{\Q}(\mathcal{E}_{p}))
\]
satisfying $\mu_{\mathrm{tot}}(\delta_{p})\le C\log 2+B$. Consequently a descent bound of the form $C_{\mathrm{desc}}\log \ell+B_{\mathrm{desc},K}$, carrying no tame support term, fails to hold uniformly over $K=\Q$ as the radicand $q$ varies, and the coefficient $A_{\mathrm{desc}}$ of Hypothesis~\ref{hyp:ControlledDescentWithSupport} is strictly positive.
\end{theorem}

\begin{proof}
Assume such constants $C$ and $B$ have been fixed. Since $\delta_{p}$ is height-admissible, Definition~\ref{def:HeightAdmissibleDescentDatum} gives
\[
  h_{\Delta}\bigl(\mathrm{NF}_{\Q}(\mathcal{E}_{p})\bigr)
  -
  h_{\Delta}\bigl(\mathrm{NF}_{L_{p}}((\mathcal{E}_{p})_{L_{p}})\bigr)
  \le
  \mu_{\mathrm{tot}}(\delta_{p}).
\]
By Lemma~\ref{lem:QuadraticTwistHeightGap}, this implies $6\log p-6\log 2\le C\log 2+B$ for every prime $p\ge 5$. The primes are unbounded, and the left member tends to infinity with $p$. This is impossible. The tame prime $p$ enters the family of Definition~\ref{def:QuadraticTwistDescentFamily} through the radicand, and it is the term $R_{\Q}(q,\ell)$ of Definition~\ref{def:TameKummerSupport} that carries it, so that $A_{\mathrm{desc}}$ is positive.
\end{proof}

\begin{remark}[The source of the obstruction]\label{rem:SourceOfUniformDescentObstruction}
The obstruction does not lie in wild ramification at the prime $\ell$, since in the family above one has $\ell=2$ throughout. The growing term is the tame prime $p$ that enters through $q=p$, and the quadratic twist makes this prime visible in the descent from $L_{p}$ back to $\Q$. Thus a bound logarithmic in $\ell$ alone cannot measure the descent cost, and the tame support of Definition~\ref{def:TameKummerSupport} is the quantity that does.
\end{remark}

\subsubsection{Local scaling under extension}\label{subsubsec:LocalScalingUnderExtension}

The model part of the descent cost is governed by the scaling needed to pass from the base-changed $K$-minimal equation to an $L$-minimal equation. The next local bound is elementary away from residue characteristics $2$ and $3$.

\begin{definition}[Local scaling exponent]\label{def:LocalScalingExponent}
Let $F$ be a non-archimedean local field, let $F'/F$ be a finite extension with ramification index $e$, and let $E/F$ be given by a minimal integral Weierstrass equation $\mathcal{W}$. After base change to $F'$, let $\mathcal{W}'_{\min}$ be a minimal Weierstrass equation for $E_{F'}$. The \emph{local scaling exponent} $r(E,F'/F)$ is the non-negative integer determined by
\[
  v_{F'}(\Delta_{\mathcal{W}'_{\min}})
  =
  e\,v_{F}(\Delta_{\mathcal{W}})
  -
  12\,r(E,F'/F).
\]
Equivalently, $r(E,F'/F)$ is the valuation of the dilation factor that is removed from the base-changed equation during minimization over $F'$.
\end{definition}

\begin{lemma}[Prime-to-six local scaling bound]\label{lem:PrimeToSixLocalScalingBound}
Let $F$ be a finite extension of $\Q_{p}$ with $p\ge 5$. Let $F'/F$ be a finite extension with ramification index $e$. For every elliptic curve $E/F$ given by a minimal integral Weierstrass equation, the local scaling exponent of Definition~\ref{def:LocalScalingExponent} satisfies
\[
  0\le r(E,F'/F)\le e-1.
\]
\end{lemma}

\begin{proof}
Let $c_{4}$, $c_{6}$, and $\Delta$ be the invariants of a minimal integral Weierstrass equation for $E/F$. In residue characteristic at least $5$, the minimality criterion for Weierstrass equations gives
\[
  \min
  \left\{
    \frac{v_{F}(c_{4})}{4},
    \frac{v_{F}(c_{6})}{6},
    \frac{v_{F}(\Delta)}{12}
  \right\}
  <1
\]
with the zero invariants omitted from the minimum. The discriminant is non-zero, so the displayed minimum is finite.

Suppose that a scaling by a uniformizer of $F'$ to exponent $r$ appears in the minimization over $F'$. Since the scaled invariants must remain integral, one has $4r\le e\,v_{F}(c_{4})$ whenever $c_{4}\ne 0$, one has $6r\le e\,v_{F}(c_{6})$ whenever $c_{6}\ne 0$, and one has $12r\le e\,v_{F}(\Delta)$. Thus
\[
  r
  \le
  e
  \min
  \left\{
    \frac{v_{F}(c_{4})}{4},
    \frac{v_{F}(c_{6})}{6},
    \frac{v_{F}(\Delta)}{12}
  \right\}
  <e.
\]
Since $r$ is an integer, it follows that $r\le e-1$.
\end{proof}

The primes above $2$ and $3$ require the full local minimization procedure. Since there are only finitely many such places for a fixed number field $K$, the remaining input can be isolated as a finite local certificate.

\begin{hypothesis}[Small-prime local scaling certificate]\label{hyp:SmallPrimeLocalScalingCertificate}
Let $K$ be a fixed number field. For every finite place $v$ of $K$ above $2$ or $3$, there exists an integer $\eta_{v}\ge 0$ with the following property. For every finite extension $L_{w}/K_{v}$ with ramification index $e_{w/v}$, and for every elliptic curve over $K_{v}$ given by a minimal integral Weierstrass equation, the local scaling exponent satisfies
\[
  r(E,L_{w}/K_{v})
  \le
  e_{w/v}-1+\eta_{v}.
\]
\end{hypothesis}

\begin{remark}[Status of the small-prime certificate]\label{rem:SmallPrimeCertificateStatus}
Hypothesis~\ref{hyp:SmallPrimeLocalScalingCertificate} is local and finite once $K$ is fixed. In a Universe-P presentation, it is supplied as a finite list of certificates for the places over $2$ and $3$. The prime-to-six places are covered by Lemma~\ref{lem:PrimeToSixLocalScalingBound}.
\end{remark}

\subsubsection{A support-controlled descent bound}\label{subsubsec:SupportControlledDescentTheorem}

\begin{proposition}[Support-controlled descent bound]\label{prop:SupportControlledDescentBound}
Fix a number field $K$, and assume Hypothesis~\ref{hyp:SmallPrimeLocalScalingCertificate}. Let $L=K(\sqrt[\ell]{q})$ with $\ell\ge 2$ and $q\in K^{\times}$. Then for every elliptic package $\mathcal{E}$ over $K$, the descent completion contains a height-admissible descent datum
\[
  \delta_{\mathcal{E},L/K}
  \colon
  (L,\mathrm{NF}_{L}(\mathcal{E}_{L}))
  \longrightarrow
  (K,\mathrm{NF}_{K}(\mathcal{E}))
\]
satisfying
\[
  \mu_{\mathrm{tot}}(\delta_{\mathcal{E},L/K})
  \le
  (12+C_{\mathrm{kum}})\log \ell
  +
  12\,R_{K}(q,\ell)
  +
  B_{K,6}.
\]
Here
\[
  B_{K,6}
  \coloneqq
  \frac{12}{[K:\Q]}
  \sum_{v\mid 6}
  \eta_{v}\log \mathrm{N}v
\]
with $\eta_{v}=0$ for places $v\nmid 6$.
\end{proposition}

\begin{proof}
Start with the normalized $K$-model $\mathrm{NF}_{K}(\mathcal{E})$, and base change it to $L$. At each place $w$ of $L$ above a place $v$ of $K$, let $r_{w}$ be the local scaling exponent needed to pass from this base-changed model to the $L_{w}$-minimal model. The model part of the reverse descent cost is
\[
  \mu_{\mathrm{mod}}(\delta_{\mathcal{E},L/K})
  =
  \frac{12}{[L:\Q]}
  \sum_{w}
  r_{w}\log \mathrm{N}w.
\]
Unramified extensions preserve minimality of Weierstrass equations, so places unramified in $L/K$ have $r_{w}=0$. The polynomial $T^{\ell}-q$ has discriminant divisible only by primes above $\ell$ and by primes in the support of $q$. Hence $L/K$ is unramified outside the places above $\ell$ and the tame support of Definition~\ref{def:TameKummerSupport}.

For $v\nmid 6$, Lemma~\ref{lem:PrimeToSixLocalScalingBound} gives $r_{w}\le e_{w/v}-1$. For $v\mid 6$, Hypothesis~\ref{hyp:SmallPrimeLocalScalingCertificate} gives $r_{w}\le e_{w/v}-1+\eta_{v}$. Consequently, for each ramified place $v$ of $K$,
\[
  \frac{12}{[L:\Q]}
  \sum_{w\mid v}
  r_{w}\log \mathrm{N}w
  \le
  \frac{12}{[L:\Q]}
  \sum_{w\mid v}
  e_{w/v}\log \mathrm{N}w
  +
  \frac{12\eta_{v}}{[L:\Q]}
  \sum_{w\mid v}
  \log \mathrm{N}w.
\]
The first term on the right is bounded by $\tfrac{12}{[K:\Q]}\log \mathrm{N}v$, since $\sum_{w\mid v}e_{w/v}\log\mathrm{N}w=[L:K]\log\mathrm{N}v$. The second term contributes only at the finite set of places $v\mid 6$, and its total contribution is bounded by $B_{K,6}$.

Summing over the ramified places away from $\ell$ gives the term $12R_{K}(q,\ell)$. Summing over the places $v\mid \ell$ gives
\[
  \frac{12}{[K:\Q]}
  \sum_{v\mid \ell}\log \mathrm{N}v
  \le
  12\log \ell.
\]
Therefore $\mu_{\mathrm{mod}}(\delta_{\mathcal{E},L/K})\le 12R_{K}(q,\ell)+12\log \ell+B_{K,6}$. The Kummer detachment term satisfies $\mu_{\mathrm{kum}}(\delta_{\mathcal{E},L/K})\le C_{\mathrm{kum}}\log \ell$ by Definition~\ref{def:KummerDefect}. Adding the two estimates gives the displayed bound for $\mu_{\mathrm{tot}}$.
\end{proof}

\begin{remark}[The hypothesis is proved away from the small primes]\label{rem:SupportHypothesisProved}
Proposition~\ref{prop:SupportControlledDescentBound} proves Hypothesis~\ref{hyp:ControlledDescentWithSupport} under the finite local certificate of Hypothesis~\ref{hyp:SmallPrimeLocalScalingCertificate}, with $C_{\mathrm{desc}}=12+C_{\mathrm{kum}}$, $A_{\mathrm{desc}}=12$, and $B_{\mathrm{desc},K}=B_{K,6}$. The bound is therefore explicit, and the only local input that remains beyond the prime-to-six estimate is concentrated at the finitely many places above $2$ and $3$. Together with Theorem~\ref{thm:NoUniformControlledDescent}, this leaves the support-aware form as the correct shape of the descent estimate, with the coefficient $A_{\mathrm{desc}}=12$ both necessary and sufficient.
\end{remark}

\begin{corollary}[Recovery of the purely logarithmic form under bounded tame support]
\label{cor:RecoveryControlledDescentLogCost}
Assume Hypothesis~\ref{hyp:ControlledDescentWithSupport}. Let $\mathcal{L}$ be a family of Kummer extensions $L=K(\sqrt[\ell]{q})$ for which there is a constant $B_{\mathrm{supp},K}$ with $R_{K}(q,\ell)\le B_{\mathrm{supp},K}$ for every member of the family. Then for every member of $\mathcal{L}$ the descent bound assumes the purely logarithmic form
\[
  \mu_{\mathrm{tot}}(\delta_{\mathcal{E},L/K})
  \le
  C_{\mathrm{desc}}\log \ell
  +
  B_{\mathrm{desc},K}^{\prime}
\]
with $B_{\mathrm{desc},K}^{\prime}=A_{\mathrm{desc}}B_{\mathrm{supp},K}+B_{\mathrm{desc},K}$.
\end{corollary}

\begin{proof}
Substitute the bound $R_{K}(q,\ell)\le B_{\mathrm{supp},K}$ into the inequality of Hypothesis~\ref{hyp:ControlledDescentWithSupport}. The support term is thereby absorbed into the constant, and the bound assumes the displayed purely logarithmic form.
\end{proof}

\begin{remark}[Two regimes of the descent bound]\label{rem:DescentRegimes}
The estimate of Hypothesis~\ref{hyp:ControlledDescentWithSupport} is read in two regimes. For a family of bounded tame support, as in Corollary~\ref{cor:RecoveryControlledDescentLogCost}, the support term collapses into a constant and the bound is logarithmic in the Kummer degree. For a family of unbounded support, the term $R_{K}(q,\ell)$ is carried through the later inequalities, where it measures the cost of the primes that the radicand introduces.
\end{remark}

\subsubsection{Arithmetic formalization}\label{subsubsec:ControlledDescentACA}

\begin{proposition}[Formalization in \texorpdfstring{\(\mathsf{ACA}_{0}\)}{ACA0}]
\label{prop:ControlledDescentACAFormalization}
The constructions and proofs of this subsection are formalizable in $\mathsf{ACA}_{0}$, once the finite local certificates of Hypothesis~\ref{hyp:SmallPrimeLocalScalingCertificate} are supplied as part of the certified field code.
\end{proposition}

\begin{proof}
The quadratic-twist family of Definition~\ref{def:QuadraticTwistDescentFamily} is given by the primitive recursive map
\[
  p\longmapsto [0,0,0,-p^{2},0]
\]
on Weierstrass coefficients. The discriminant identity $\Delta(E_{p})=64p^{6}$ is an identity among integers. The height estimates of Lemma~\ref{lem:QuadraticTwistHeightGap} are therefore comparisons of logarithmic integer data in the sense of Definition~\ref{def:LogIntDatum}. The unboundedness of the primes, used in Theorem~\ref{thm:NoUniformControlledDescent}, is provable in elementary arithmetic and hence in $\mathsf{ACA}_{0}$.

The tame support $R_{K}(q,\ell)$ of Definition~\ref{def:TameKummerSupport} is computed from the valuations of the finite element $q$ in the certified field code of Definition~\ref{def:CertifiedFieldCode}. The assertion that $K(\sqrt[\ell]{q})/K$ is unramified outside the primes above $\ell$ and the support of $q$ follows from the discriminant of the polynomial $T^{\ell}-q$, which is a finite algebraic computation. The local scaling exponents of Definition~\ref{def:LocalScalingExponent} are integer valuations of discriminants before and after minimization. Lemma~\ref{lem:PrimeToSixLocalScalingBound} is a finite valuation argument using $c_{4}$, $c_{6}$, and $\Delta$. The remaining places above $2$ and $3$ are handled by the finite list of certificates in Hypothesis~\ref{hyp:SmallPrimeLocalScalingCertificate}.

All sums in Proposition~\ref{prop:SupportControlledDescentBound} range over finite lists of prime ideals appearing in the certified field code and in the factorization of $q\ell$. The real quantities are represented by logarithmic integer data or by fast Cauchy names, as in Definition~\ref{def:UniverseP}. No comprehension beyond arithmetical comprehension is invoked.
\end{proof}

%==================================================================
\section{An intensional Szpiro inequality}\label{sec:IntensionalSzpiro}
%==================================================================

We now arrive at the arithmetic statement that the framework was built to carry. The Szpiro inequality, in its classical reading over a fixed number field, compares the minimal discriminant of an elliptic curve with its conductor. In Universe~A the curve is replaced by a presentation, the discriminant by the discriminant height of Definition~\ref{def:DiscriminantHeight}, and the comparison gains a third term that keeps the cost of the chosen presentation. The dilation tower of Section~\ref{subsec:DilationTower} showed that this height moves along an isomorphism, by an amount the model defect measures exactly, so that a scalar comparable to it cannot be invariant under all isomorphisms of packages. The log-volume of the inequality is therefore read as a quantity that transforms under transport. We make this reading precise, prove that the discriminant height itself serves as the log-volume, and then reduce the remaining arithmetic input to a single reference profile. The inequality of Proposition~\ref{prop:NormalizedToIntensionalSzpiro} is the intensional shape of Szpiro, and the classical statement is recovered as its costless, normalized specialization.

\subsection{The transport Szpiro log-volume}\label{subsec:TransportLogVolume}

The dilation tower indicates the correct shape for the log-volume. It should be invariant under coherent changes of the transport, namely under $2$-isomorphism of $1$-morphisms, while it transforms along positive-cost $1$-morphisms by an inequality.

\begin{definition}[Transport Szpiro log-volume]\label{def:TransportSzpiroLogVolume}
Let $K$ be a number field, and let $\mathfrak{EllPkg}_K$ be a quantitative $(2,1)$-category of elliptic packages as in Definition~\ref{def:EllPkgTwoOne}. A \emph{transport Szpiro log-volume} over $K$ is a numerical assignment
\[
  V \colon \operatorname{Ob}(\mathfrak{EllPkg}_K)\longrightarrow \R_{\ge 0}
\]
for which there exist constants $A,B>0$ and $C_0\ge 0$, depending only upon $K$, such that the following conditions hold:
\begin{enumerate}
\item[\textup{(i)}] The transport law factors through the homotopy $1$-category. Equivalently, if $\alpha \colon f \Rightarrow g$ is an invertible $2$-morphism between parallel $1$-morphisms, then the inequalities attached to $f$ and to $g$ have the same numerical cost.
\item[\textup{(ii)}] For every elliptic package $\mathcal{E}$ one has
\[
  A^{-1}h_\Delta(\mathcal{E})-B
  \le
  V(\mathcal{E})
  \le
  Ah_\Delta(\mathcal{E})+B.
\]
\item[\textup{(iii)}] For every $1$-morphism $f\colon\mathcal{E}_1\to\mathcal{E}_2$ in $\mathfrak{EllPkg}_K$ one has
\[
  V(\mathcal{E}_2)
  \le
  V(\mathcal{E}_1)+C_0\,\mu(f).
\]
\end{enumerate}
\end{definition}

\begin{theorem}[The discriminant height is a transport log-volume]
\label{thm:DiscriminantHeightAsTransportLogVolume}
For every number field $K$, the assignment
\[
  \mathrm{Vol}^{A}_{\mathrm{Szp}}(\mathcal{E})
  \coloneqq
  h_\Delta(\mathcal{E})
\]
is a transport Szpiro log-volume in the sense of Definition~\ref{def:TransportSzpiroLogVolume}. One may take $A=1$, $B=0$, and $C_0=1$.
\end{theorem}

\begin{proof}
The comparability condition of Definition~\ref{def:TransportSzpiroLogVolume}\textup{(ii)} is immediate from the identification $\mathrm{Vol}^{A}_{\mathrm{Szp}}=h_\Delta$, with $A=1$ and $B=0$. The transport law of Definition~\ref{def:TransportSzpiroLogVolume}\textup{(iii)} is exactly Lemma~\ref{lem:TransportDiscriminant}, with $C_0=1$. Finally, the cost $\mu$ is invariant under $2$-isomorphism by Definition~\ref{def:EllPkgTwoOne}\textup{(iv)}, so that the transport law depends only upon the image of the transporting morphism in $\Ho(\mathfrak{EllPkg}_K)$, as required in Definition~\ref{def:TransportSzpiroLogVolume}\textup{(i)}.
\end{proof}

\begin{corollary}[Zero-defect isomorphisms recover invariance]\label{cor:ZeroDefectSubcategoryLogVolume}
Let $\mathfrak{EllPkg}_K^{0}$ be the wide subgroupoid of $\mathfrak{EllPkg}_K$ whose isomorphisms $f$ satisfy $\mu(f)=\mu(f^{-1})=0$. Then $h_\Delta$ is invariant on $\mathfrak{EllPkg}_K^{0}$.
\end{corollary}

\begin{proof}
If $f\colon\mathcal{E}\to\mathcal{E}'$ lies in $\mathfrak{EllPkg}_K^{0}$, then Corollary~\ref{cor:TwoSidedControlTransport} gives
\[
  \bigl|h_\Delta(\mathcal{E}')-h_\Delta(\mathcal{E})\bigr|
  \le
  \mu(f)+\mu(f^{-1})
  =
  0
\]
so that $h_\Delta(\mathcal{E}')=h_\Delta(\mathcal{E})$.
\end{proof}

\begin{remark}[The log-volume is the discriminant height]\label{rem:CorrectedLogVolumeDatum}
Theorem~\ref{thm:DiscriminantHeightAsTransportLogVolume} settles the structural part of the inequality without an additional hypothesis. The log-volume need not be assumed as an extra object. It may be taken to be the discriminant height itself, invariant under $2$-isomorphism of transports by Proposition~\ref{prop:CoherenceStability} and transforming along positive-cost $1$-morphisms by Lemma~\ref{lem:TransportDiscriminant}. The reading invariant under all package isomorphisms is the one the dilation tower of Proposition~\ref{prop:NoInvariantLogVolume} excludes, and the transport reading is the one compatible with Definition~\ref{def:EllPkgTwoOne}.
\end{remark}

\subsection{The reduced normalized profile}\label{subsec:ReducedProfile}

Once Theorem~\ref{thm:DiscriminantHeightAsTransportLogVolume} is in place, the remaining arithmetic input lies in the choice of a reference profile. It is worth saying in plain terms what the symbols of the profile stand for. The reference profile endofunctor $\mathsf{Q}_{K,\varepsilon}$ produces, for each package, an auxiliary presentation whose discriminant height is held below the conductor directly. The family $\tau_{\mathcal{E},\varepsilon}$ is the transport from that reference presentation to the normalized one, and its cost is the quantity the $\varepsilon$-term absorbs. The two constants are the bounded error terms of the classical inequality, each depending only upon the field $K$. The constant $B_{q,K}$ is the slack in the comparison of the reference height with six times the conductor complexity, and $B_{\varepsilon,K}$ is the constant summand of the transport cost that remains once the part proportional to the conductor has been set against $\varepsilon$. No calibration constant between a log-volume and the discriminant height is needed, since Theorem~\ref{thm:DiscriminantHeightAsTransportLogVolume} identifies the two.

\begin{hypothesis}[Reduced normalized profile]\label{hyp:ReducedNormalizedProfile}
Let $K$ be a number field and fix a normalization procedure $\mathrm{NF}_K$ in the sense of Definition~\ref{def:AccumulatedDefect}. For every real number $\varepsilon>0$, there exist a reference profile endofunctor
\[
  \mathsf{Q}_{K,\varepsilon}
  \colon
  \mathbf{EllPkg}_K\longrightarrow\mathbf{EllPkg}_K
\]
and a family of morphisms
\[
  \tau_{\mathcal{E},\varepsilon}
  \colon
  \mathsf{Q}_{K,\varepsilon}(\mathcal{E})
  \longrightarrow
  \mathrm{NF}_K(\mathcal{E})
\]
together with constants $B_{q,K}\ge 0$ and $B_{\varepsilon,K}\ge 0$ such that for every elliptic package $\mathcal{E}$ over $K$ one has the reference-side bound
\[
  h_\Delta\bigl(\mathsf{Q}_{K,\varepsilon}(\mathcal{E})\bigr)
  \le
  6\,n(\mathcal{E})+B_{q,K}
\]
and the transport bound
\[
  \mu\bigl(\tau_{\mathcal{E},\varepsilon}\bigr)
  \le
  \varepsilon\,n(\mathcal{E})+B_{\varepsilon,K}.
\]
When functorial comparison is required, the family $\tau_{\mathcal{E},\varepsilon}$ is assumed to extend to a pseudonatural transformation $\tau_{\varepsilon}\colon\mathsf{Q}_{K,\varepsilon}\Rightarrow\mathrm{NF}_K$ in the sense of Definition~\ref{def:PseudoNatural}, so that for every $1$-morphism $f\colon\mathcal{E}\to\mathcal{E}'$ there is an invertible $2$-morphism
\[
  (\tau_\varepsilon)_f \colon
  \tau_{\mathcal{E}',\varepsilon}\circ \mathsf{Q}_{K,\varepsilon}(f)
  \Rightarrow
  \mathrm{NF}_K(f)\circ \tau_{\mathcal{E},\varepsilon}.
\]
\end{hypothesis}

\begin{remark}[What remains after the reduction]\label{rem:WhatRemainsAfterReduction}
Theorem~\ref{thm:DiscriminantHeightAsTransportLogVolume} removes the need to posit a separate log-volume satisfying the transport law. The remaining Diophantine content is concentrated in Hypothesis~\ref{hyp:ReducedNormalizedProfile}. It asks for a reference package whose discriminant height is controlled by the conductor complexity, together with a transport from that reference package to the normalized package whose cost is absorbed by the $\varepsilon$-term.
\end{remark}

\subsection{From the profile to the inequality}\label{subsec:ProfileToInequality}

\begin{proposition}[The intensional Szpiro inequality]\label{prop:NormalizedToIntensionalSzpiro}
Assume Hypothesis~\ref{hyp:ReducedNormalizedProfile}. Then for every real number $\varepsilon>0$ there exists a constant $C_{\varepsilon,K}\in\R$ such that for every elliptic package $\mathcal{E}$ over $K$ one has
\begin{equation}\label{eq:IntensionalSzpiroBasic}
  h_\Delta(\mathcal{E})
  \le
  (6+\varepsilon)\,n(\mathcal{E})
  +
  \defect(\mathcal{E})
  +
  C_{\varepsilon,K}.
\end{equation}
One may take $C_{\varepsilon,K}=B_{q,K}+B_{\varepsilon,K}$.
\end{proposition}

\begin{proof}
Fix $\varepsilon>0$ and let $\mathcal{E}$ be an elliptic package over $K$. Applying the transport inequality of Lemma~\ref{lem:TransportDiscriminant} to the distinguished normalization morphism $\nu_{\mathcal{E}}\colon\mathrm{NF}_K(\mathcal{E})\to\mathcal{E}$ of Definition~\ref{def:AccumulatedDefect} gives
\[
  h_\Delta(\mathcal{E})
  \le
  h_\Delta\bigl(\mathrm{NF}_K(\mathcal{E})\bigr)
  +
  \defect(\mathcal{E})
\]
where $\mu(\nu_{\mathcal{E}})=\defect(\mathcal{E})$ by Definition~\ref{def:AccumulatedDefect}. Applying Lemma~\ref{lem:TransportDiscriminant} again to $\tau_{\mathcal{E},\varepsilon}\colon\mathsf{Q}_{K,\varepsilon}(\mathcal{E})\to\mathrm{NF}_K(\mathcal{E})$ gives
\[
  h_\Delta\bigl(\mathrm{NF}_K(\mathcal{E})\bigr)
  \le
  h_\Delta\bigl(\mathsf{Q}_{K,\varepsilon}(\mathcal{E})\bigr)
  +
  \mu\bigl(\tau_{\mathcal{E},\varepsilon}\bigr).
\]
The two bounds of Hypothesis~\ref{hyp:ReducedNormalizedProfile} then yield
\[
  h_\Delta(\mathcal{E})
  \le
  6\,n(\mathcal{E})+B_{q,K}
  +
  \varepsilon\,n(\mathcal{E})+B_{\varepsilon,K}
  +
  \defect(\mathcal{E})
\]
which is the stated inequality with $C_{\varepsilon,K}=B_{q,K}+B_{\varepsilon,K}$.
\end{proof}

\begin{lemma}[Objectwise bounds are insensitive to pseudonaturality]\label{lem:ObjectwiseInsensitive}
The inequality \eqref{eq:IntensionalSzpiroBasic} is derived in Proposition~\ref{prop:NormalizedToIntensionalSzpiro} from the two objectwise bounds of Hypothesis~\ref{hyp:ReducedNormalizedProfile} together with two applications of the transport inequality of Lemma~\ref{lem:TransportDiscriminant}. In particular, it does not use the pseudonatural comparison that the hypothesis supplies when functorial transport is required.
\end{lemma}

\begin{proof}
This is immediate from the proof of Proposition~\ref{prop:NormalizedToIntensionalSzpiro}, which invokes only the reference-side bound, the transport bound, and the transport inequality for $\nu_{\mathcal{E}}$ and $\tau_{\mathcal{E},\varepsilon}$.
\end{proof}

\begin{proposition}[Elliptic packages as a quantitative category of Szpiro form]\label{prop:EllPkgQuantitativeSzpiro}
Let $K$ be a number field. Consider the category $\mathbf{EllPkg}_K$ of elliptic packages over $K$, equipped with the model defect $\mu$ of Definition~\ref{def:ModelDefect} and the discriminant height $h_\Delta$ of Definition~\ref{def:DiscriminantHeight}. Then the following hold:
\begin{enumerate}
\item[(i)] The pair $(\mathbf{EllPkg}_K,\mu)$ forms a quantitative category of isomorphisms in the sense of Definition~\ref{def:QuantitativeCategory}.
\item[(ii)] The discriminant height satisfies the transport inequality of Lemma~\ref{lem:TransportDiscriminant}, so that for any morphism $f \colon \mathcal{E} \to \mathcal{E}'$ one has $h_\Delta(\mathcal{E}') \le h_\Delta(\mathcal{E}) + \mu(f)$.
\item[(iii)] Under Hypothesis~\ref{hyp:ReducedNormalizedProfile}, the basic intensional inequality \eqref{eq:IntensionalSzpiroBasic} holds, in the form
  \[
    h_\Delta(\mathcal{E})
    \le
    (6+\varepsilon)\,n(\mathcal{E})
    +
    \defect(\mathcal{E})
    +
    C_{\varepsilon,K}.
  \]
It keeps the place where the normalization defect enters, through the transport $\nu_{\mathcal{E}}$ of Definition~\ref{def:AccumulatedDefect}.
\end{enumerate}
\end{proposition}

\begin{proof}
Part \textup{(i)} is contained in Definition~\ref{def:ModelDefect}, part \textup{(ii)} is Lemma~\ref{lem:TransportDiscriminant}, and part \textup{(iii)} is Proposition~\ref{prop:NormalizedToIntensionalSzpiro}.
\end{proof}

\begin{remark}[Comparison with the classical Szpiro inequality]\label{rem:SzpiroComparison}
When the package is normalized by the chosen procedure, one has $\defect(\mathcal{E}) = 0$, and the inequality \eqref{eq:IntensionalSzpiroBasic} assumes the classical shape up to the constant $C_{\varepsilon,K}$. The defect term is precisely the quantity that the extensional reading discards. By keeping it, the intensional inequality holds two quantities together -- the comparison of discriminant and conductor on the one hand, and on the other the price of the presentation through which that comparison was read.
\end{remark}

\subsection{Formalization in \texorpdfstring{\(\mathsf{ACA}_0\)}{ACA0}}\label{subsec:ACA0LogVolume}

The arguments of this section stay within the arithmetic universe of Section~\ref{sec:UniverseP}. They use finite Weierstrass coefficients, finite coordinate changes, integer valuations, and inequalities between logarithmic integer data.

\begin{proposition}[Arithmetic formalization of the log-volume reduction]
\label{prop:ACA0LogVolumeFormalization}
The statements and proofs of Lemma~\ref{lem:DilationTowerGrowth}, Proposition~\ref{prop:NoInvariantLogVolume}, Theorem~\ref{thm:DiscriminantHeightAsTransportLogVolume}, and Proposition~\ref{prop:NormalizedToIntensionalSzpiro} are formalizable in $\mathsf{ACA}_0$ after replacing real-valued quantities by the coding conventions of Definition~\ref{def:LogIntDatum} and Definition~\ref{def:EllPkgCode}.
\end{proposition}

\begin{proof}
The dilation tower of Definition~\ref{def:DilationTowerAtBadPrime} is given by the primitive recursive map
\[
  m
  \longmapsto
  [0,0,37^{3m},-37^{4m},0]
\]
on Weierstrass coefficients. The discriminant computation is an identity among integers, $\Delta_{\mathcal{W}_m}=37^{12m+1}$, so that the height calculation in Lemma~\ref{lem:DilationTowerGrowth} becomes the comparison of the logarithmic integer datum $37^{12m+1}$ with $1$, and the defect calculation becomes the logarithmic integer datum $37^{12m}$.

The contradiction in Proposition~\ref{prop:NoInvariantLogVolume} uses only the archimedean growth of the sequence $(12m+1)\log 37$. For real numbers presented by fast Cauchy names, $\mathsf{ACA}_0$ proves the needed bounded-search statement. Namely, for fixed coded reals $A>0$, $B$, and $V(\mathcal{E}_0)$, there is an integer $m$ such that $A^{-1}(12m+1)\log 37-B>V(\mathcal{E}_0)$.

Theorem~\ref{thm:DiscriminantHeightAsTransportLogVolume} is a direct transcription of Lemma~\ref{lem:TransportDiscriminant}, which in code form is an inequality between finite sums of local valuation contributions. Proposition~\ref{prop:NormalizedToIntensionalSzpiro} uses only two applications of this transport inequality together with the two bounds supplied by Hypothesis~\ref{hyp:ReducedNormalizedProfile}. All these are finite operations on the package codes of Definition~\ref{def:EllPkgCode}, with real values handled as fast Cauchy names in the sense of Definition~\ref{def:UniverseP}.
\end{proof}

\subsection{Partial reductions for the reduced normalized profile}
\label{subsec:PartialReducedNormalizedProfile}

The reduced normalized profile of Hypothesis~\ref{hyp:ReducedNormalizedProfile} carries two demands, one arithmetical and one intensional. The first asks for a presentation whose discriminant height lies close to six times the conductor complexity. The second asks that the transport from that presentation to the chosen normalization have small cost. This subsection does not prove the full profile. It discharges the intensional demand within the framework, reduces the arithmetical demand to a single linear envelope, and then localizes that envelope to one explicit quantity, the height of the denominator of the $j$-invariant. Everything apart from a bound on that one quantity is proved unconditionally and within $\mathsf{ACA}_{0}$.

\begin{definition}[Normalized height and remainder]
\label{def:NormalizedHeightRemainder}
Fix a number field $K$ and a normalization procedure $\mathrm{NF}_{K}$ in the sense of Definition~\ref{def:AccumulatedDefect}. The \emph{normalized discriminant height} of an elliptic package $\mathcal{E}$ is
\[
  H^{\mathrm{nf}}_{K}(\mathcal{E})
  \coloneqq
  h_{\Delta}\bigl(\mathrm{NF}_{K}(\mathcal{E})\bigr).
\]
For a real number $a>0$, its \emph{$a$-remainder} is
\[
  \mathcal{R}_{K,a}(\mathcal{E})
  \coloneqq
  H^{\mathrm{nf}}_{K}(\mathcal{E})-a\,n(\mathcal{E}).
\]
For a finite coded list $\mathcal{F}=(\mathcal{E}_{0},\ldots,\mathcal{E}_{m-1})$ of package codes in the sense of Definition~\ref{def:EllPkgCode}, put
\[
  B_{\mathcal{F},a,K}
  \coloneqq
  \max_{0\le i<m}
  \max\bigl\{0,\mathcal{R}_{K,a}(\mathcal{E}_{i})\bigr\}.
\]
\end{definition}

\subsubsection{The intensional demand is discharged}\label{subsubsec:IntensionalDemandFree}

\begin{proposition}[The profile forces a normalized Szpiro bound]
\label{prop:ReducedProfileForcesNormalizedSzpiro}
Assume Hypothesis~\ref{hyp:ReducedNormalizedProfile}. Then for every real number $\varepsilon>0$ and every elliptic package $\mathcal{E}$ over $K$, one has
\[
  H^{\mathrm{nf}}_{K}(\mathcal{E})
  \le
  (6+\varepsilon)\,n(\mathcal{E})+B_{q,K}+B_{\varepsilon,K}.
\]
\end{proposition}

\begin{proof}
Apply the transport inequality of Lemma~\ref{lem:TransportDiscriminant} to the morphism $\tau_{\mathcal{E},\varepsilon}\colon\mathsf{Q}_{K,\varepsilon}(\mathcal{E})\to\mathrm{NF}_{K}(\mathcal{E})$ of Hypothesis~\ref{hyp:ReducedNormalizedProfile}. This gives
\[
  h_{\Delta}\bigl(\mathrm{NF}_{K}(\mathcal{E})\bigr)
  \le
  h_{\Delta}\bigl(\mathsf{Q}_{K,\varepsilon}(\mathcal{E})\bigr)
  +
  \mu\bigl(\tau_{\mathcal{E},\varepsilon}\bigr).
\]
The reference-side bound and the transport bound of Hypothesis~\ref{hyp:ReducedNormalizedProfile} give the displayed estimate.
\end{proof}

\begin{lemma}[The identity profile]
\label{lem:IdentityProfileReduction}
Suppose that, for some real numbers $a>0$ and $B_{a,K}\ge 0$, one has $H^{\mathrm{nf}}_{K}(\mathcal{E})\le a\,n(\mathcal{E})+B_{a,K}$ for every elliptic package $\mathcal{E}$ over $K$. Then the choice
\[
  \mathsf{Q}_{K,a}(\mathcal{E})\coloneqq \mathrm{NF}_{K}(\mathcal{E})
  \qquad\text{and}\qquad
  \tau_{\mathcal{E},a}\coloneqq \mathrm{id}_{\mathrm{NF}_{K}(\mathcal{E})}
\]
satisfies $h_{\Delta}(\mathsf{Q}_{K,a}(\mathcal{E}))\le a\,n(\mathcal{E})+B_{a,K}$ and $\mu(\tau_{\mathcal{E},a})=0$. Read $\mathrm{NF}_{K}$ as an endofunctor of the chosen $(2,1)$-category. Then these identity morphisms form the identity pseudonatural comparison from $\mathrm{NF}_{K}$ to itself, in the sense of Definition~\ref{def:PseudoNatural}.
\end{lemma}

\begin{proof}
The first inequality is the assumed normalized bound. The equality $\mu(\mathrm{id})=0$ follows from Definition~\ref{def:ModelDefect}, since the model defect of an identity morphism vanishes. The pseudonatural comparison is the unit comparison for the endofunctor $\mathrm{NF}_{K}$.
\end{proof}

\begin{remark}[Isolation of the arithmetical core]
\label{rem:ReducedProfileArithmeticCore}
Proposition~\ref{prop:ReducedProfileForcesNormalizedSzpiro} and Lemma~\ref{lem:IdentityProfileReduction} pass in opposite directions between the reduced profile and a normalized linear envelope, so that the arithmetical content of Hypothesis~\ref{hyp:ReducedNormalizedProfile} is exactly such an envelope on $H^{\mathrm{nf}}_{K}$. The missing input is therefore not the coherence datum of Definition~\ref{def:PseudoNatural}. By Lemma~\ref{lem:ObjectwiseInsensitive} that datum is silent for the objectwise inequality, and the identity profile supplies it at no cost. The missing input is the normalized remainder $\mathcal{R}_{K,6+\varepsilon}(\mathcal{E})$, and it is here that the abc-level strength would enter.
\end{remark}

\subsubsection{From a normalized envelope to the inequality}\label{subsubsec:EnvelopeToInequality}

\begin{proposition}[The defect-normalized bound]
\label{prop:DefectNormalizedBound}
For every elliptic package $\mathcal{E}$ over $K$, one has
\[
  h_{\Delta}(\mathcal{E})
  \le
  H^{\mathrm{nf}}_{K}(\mathcal{E})+
  \defect(\mathcal{E}).
\]
Consequently, for every finite coded list $\mathcal{F}=(\mathcal{E}_{0},\ldots,\mathcal{E}_{m-1})$ and every real number $a>0$, one has
\[
  h_{\Delta}(\mathcal{E}_{i})
  \le
  a\,n(\mathcal{E}_{i})+
  \defect(\mathcal{E}_{i})+
  B_{\mathcal{F},a,K}
\]
for each $0\le i<m$.
\end{proposition}

\begin{proof}
Apply Lemma~\ref{lem:TransportDiscriminant} to the distinguished normalization morphism $\nu_{\mathcal{E}}\colon\mathrm{NF}_{K}(\mathcal{E})\to\mathcal{E}$ of Definition~\ref{def:AccumulatedDefect}. Since $\mu(\nu_{\mathcal{E}})=\defect(\mathcal{E})$, this gives the first inequality. For the finite list, Definition~\ref{def:NormalizedHeightRemainder} gives $H^{\mathrm{nf}}_{K}(\mathcal{E}_{i})\le a\,n(\mathcal{E}_{i})+B_{\mathcal{F},a,K}$ for each index $i$, and substitution into the first inequality gives the claim.
\end{proof}

\begin{hypothesis}[Normalized linear envelope]
\label{hyp:NormalizedLinearEnvelope}
There exist real numbers $a_{K}>0$ and $B_{K}\ge 0$, depending only upon $K$ and the normalization procedure $\mathrm{NF}_{K}$, such that for every elliptic package $\mathcal{E}$ over $K$ one has
\[
  H^{\mathrm{nf}}_{K}(\mathcal{E})
  \le
  a_{K}\,n(\mathcal{E})+B_{K}.
\]
\end{hypothesis}

\begin{proposition}[A global intensional bound from a linear envelope]
\label{prop:LinearEnvelopeGivesIntensionalBound}
Assume Hypothesis~\ref{hyp:NormalizedLinearEnvelope}. Then every elliptic package $\mathcal{E}$ over $K$ satisfies
\[
  h_{\Delta}(\mathcal{E})
  \le
  a_{K}\,n(\mathcal{E})+
  \defect(\mathcal{E})+B_{K}.
\]
If, for every $\varepsilon>0$, the envelope may be taken with $a_{K}=6+\varepsilon$, then one recovers the intensional Szpiro inequality of Proposition~\ref{prop:NormalizedToIntensionalSzpiro} with $C_{\varepsilon,K}=B_{K}$.
\end{proposition}

\begin{proof}
The first assertion is Proposition~\ref{prop:DefectNormalizedBound} followed by Hypothesis~\ref{hyp:NormalizedLinearEnvelope}. The final assertion is the same inequality with $a_{K}=6+\varepsilon$.
\end{proof}

\subsubsection{The potentially good part is unconditional}\label{subsubsec:PotentiallyGoodPart}

The normalized linear envelope is a single inequality, and yet it is not innocent. We separate the part of it that the local theory settles outright from the part that carries the arithmetic weight. The separation runs along the dichotomy between potentially good and potentially multiplicative reduction.

\begin{definition}[Multiplicative excess and the potentially good defect]
\label{def:MultiplicativeExcess}
Let $\mathcal{E}$ be an elliptic package over $K$ with underlying curve $E$ of $j$-invariant $j_{E}$. The \emph{multiplicative excess height} of $\mathcal{E}$ is
\[
  \Theta_{K}(\mathcal{E})
  \coloneqq
  \frac{1}{[K:\Q]}
  \sum_{v}
  \max\bigl\{0,-\mathrm{ord}_{v}(j_{E})\bigr\}\,\log \mathrm{N}v
\]
the sum running over the finite places of $K$, so that $\Theta_{K}(\mathcal{E})$ is the logarithmic height of the denominator of $j_{E}$. For a finite place $v$, the \emph{potentially good defect} of $\mathcal{E}$ at $v$ is
\[
  \delta_{v}(\mathcal{E})
  \coloneqq
  \mathrm{ord}_{v}\bigl(\Delta_{\mathrm{NF}_{K}(\mathcal{E}),v}\bigr)
  -
  \max\bigl\{0,-\mathrm{ord}_{v}(j_{E})\bigr\}
\]
the difference between the minimal discriminant exponent and the denominator exponent of $j_{E}$ at $v$.
\end{definition}

\begin{lemma}[The defect is non-negative and locally bounded]
\label{lem:PotentiallyGoodDefectBound}
For every elliptic package $\mathcal{E}$ over $K$ and every finite place $v$ of $K$, one has $\delta_{v}(\mathcal{E})\ge 0$. If the residue characteristic at $v$ is at least $5$, then
\[
  \delta_{v}(\mathcal{E})\le 5\,\mathfrak{f}_{v}(\mathcal{E})
\]
where $\mathfrak{f}_{v}(\mathcal{E})$ is the conductor exponent of $E$ at $v$. If the residue characteristic at $v$ is $2$ or $3$, then $\delta_{v}(\mathcal{E})\le\kappa_{v}$ for a constant $\kappa_{v}$ depending only upon $v$.
\end{lemma}

\begin{proof}
Write $\Delta$ for the minimal discriminant and $c_{4}$ for the corresponding invariant of the model $\mathrm{NF}_{K}(\mathcal{E})$, so that $\mathrm{ord}_{v}(j_{E})=3\,\mathrm{ord}_{v}(c_{4})-\mathrm{ord}_{v}(\Delta)$. The model is integral, so $\mathrm{ord}_{v}(c_{4})\ge 0$, whence $-\mathrm{ord}_{v}(j_{E})\le\mathrm{ord}_{v}(\Delta)$ and therefore $\max\{0,-\mathrm{ord}_{v}(j_{E})\}\le\mathrm{ord}_{v}(\Delta)$. This gives $\delta_{v}(\mathcal{E})\ge 0$.

For the upper bound, read the N\'eron--Kodaira classification of the special fiber through Tate's algorithm (see Silverman~\cite{SilvermanAEC} and Cremona~\cite{Cremona1997}). At residue characteristic at least $5$ the reduction is one of the following. At good reduction $\mathrm{ord}_{v}(\Delta)=0$ and $\mathfrak{f}_{v}=0$, so $\delta_{v}=0$. At multiplicative reduction of type $I_{n}$ one has $\mathrm{ord}_{v}(\Delta)=n=-\mathrm{ord}_{v}(j_{E})$ and $\mathfrak{f}_{v}=1$, so $\delta_{v}=0$. At additive reduction of potentially multiplicative type $I_{n}^{*}$ one has $\mathrm{ord}_{v}(\Delta)=n+6$, $-\mathrm{ord}_{v}(j_{E})=n$, and $\mathfrak{f}_{v}=2$, so $\delta_{v}=6\le 5\mathfrak{f}_{v}$. At the remaining additive types, which are of potentially good reduction, $j_{E}$ is integral and so $\max\{0,-\mathrm{ord}_{v}(j_{E})\}=0$, while $\mathrm{ord}_{v}(\Delta)$ takes one of the values $2,3,4,6,8,9,10$ and $\mathfrak{f}_{v}=2$. In each case $\delta_{v}=\mathrm{ord}_{v}(\Delta)\le 10=5\mathfrak{f}_{v}$. This proves $\delta_{v}\le 5\mathfrak{f}_{v}$ at every place of residue characteristic at least $5$. At the residue characteristics $2$ and $3$ the reduction is again classified by a finite list of types, on which $\mathrm{ord}_{v}(\Delta)$ and $\max\{0,-\mathrm{ord}_{v}(j_{E})\}$ take a bounded difference, so that $\delta_{v}\le\kappa_{v}$ for a constant determined by $v$ alone.
\end{proof}

\begin{proposition}[An unconditional decomposition of the normalized height]
\label{prop:UnconditionalNormalizedDecomposition}
Fix a number field $K$. There is a constant
\[
  C_{K,6}
  \coloneqq
  \frac{1}{[K:\Q]}
  \sum_{v\mid 6}
  \kappa_{v}\log \mathrm{N}v
\]
depending only upon $K$, such that every elliptic package $\mathcal{E}$ over $K$ satisfies
\[
  H^{\mathrm{nf}}_{K}(\mathcal{E})
  \le
  5\,n(\mathcal{E})+\Theta_{K}(\mathcal{E})+C_{K,6}.
\]
\end{proposition}

\begin{proof}
By Definition~\ref{def:MultiplicativeExcess} the minimal discriminant exponent at each finite place $v$ decomposes as
\[
  \mathrm{ord}_{v}\bigl(\Delta_{\mathrm{NF}_{K}(\mathcal{E}),v}\bigr)
  =
  \max\bigl\{0,-\mathrm{ord}_{v}(j_{E})\bigr\}+\delta_{v}(\mathcal{E}).
\]
Multiply by $\log\mathrm{N}v$, sum over the finite places, and divide by $[K:\Q]$. The first term assembles into $\Theta_{K}(\mathcal{E})$ by Definition~\ref{def:MultiplicativeExcess}, and the discriminant height of the normalized model has no archimedean contribution, so the left member is $H^{\mathrm{nf}}_{K}(\mathcal{E})$. There remains the defect sum
\[
  \frac{1}{[K:\Q]}\sum_{v}\delta_{v}(\mathcal{E})\log\mathrm{N}v.
\]
Split this sum at the residue characteristics $2$ and $3$. By Lemma~\ref{lem:PotentiallyGoodDefectBound}, the places $v\nmid 6$ contribute at most
\[
  \frac{5}{[K:\Q]}\sum_{v\nmid 6}\mathfrak{f}_{v}(\mathcal{E})\log\mathrm{N}v
  \le
  5\,n(\mathcal{E})
\]
since $\sum_{v}\mathfrak{f}_{v}\log\mathrm{N}v=\log N_{E}=[K:\Q]\,n(\mathcal{E})$ by Definition~\ref{def:ConductorComplexity}. The places $v\mid 6$ contribute at most $C_{K,6}$. Adding the three pieces gives the displayed bound.
\end{proof}

\begin{remark}[Where the difficulty is concentrated]
\label{rem:LocalizationOfDifficulty}
Proposition~\ref{prop:UnconditionalNormalizedDecomposition} is unconditional. It says that the normalized discriminant height exceeds five times the conductor complexity by no more than the multiplicative excess, up to a constant of the field. The first term lies below the Szpiro slope of six, with room to spare, so the whole weight of a normalized linear envelope rests upon a bound for $\Theta_{K}(\mathcal{E})$. The multiplicative excess is the logarithmic height of the denominator of $j_{E}$, and it grows with the depth of the potentially multiplicative reduction, which is the locus that the dilation tower of Section~\ref{subsec:DilationTower} and the obstruction of Theorem~\ref{thm:DescentObstruction} have flagged throughout as the place that no identification may erase.
\end{remark}

\begin{hypothesis}[Multiplicative excess envelope]
\label{hyp:MultiplicativeExcessEnvelope}
There exist real numbers $b_{K}\ge 0$ and $B'_{K}\ge 0$, depending only upon $K$, such that every elliptic package $\mathcal{E}$ over $K$ satisfies
\[
  \Theta_{K}(\mathcal{E})\le b_{K}\,n(\mathcal{E})+B'_{K}.
\]
\end{hypothesis}

\begin{proposition}[Any excess envelope gives a Szpiro bound]
\label{prop:ExcessEnvelopeGivesNormalizedEnvelope}
Assume Hypothesis~\ref{hyp:MultiplicativeExcessEnvelope}. Then the normalized linear envelope of Hypothesis~\ref{hyp:NormalizedLinearEnvelope} holds with
\[
  a_{K}=5+b_{K}
  \qquad\text{and}\qquad
  B_{K}=B'_{K}+C_{K,6}.
\]
Consequently every elliptic package $\mathcal{E}$ over $K$ satisfies
\[
  h_{\Delta}(\mathcal{E})
  \le
  (5+b_{K})\,n(\mathcal{E})+\defect(\mathcal{E})+B'_{K}+C_{K,6}.
\]
\end{proposition}

\begin{proof}
Insert the bound of Hypothesis~\ref{hyp:MultiplicativeExcessEnvelope} into the unconditional decomposition of Proposition~\ref{prop:UnconditionalNormalizedDecomposition}. This gives $H^{\mathrm{nf}}_{K}(\mathcal{E})\le(5+b_{K})n(\mathcal{E})+B'_{K}+C_{K,6}$, which is the normalized linear envelope with the stated constants. The displayed intensional inequality follows by Proposition~\ref{prop:LinearEnvelopeGivesIntensionalBound}.
\end{proof}

\begin{remark}[Any bound suffices, and the sharp one is abc-level]
\label{rem:AnyBoundSuffices}
Proposition~\ref{prop:ExcessEnvelopeGivesNormalizedEnvelope} converts any linear bound on the multiplicative excess into an intensional Szpiro inequality, the Szpiro slope becoming $5+b_{K}$. A coarse bound, with $b_{K}$ far from optimal, gives a coarse Szpiro inequality, while the sharp inequality at slope $6+\varepsilon$ asks for $b_{K}=1+\varepsilon$, that is, a bound on the height of the $j$-denominator by the conductor complexity with constant approaching one. This last is the modified Szpiro inequality, equivalent to the $abc$ conjecture of Section~\ref{subsec:SzpiroAbc}. The framework leaves no other gap. Once any envelope for $\Theta_{K}$ is granted, the remaining passage to the intensional inequality is proved.
\end{remark}

\subsubsection{Formal strength}\label{subsubsec:PartialProfileACA0}

\begin{remark}[The partial reductions stay within \texorpdfstring{\(\mathsf{ACA}_{0}\)}{ACA0}]
\label{rem:PartialProfileACA0}
The proofs of this subsection use finite package codes, finite sums of valuations, identity morphisms, and the transport inequality of Lemma~\ref{lem:TransportDiscriminant}. The Kodaira type at a place, and with it the defect $\delta_{v}$ and the exponents that enter $\Theta_{K}$, are computed by Tate's algorithm, a primitive recursive procedure on the package code of Definition~\ref{def:EllPkgCode}. For a finite coded list, the maximum in Definition~\ref{def:NormalizedHeightRemainder} is obtained by primitive recursion. Thus Propositions~\ref{prop:DefectNormalizedBound} and~\ref{prop:UnconditionalNormalizedDecomposition}, together with Lemma~\ref{lem:PotentiallyGoodDefectBound}, lie within the same $\mathsf{ACA}_{0}$ treatment as Proposition~\ref{prop:ACA0LogVolumeFormalization}, with real values handled as fast Cauchy names in the sense of Definition~\ref{def:UniverseP}. The remaining input, whether the normalized linear envelope of Hypothesis~\ref{hyp:NormalizedLinearEnvelope} or the multiplicative excess envelope of Hypothesis~\ref{hyp:MultiplicativeExcessEnvelope}, is a new arithmetic statement over the package codes, and no appeal to $\mathsf{ZFC}$ is involved.
\end{remark}

\subsection{The Szpiro--abc link in the intensional framework}\label{subsec:SzpiroAbc}

\begin{remark}[The Szpiro--\texorpdfstring{$abc$}{abc} link in Universes~H and~A]\label{rem:SzpiroAbcUniverses}
The schematic implication ``Szpiro $\Rightarrow abc$'' does not, in itself, commit one to the Hilbertian Universe~H. The argument requires only two ingredients. The first is a functorial way to associate to a primitive triple $(a,b,c)$ with $a+b=c$ an elliptic object $E_{(a,b,c)}$ endowed with enough structure for its discriminant, conductor, and height invariants to be controlled in terms of $\operatorname{rad}(abc)$. The second is an inequality of Szpiro type for such elliptic objects, formulated in whatever environment provides those invariants.

In Universe~H, these requirements are met by working with elliptic curves up to isomorphism over a fixed number field and with the classical invariants $\Delta_E$, $N_E$, and standard heights. In Universe~A, the role of $E_{(a,b,c)}$ may be played instead by an intensional package $X_{(a,b,c)}$ in a quantitative category in the sense of Definition~\ref{def:QuantitativeCategory}, and the Szpiro-type inequality takes the form of a transport bound
\[
  H\bigl(X_{(a,b,c)}^{\mathrm{tgt}}\bigr) \le H\bigl(X_{(a,b,c)}^{\mathrm{src}}\bigr) + c\!\left(f_{(a,b,c)}\right)
\]
for an appropriate coercive morphism $f_{(a,b,c)}$ and a height functor $H$. Provided that the construction $(a,b,c) \mapsto X_{(a,b,c)}$ exhibits sufficient control, and that the cost $c(f_{(a,b,c)})$ is uniformly bounded in terms of $\operatorname{rad}(abc)$ subject to the standard $\varepsilon$-loss, then the deduction of an $abc$ bound from the Szpiro inequality proceeds entirely within the internal logic of Universe~A. The obstruction of Lemma~\ref{lemma:ToyNoBridge} and Theorem~\ref{thm:DescentObstruction} concerns only the further prospect of reading the same labeled comparison simultaneously as a transport statement in Universe~A and as a classical inequality in Universe~H. It does not impede the internal Universe~A deduction, where the cost is kept rather than collapsed.
\end{remark}

%==================================================================
\section{An explicit example}\label{sec:ExplicitExample}
%==================================================================

The quantitative category of elliptic packages affords an instance in which one may follow the variation of the model defect in a numerically accessible way. It is natural to choose a curve whose conductor is governed by a prime that is not among the very smallest, so that the local behavior is influenced by a slightly larger place. The computation below draws the discriminant height, the conductor complexity, and the model defect into a single inequality, and shows that the transport bound of Lemma~\ref{lem:TransportDiscriminant} is sharp for the particular coordinate change chosen.

Consider the elliptic curve
\[
  E \colon y^2 + y = x^3 - x
\]
defined over $K = \Q$. Identified as $37\mathrm{a}1$ within the Cremona database~\cite{Cremona1997}, this equation is the global minimal Weierstrass model in the sense of Silverman~\cite{SilvermanAEC}. It has conductor $N_E = 37$ and minimal discriminant $\Delta_{\min}(E) = 37$, so that the unique prime of bad reduction is the prime $37$. Let $\mathcal{E}_0^{\mathrm{can}}$ denote the canonical normalized package $\mathrm{NF}_\Q$ attached to $E$ in the sense of Definition~\ref{def:AccumulatedDefect}. By construction this package carries the global minimal Weierstrass model
\[
  \mathcal{W}_0 \colon y^2 + y = x^3 - x
\]
together with a natural choice of local sections and a fixed label.

To isolate the contribution of the bad place, it is convenient to consider the truncated discriminant height of Definition~\ref{def:DiscriminantHeight} at the prime $v = 37$. One sets
\[
  h_\Delta^{(37)}(\mathcal{E}) \coloneqq
  \log \max\bigl(1,\lvert \Delta_{\mathcal{W},37} \rvert_{37}^{-1}\bigr)
\]
for any elliptic package $\mathcal{E}$ whose underlying curve is $E$ and whose Weierstrass model at $37$ is denoted $\mathcal{W}$. For the canonical package this gives
\begin{equation}\label{eq:hDelta2-canonical}
  h_\Delta^{(37)}\bigl(\mathcal{E}_0^{\mathrm{can}}\bigr)
  =
  \log \max\bigl(1,\lvert \Delta_{\mathcal{W}_0,37} \rvert_{37}^{-1}\bigr)
  =
  \log 37
\end{equation}
since $\Delta_{\mathcal{W}_0,37} = 37$ and the normalized absolute value satisfies $\lvert 37 \rvert_{37} = 37^{-1}$.

A second elliptic package is now produced by modifying the Weierstrass model at the place $37$ through a simple dilation. Over $\Q_{37}$ consider the change of variables
\begin{equation}\label{eq:explicit-change-of-variables-37}
  x = 37^{-2} x'
  \qquad\text{and}\qquad
  y = 37^{-3} y'
\end{equation}
on the affine plane. After replacing $(x,y)$ by $(37^{-2} x', 37^{-3} y')$ in the equation of $E$ and simplifying, one obtains a new Weierstrass equation $\mathcal{W}_1$ for $E$ over $\Q_{37}$. The coefficients of $\mathcal{W}_1$ need not be written out, for what matters is the transformation of the discriminant, governed by Lemma~\ref{lem:LocalDiscriminantChange}. That lemma yields $\Delta_{\mathcal{W}_1,37} = u^{-12}\,\Delta_{\mathcal{W}_0,37}$ with $u = 37^{-1}$, so that
\[
  \Delta_{\mathcal{W}_1,37} = (37^{-1})^{-12} \cdot 37 = 37^{12}\cdot 37 = 37^{13}
\]
and $\lvert \Delta_{\mathcal{W}_1,37} \rvert_{37} = 37^{-13}$, hence $\lvert \Delta_{\mathcal{W}_1,37} \rvert_{37}^{-1} = 37^{13}$. The truncated discriminant height at $37$ for the new model is therefore
\begin{equation}\label{eq:hDelta2-new}
  h_\Delta^{(37)}(\mathcal{E}_1)
  =
  \log \max\bigl(1,\lvert \Delta_{\mathcal{W}_1,37} \rvert_{37}^{-1}\bigr)
  =
  13 \log 37.
\end{equation}

The change of variables~\eqref{eq:explicit-change-of-variables-37} defines a morphism of elliptic packages $f \colon \mathcal{E}_0^{\mathrm{can}} \to \mathcal{E}_1$ in the sense of Definition~\ref{def:EllipticPackage}. At $v=37$ the factor $c_{37}(f)$ of Definition~\ref{def:EllipticPackage} is $c_{37}(f) = u^{-12} = (37^{-1})^{-12} = 37^{12}$. Thus, by Definition~\ref{def:ModelDefect},
\begin{equation}\label{eq:explicit-mu}
  \mu(f)
  =
  \log^{+}\!\bigl(\lvert c_{37}(f) \rvert_{37}^{-1}\bigr)
  =
  \log^{+}\!\bigl(\lvert 37^{12} \rvert_{37}^{-1}\bigr)
  =
  \log(37^{12})
  =
  12 \log 37
\end{equation}
where we have used that $\lvert 37^{12}\rvert_{37}=37^{-12}$, hence $\lvert 37^{12}\rvert_{37}^{-1}=37^{12}>1$. Lemma~\ref{lem:TransportDiscriminant} then gives, after restriction to the contribution at $v=37$,
\[
  h_\Delta^{(37)}(\mathcal{E}_1)
  \le
  h_\Delta^{(37)}\bigl(\mathcal{E}_0^{\mathrm{can}}\bigr)
  +
  \mu(f).
\]
Inserting the values from~\eqref{eq:hDelta2-canonical}, \eqref{eq:hDelta2-new}, and~\eqref{eq:explicit-mu} yields the sharp identity
\[
  13\log 37
  \le
  \log 37 + 12\log 37
  =
  13\log 37.
\]
For this choice of local coordinate change at $37$, the transport inequality of Lemma~\ref{lem:TransportDiscriminant} holds with equality at the truncated level $v=37$, since the entire change in $h_\Delta^{(37)}$ is accounted for exactly by the defect term $\mu(f)$ computed in~\eqref{eq:explicit-mu}.

The conductor complexity of $\mathcal{E}_0^{\mathrm{can}}$ and of $\mathcal{E}_1$ in the sense of Definition~\ref{def:ConductorComplexity} is $n(\mathcal{E}_0^{\mathrm{can}}) = n(\mathcal{E}_1) = \log N_E = \log 37$. The three quantities $h_\Delta(\mathcal{E})$, $n(\mathcal{E})$, and $\defect(\mathcal{E})$ of Section~\ref{sec:IntensionalSzpiro} are thus brought together within a single frame of reference. The discriminant height changes under a visible modification of the local model at $37$, while the conductor complexity stays anchored to the underlying curve $E$. The morphism $f$ keeps the passage from one elliptic package to the other, and its cost $\mu(f)$ gives an explicit upper bound for the accumulated defect in the sense of Definition~\ref{def:AccumulatedDefect}. Any morphism $g \colon \mathcal{E}_0^{\mathrm{can}} \to \mathcal{E}_1$ in $\mathbf{EllPkg}_\Q$ provides an upper bound $\defect(\mathcal{E}_1) \le \mu(g)$, so that in particular $\defect(\mathcal{E}_1) \le \mu(f) = 12 \log 37$.

Under Hypothesis~\ref{hyp:ReducedNormalizedProfile}, Proposition~\ref{prop:NormalizedToIntensionalSzpiro} yields that for every real number $\varepsilon > 0$ there exists a constant $C_{\varepsilon,\Q}$ such that
\begin{equation}\label{eq:explicit-intensional-szpiro}
  h_\Delta(\mathcal{E}_1)
  \le
  (6+\varepsilon)\,n(\mathcal{E}_1)
  +
  \defect(\mathcal{E}_1)
  +
  C_{\varepsilon,\Q}.
\end{equation}
Combining the sharp identity above with the relation $n(\mathcal{E}_1)=\log 37$ and the estimate $\defect(\mathcal{E}_1)\le 12\log 37$ shows that \eqref{eq:explicit-intensional-szpiro} may be realized with a suitable choice of $C_{\varepsilon,\Q}$. The demonstrative value of the calculation is the simultaneous integration of the height, the conductor, and the model defect into the architecture of a single inequality.

This higher-conductor example over $\Q$ is elementary from a Diophantine perspective. It serves, however, to exhibit the behavior of the discriminant height under a controlled deformation at a non-trivial prime, with a minimum of technical overhead. In this way it illustrates how an intensional Szpiro-type relation in the sense of Proposition~\ref{prop:NormalizedToIntensionalSzpiro} arises from transport inequalities in a quantitative category of elliptic packages.

%==================================================================
\section{Descent to a pre-set-theoretic universe}\label{sec:UniverseP}
%==================================================================

Arithmetic practice remains attentive to numbers, primes, and finite packets of coefficients. Many of the statements that matter -- such as the Szpiro inequality and the $abc$ conjecture -- permit a formulation that appeals solely to such finite data, without invoking any general facility to form infinite sets of arbitrary objects. It is appropriate, therefore, to circumscribe an arithmetic universe that forbids the completed totalities which would erase intensional distinctions, while remaining a natural environment for working number theory. This universe should be regarded as a convenient internal presentation of objects and uniformity statements that are intended to be formalizable in $\mathsf{ACA}_0$, rather than as a separate foundational system.

\subsection{A pre-set-theoretic arithmetic universe}\label{subsec:PreSetArithmetic}

\begin{definition}[The pre-set-theoretic arithmetic universe]\label{def:UniverseP}
Fix $\mathsf{ACA}_0$ as the ambient theory, in the language of second-order arithmetic. The \emph{pre-set-theoretic arithmetic universe}, denoted $\mathfrak{P}$ and called Universe~P, is the category of arithmetically presented codes defined internally to $\mathsf{ACA}_0$. An object $X$ of $\mathfrak{P}$ consists of a code $c_X \subseteq \mathbb{N}$ representing a countable many-sorted presentation that includes the following data:
\begin{enumerate}
\item[(i)] for each sort $\sigma$, a domain $D_X^\sigma \subseteq \mathbb{N}$ defined by an arithmetic predicate relative to $c_X$,
\item[(ii)] for each basic relation symbol $R$, an arithmetic predicate relative to $c_X$ representing a subset of a finite product of these domains,
\item[(iii)] for each basic function symbol $F$, an arithmetic graph relative to $c_X$ representing a total function on the domains,
\item[(iv)] for each quantitative datum in the presentation, either a rational code or a fast Cauchy name of rationals, arithmetical in $c_X$, when a real-valued quantity is required.
\end{enumerate}
A morphism $f \colon X \to Y$ is a code $c_f \subseteq \mathbb{N}$ whose graph is arithmetical in $c_X \oplus c_Y \oplus c_f$, and one requires that $\mathsf{ACA}_0$ prove this code defines a total function respecting the structure of $X$ and $Y$. Composition is derived from the composition of coded graphs, and the identity morphisms are coded diagonals -- for an object $X$, the identity is the arithmetic predicate identifying the set of pairs $\langle n, n \rangle$ for every $n$ in the domains of $X$. Self-identification within this world is not a background fact, and it is rather a program requiring a definite arithmetic representation.
\end{definition}

\begin{proposition}[Formalizability of the $\mathfrak{P}$-constructions in \texorpdfstring{$\mathsf{ACA}_0$}{ACA0}]\label{prop:ACAFormalizable}
All constructions on $\mathfrak{P}$ used in the descent of Proposition~\ref{prop:APDescentToP} are formalizable in $\mathsf{ACA}_0$. More precisely, within $\mathsf{ACA}_0$ one can form and manipulate:
\begin{enumerate}
\item[(i)] finite tuples, words, matrices, and other standard combinatorial codes,
\item[(ii)] arithmetically defined domains, graphs, and subobjects associated with a presentation code,
\item[(iii)] pullbacks, finite products, and other finite-limit constructions on arithmetically presented objects,
\item[(iv)] rational-valued cost functions and real-valued quantities represented by fast Cauchy names of rationals,
\item[(v)] the comparisons, bounded searches, and finite minimization procedures required to descend from extensional cost data to combinatorial bounds.
\end{enumerate}
Hence the transcription of Proposition~\ref{prop:APDescentToP} may be read as a theorem internal to $\mathsf{ACA}_0$, once its hypotheses are expressed in the coding formalism of $\mathfrak{P}$.
\end{proposition}

\begin{proof}
All finite coding operations used in the definition of $\mathfrak{P}$ are primitive recursive. The only non-trivial set-existence principle required in passing from a presentation to its associated domains, graphs, and subobjects is the formation of sets defined by arithmetic predicates, which is exactly what arithmetical comprehension provides. Real-valued quantities are represented by fast Cauchy names of rationals, and the arithmetic manipulation of such names is standard in $\mathsf{ACA}_0$. The descent argument uses only these coding operations, arithmetically definable graphs, finite search procedures, and comparisons of coded rational or Cauchy data. No stronger comprehension scheme, and no higher-order set-existence principle, is invoked.
\end{proof}

\subsection{Heights and defects as finite integer data}\label{subsec:LogIntData}

\begin{definition}[Logarithmic integer datum]\label{def:LogIntDatum}
Let $K$ be a number field presented by finite data in Universe~P. A \emph{logarithmic integer datum over $K$} is a positive integer $M$. It represents the real number $\frac{1}{[K:\Q]}\,\log M$ and is read as an element of $\R_{\ge 0}$ whenever a numerical evaluation is required. A comparison between two logarithmic integer data $M_1$ and $M_2$ over the same field reduces to the integer inequality $M_1 \le M_2$. More generally, an inequality of the form
\[
  \frac{1}{[K:\Q]}\,\log M_1
  \le
  \bigl(6+\tfrac{a}{b}\bigr)\,\frac{1}{[K:\Q]}\,\log M_2
  +
  \frac{1}{[K:\Q]}\,\log M_3
  +
  C
\]
with $a,b$ positive integers and $C$ a real constant admits an \emph{exponentiated integer reformulation}. One multiplies through by $b\,[K:\Q]$ and exponentiates, producing the comparison
\[
  M_1^b \le D \cdot M_2^{6b+a} \cdot M_3^b
\]
where $D$ is any positive integer satisfying $D \ge \lceil e^{\,b\,[K:\Q]\,C} \rceil$. This last expression is a comparison of positive integers and lies entirely within the arithmetic of Universe~P. The integer $D$ is determined by a finite computation once $K$ and the rational precision $a/b$ have been fixed.
\end{definition}

\begin{remark}[Heights and defects as logarithmic integer data]\label{rem:HeightsAsLogIntData}
The three principal quantities of the intensional Szpiro inequality admit a natural representation as logarithmic integer data, at least for their non-archimedean contributions. For an elliptic package $\mathcal{E}$ over $K$ with Weierstrass model $\mathcal{W}$, the non-archimedean part of the discriminant height $h_\Delta(\mathcal{E})$ of Definition~\ref{def:DiscriminantHeight} satisfies
\[
  [K:\Q]\,h_\Delta^{\mathrm{fin}}(\mathcal{E})
  =
  \log\!\Bigl(\prod_{v\,\text{finite}} p_v^{\,f_v\,\max(0,\,\mathrm{ord}_v(\Delta_{\mathcal{W}}))}\Bigr)
\]
where $p_v$ is the rational prime below $v$ and $f_v$ is the residue degree. The product on the right is a positive integer computable from the Weierstrass coefficients, and the archimedean contribution is bounded by a constant depending only upon $K$ and upon the number of archimedean places, hence may be pre-computed and bounded by a rational number that is part of the field code. The conductor complexity $n(\mathcal{E}) = (1/[K:\Q])\,\log N_E$ is a logarithmic integer datum in the strict sense, since $N_E$ is a positive integer. For a morphism $f$ of elliptic packages, each local factor $|c_v(f)|_v^{-1}$ is a power of the residue characteristic, so that $[K:\Q]\,\mu(f) = \log M_f$ for a positive integer $M_f$ determined by the local dilation data.
\end{remark}

\begin{definition}[Certified number field code]\label{def:CertifiedFieldCode}
Within Universe~P, a \emph{certified number field code} for a number field $K$ of degree $d = [K:\Q]$ is a finite tuple of integer data comprising the following items:
\begin{enumerate}
\item[\textit{(i)}] a monic irreducible polynomial $f(x) \in \Z[x]$ of degree $d$, together with a fixed ordering of its roots in $\C$ that determines the real and complex embeddings of $K$,
\item[\textit{(ii)}] a \emph{certified integral basis}, a list of $d$ elements $(\omega_1,\dots,\omega_d)$ of $K$, each presented as a polynomial in the root $\alpha$ of $f$ with rational coefficients, such that $\mathcal{O}_K = \Z\omega_1 + \cdots + \Z\omega_d$, the certification consisting of the field discriminant $d_K \in \Z$ and a finite proof witness that confirms the maximality of the order $\Z[\omega_1,\dots,\omega_d]$ (for instance a transcript of the steps of the Round Four algorithm of Ford--Zassenhaus as in~\cite{Cohen1993}),
\item[\textit{(iii)}] a \emph{ramification table}, for each rational prime $p$ dividing $d_K$ a finite list of tuples $(e_{v}, f_{v}, \pi_{v})$ keeping the ramification index, the residue degree, and a uniformizer for each prime ideal $v$ of $\mathcal{O}_K$ above $p$, the uniformizer being encoded as a $\Z$-linear combination of the integral basis,
\item[\textit{(iv)}] a rational upper bound $B_K^{\mathrm{arch}} \in \Q_{>0}$ for the archimedean contribution to the discriminant height $h_\Delta^{\mathrm{arch}}$ of any elliptic package over $K$, computed from the embeddings in \textit{(i)} and kept as part of the field code.
\end{enumerate}
The purpose of this enriched presentation is to ensure that the computation of conductors, Tamagawa numbers, local dilation factors, and normalization morphisms proceeds within Universe~P by deterministic operations on the certified data, without requiring the execution of maximal-order algorithms at runtime.
\end{definition}

The descent from Universe~A to Universe~P admits an explicit transcription. The principle is lucid. An elliptic package comprises finitely many coefficients, finitely many local parameters, and a finite label. A morphism consists of finitely many coordinate changes, and a defect appears as a sum of finitely many local contributions. All of this registers as finite arithmetic data, manipulable without quantification over any global class of objects.

\begin{definition}[Arithmetic code for an elliptic package]\label{def:EllPkgCode}
Let $K$ be a number field equipped with a certified number field code in the sense of Definition~\ref{def:CertifiedFieldCode}. An \emph{elliptic package code over $K$} is a finite tuple of integers encoding the following data:
\begin{enumerate}
\item[\textit{(i)}] the coefficients $[a_1,a_2,a_3,a_4,a_6]$ of a Weierstrass equation for $E$, each expressed as a $\Z$-linear combination of the certified integral basis $(\omega_1,\dots,\omega_d)$ from the field code,
\item[\textit{(ii)}] a finite list representing the set $S$ of marked places, by listing the prime ideals above the corresponding rational primes through the indices in the ramification table of the field code and a tag for each archimedean place,
\item[\textit{(iii)}] the finite parameters describing the local sections or level structures $\{\sigma_v\}_{v\in S}$, each encoded as a tuple of elements of $\mathcal{O}_K$ relative to the certified basis,
\item[\textit{(iv)}] a finite label $\lambda$, transcribed as a finite string over a fixed alphabet.
\end{enumerate}
A \emph{morphism code} from a package code $d$ to a package code $d'$ over the same field $K$ is a finite tuple encoding the coordinate-change parameters $(u_v, r_v, s_v, t_v)$ that appear in~\eqref{eq:local-change-of-variables} for each relevant place, each presented as an element of the local field $K_v$ through its expansion in the certified basis. When one works in the variable-base category $\mathbf{EllPkg}_{\mathrm{var}}$ of Definition~\ref{def:EllPkgVar}, then a morphism code from a package over $K$ to a package over $L = K(\sqrt[\ell]{q})$ carries in addition a \emph{field extension code} -- the minimal polynomial of $\sqrt[\ell]{q}$ over $K$, a certified integral basis for $\mathcal{O}_L$ extending the basis of $\mathcal{O}_K$, and the polynomial-basis matrix expressing the old basis elements in terms of the new. The Kummer detachment cost $\mu_{\mathrm{kum}}(f_{\mathrm{rad}})$ of Definition~\ref{def:KummerDefect} is then computed from the degree $\ell$ alone, and the descent cost of Hypothesis~\ref{hyp:ControlledDescentWithSupport} is determined by the ramification data of the certified extension code, together with the tame support of the radicand. One writes $\mathsf{PkgCode}_K$ for the finite-type domain of such codes and $\mathsf{PkgCode}_{\mathrm{var}}$ for the variable-base variant.
\end{definition}

\begin{remark}[Computability of invariants from package codes]\label{rem:PkgCodeComputability}
Universe~P admits function symbols for arithmetic invariants computed from codes. The following quantities are determined by finite, deterministic procedures operating on the data of Definitions~\ref{def:EllPkgCode} and~\ref{def:CertifiedFieldCode}. The non-archimedean part of $h_\Delta$ is computed from the Weierstrass coefficients and the ramification table by evaluating valuations at each prime ideal of $\mathcal{O}_K$ dividing the discriminant, and by Remark~\ref{rem:HeightsAsLogIntData} this part is a logarithmic integer datum, while the archimedean part is bounded by the pre-computed constant $B_K^{\mathrm{arch}}$ from the certified field code. The conductor $N_E$ is computed by Ogg's formula from the local reduction types at each prime in the ramification table, so that the conductor complexity is a logarithmic integer datum in the strict sense. The defect $\mu(f)$ of a morphism code $f$ is computed from the local dilation factors as in Definition~\ref{def:ModelDefect}, each $|u_v|_v$ being a power of the residue characteristic read from the uniformizer data. The normalization map $\mathrm{NF}_K$ and the canonical morphism $\nu_{\mathcal{E}}$ of Definition~\ref{def:AccumulatedDefect} are executed by Tate's algorithm following the deterministic discipline of Remark~\ref{rem:DeterministicNormalizationDiscipline}. No operation above requires quantification over an infinite collection or appeals to the existence of a completed totality, since all inputs and outputs are finite tuples of integers and all intermediate steps are bounded by the degree $[K:\Q]$ and by the number of primes dividing the conductor.
\end{remark}

\subsection{The transcription}\label{subsec:Transcription}

\begin{proposition}[Transcription of intensional Szpiro into Universe~P]\label{prop:APDescentToP}
Fix a number field $K$ and a normalization procedure $\mathrm{NF}_K$ in the sense of Definition~\ref{def:AccumulatedDefect}, both presented by finite data. Then there exists a Universe~P formula $\mathrm{Szp}^P_{\varepsilon,K}(d)$ with $d \in \mathsf{PkgCode}_K$ as its only free variable, possessing the following property. For every elliptic package $\mathcal{E}$ over $K$ with code $d$, the inequality
\[
h_\Delta(\mathcal{E})
\le
(6+\varepsilon)\,n(\mathcal{E})
+
\defect(\mathcal{E})
+
C_{\varepsilon,K}
\]
from~\eqref{eq:IntensionalSzpiroBasic} holds if and only if $\mathrm{Szp}^P_{\varepsilon,K}(d)$ holds when the terms $h_\Delta$, $n$, and $\defect$ are interpreted as the corresponding functions on codes. Consequently, the basic intensional Szpiro inequality may be phrased in Universe~P as a family of inequalities over finite codes, devoid of quantification over a totality of all elliptic curves.
\end{proposition}

\begin{proof}
By Definition~\ref{def:EllPkgCode}, a package is determined by a finite tuple $d$. The quantities $h_\Delta(\mathcal{E})$ and $n(\mathcal{E})$ are computed from $d$ by the explicit procedures of arithmetic geometry, and $\defect(\mathcal{E})$ is computed from $d$ by way of the normalization procedure of Definition~\ref{def:AccumulatedDefect} together with the defect $\mu$ of Definition~\ref{def:ModelDefect}. One may thus write a Universe~P formula asserting the inequality among these computed numbers, and no further ontological commitment is required.
\end{proof}

\begin{remark}[Universe~P as recipient]\label{rem:PAsTarget}
Universe~P acts as a natural recipient for the structures of Universe~A. Universe~A keeps transport and defect, while Universe~P stores the resulting numerical inequalities on codes and declines to enforce static identification in the sense of Definition~\ref{def:StaticDynamic} by fiat. In this manner the descent from A to P bypasses the extensional collapse of Definition~\ref{def:ExtensionalCollapse} entirely.
\end{remark}

\begin{proposition}[Soundness of descent along normalization]\label{prop:SoundnessDescentNormalization}
Fix a number field $K$ and a normalization procedure $\mathrm{NF}_K$ in the sense of Definition~\ref{def:AccumulatedDefect}, presented by finite data in Universe~P, and assume Hypothesis~\ref{hyp:ReducedNormalizedProfile}. Then for every real number $\varepsilon>0$ there exists a constant $C_{\varepsilon,K}$ such that for every package code $d \in \mathsf{PkgCode}_K$ one has the Universe~P inequality
\[
h_\Delta\bigl(\mathrm{NF}_K(d)\bigr)
\le
(6+\varepsilon)\,n\bigl(\mathrm{NF}_K(d)\bigr) + C_{\varepsilon,K}
\]
where $\mathrm{NF}_K(d)$ is the code obtained by applying the normalization algorithm to $d$, and $h_\Delta$ and $n$ are interpreted as the corresponding functions on codes.
\end{proposition}

\begin{proof}
Let $d$ be a package code and let $\mathcal{E}$ be an elliptic package over $K$ with code $d$. Consider the normalized package $\mathrm{NF}_K(\mathcal{E})$. By Definition~\ref{def:AccumulatedDefect}\textup{\ref{item:NF_idempotent}} one has $\defect(\mathrm{NF}_K(\mathcal{E}))=0$, so that the intensional inequality of Proposition~\ref{prop:NormalizedToIntensionalSzpiro} specializes to $h_\Delta(\mathrm{NF}_K(\mathcal{E})) \le (6+\varepsilon)\,n(\mathrm{NF}_K(\mathcal{E})) + C_{\varepsilon,K}$. Transcribe this inequality to codes using Definition~\ref{def:EllPkgCode} and Remark~\ref{rem:PkgCodeComputability}, and use that $\mathrm{NF}_K(d)$ is the code of $\mathrm{NF}_K(\mathcal{E})$.
\end{proof}

\subsection{How the coherence data is discharged}\label{subsec:TwoMorphismErasure}

\begin{remark}[How $2$-morphisms are resolved in the descent to finite codes]\label{rem:TwoMorphismErasure}
The passage from the labeled quantitative $(2,1)$-category $\mathfrak{A}$ of Definition~\ref{def:UniverseAFormal} to the finite codes of Definition~\ref{def:EllPkgCode} traverses the homotopy $1$-category $\Ho(\mathfrak{A})$ of Definition~\ref{def:HoOneCategory}, and this intermediate step absorbs the $2$-morphisms. The descent may be organized in three stages. In the first, by Definition~\ref{def:HoOneCategory}, the homotopy $1$-category $\Ho(\mathfrak{A})$ retains the same objects as $\mathfrak{A}$ and replaces each $1$-morphism by its $2$-isomorphism class. By Proposition~\ref{prop:CoherenceStability}\textup{(i)} the cost factors through this quotient, every invertible $2$-morphism $\alpha \colon f \Rightarrow g$ is collapsed to the identity $[f] = [g]$, and the numerical datum $c(f) = c(g)$ attached to the class is well defined, so that no cost accrues from the collapse. In the second stage, within each $2$-isomorphism class one selects a representative $1$-morphism. In the elliptic-package fragment of Universe~A of Section~\ref{sec:EllipticPackages}, this selection is governed by the normalization procedure $\mathrm{NF}_K$ of Definition~\ref{def:AccumulatedDefect}, which fixes a canonical presentation within each extensional fiber, and the choice depends only upon the finitely presented data of the source and target packages by the deterministic tie-breaking discipline of Definition~\ref{def:AccumulatedDefect}\textup{\ref{item:NF_finitary}}. In the third stage, the selected representative is transcribed into a tuple of integers by Definition~\ref{def:EllPkgCode}, and from these data one computes the cost $\mu(f)$ and the discriminant height $h_\Delta$ by the procedures of Remark~\ref{rem:PkgCodeComputability}. The net effect is that the $2$-morphisms of $\mathfrak{A}$ leave no residue in the finite code. Their role was to witness that two $1$-morphisms differed only by a coherent rearrangement of auxiliary choices, and Proposition~\ref{prop:CoherenceStability}\textup{(ii)} ensures that such rearrangements do not alter the cost. Once the cost has been kept as a rational number or a computable real in $\mathsf{PkgCode}_K$, the coherence data has been fully discharged, and the transport inequality of Lemma~\ref{lem:TransportDiscriminant} and the intensional Szpiro bound of Proposition~\ref{prop:NormalizedToIntensionalSzpiro} operate on the finite codes exactly as they would in $\Ho(\mathfrak{A})$.
\end{remark}

\begin{remark}[Compatibility with the descent obstruction]\label{rem:SoundnessVsNoBridge}
Proposition~\ref{prop:SoundnessDescentNormalization} does not contradict Lemma~\ref{lemma:ToyNoBridge} or Theorem~\ref{thm:DescentObstruction}. The obstruction requires a two-sided comparison passing through a functorial, label-transparent bridge in the sense of Definition~\ref{def:ToyAliasedStandard}, and collapses it to $H \le H + \Delta$ once the labels have been erased. The descent to Universe~P avoids this collapse by declining to perform the identification. In Universe~P the primary objects are finite codes $d \in \mathsf{PkgCode}_K$, equality between two codes is typographical, and two codes encoding isomorphic elliptic curves through different Weierstrass models remain distinct. The discriminant height computed from each code reflects the specific dilation factors that produced it, and the model defect $\mu(f)$ is itself a determination on morphism codes that keeps the structural friction between two presentations as an explicit numerical offset. Moreover the inequality that descends to Universe~P is one-sided. It states $h_{\Delta}(\mathrm{NF}_K(d)) \le (6+\varepsilon)\,n(\mathrm{NF}_K(d)) + C_{\varepsilon,K}$ for each code $d$, and it is not a comparison between two evaluations of a single invariant attached to an isomorphism class. No label-transparent bridge of the kind whose collapse Theorem~\ref{thm:DescentObstruction} governs is invoked.
\end{remark}

\begin{remark}[A brief link back to Universe~H]\label{rem:LinkBackToH}
If one chooses to read $\mathrm{NF}_K(\mathcal{E})$ as a canonical presentation of the underlying curve in the extensional sense, then Proposition~\ref{prop:SoundnessDescentNormalization} may be read as an extensional Szpiro-type inequality for that chosen representative. This form of harmonization is optional, since Universe~P is a sufficient recipient for arithmetic statements once they are expressed on finite codes.
\end{remark}

%==================================================================
\section{Formalization sketches in Coq}\label{sec:ToyFormalization}
%==================================================================

The Calculus of Inductive Constructions, as implemented in the Coq system, offers a place where one may test whether a quantitative argument truly survives the discipline of a small kernel. The kernel itself is indifferent to the arithmetic meaning of a transport. If one wishes to distinguish a transport of zero cost from a transport that carries a visible distortion, then this distinction must be written into the data.

In the present paper, the native environment is a labeled quantitative $(2,1)$-category in the sense of Definition~\ref{def:UniverseAFormal}. It is therefore natural, in a toy formalization, to keep $2$-morphisms literal and to impose the two requirements that matter most for the numerical accounting:
\begin{subtlelist}
\item Every $2$-morphism is invertible.
\item The distortion cost is invariant under $2$-isomorphism.
\end{subtlelist}
Under these conditions, coherence carries no numerical weight. It remains present as structure, while the scalar bounds attach only to $1$-morphisms.

\subsection{A minimal quantitative $(2,1)$-category in Coq}\label{subsec:CoqQuantCat}

To fix ideas, we present a small interface for a strict $(2,1)$-category equipped with a cost and a height. The code is written as explanatory scaffolding. Several laws are stated as axioms, and the interface is kept minimal, while it remains enough to express the transport inequalities that recur throughout the paper.

\begin{scriptsize}
\begin{verbatim}
(*
  Toy quantitative (2,1)-categories.

  The interface below is written to mirror the data of a labeled quantitative
  (2,1)-category in the sense of Definition~\ref{def:UniverseAFormal}.

  Conventions:
    - comp1 f g means "first f, then g"
    - 2-cells are written as Two f g, and are all invertible
*)

From Coq Require Import Reals.
From Coq Require Import Lra.

Open Scope R_scope.

Record QuantTwoOneCat := {
  Obj : Type ;
  Hom : Obj -> Obj -> Type ;

  (* 2-cells between parallel 1-morphisms *)
  Two : forall {A B : Obj}, Hom A B -> Hom A B -> Type ;

  (* 1-categorical structure *)
  id1   : forall A : Obj, Hom A A ;
  comp1 : forall {A B C : Obj}, Hom A B -> Hom B C -> Hom A C ;

  comp1_assoc :
    forall (A B C D : Obj)
           (f : Hom A B) (g : Hom B C) (h : Hom C D),
      comp1 (comp1 f g) h = comp1 f (comp1 g h) ;

  comp1_id_l :
    forall (A B : Obj) (f : Hom A B),
      comp1 (id1 A) f = f ;

  comp1_id_r :
    forall (A B : Obj) (f : Hom A B),
      comp1 f (id1 B) = f ;

  (* 2-categorical structure, strict for convenience *)
  id2    : forall {A B : Obj} (f : Hom A B), Two f f ;
  vcomp2 : forall {A B : Obj} {f g h : Hom A B},
             Two f g -> Two g h -> Two f h ;

  vcomp2_assoc :
    forall (A B : Obj) (f g h k : Hom A B)
           (a : Two f g) (b : Two g h) (c : Two h k),
      vcomp2 a (vcomp2 b c) = vcomp2 (vcomp2 a b) c ;

  vcomp2_id_l :
    forall (A B : Obj) (f g : Hom A B) (a : Two f g),
      vcomp2 (id2 f) a = a ;

  vcomp2_id_r :
    forall (A B : Obj) (f g : Hom A B) (a : Two f g),
      vcomp2 a (id2 g) = a ;

  (* whiskering *)
  whiskerL :
    forall {A B C : Obj} (f : Hom A B) {g h : Hom B C},
      Two g h -> Two (comp1 f g) (comp1 f h) ;

  whiskerR :
    forall {A B C : Obj} {f g : Hom A B} (a : Two f g) (h : Hom B C),
      Two (comp1 f h) (comp1 g h) ;

  (* (2,1)-condition: explicit inverses for 2-cells *)
  inv2 :
    forall {A B : Obj} {f g : Hom A B}, Two f g -> Two g f ;

  inv2_left :
    forall {A B : Obj} {f g : Hom A B} (a : Two f g),
      vcomp2 a (inv2 a) = id2 f ;

  inv2_right :
    forall {A B : Obj} {f g : Hom A B} (a : Two f g),
      vcomp2 (inv2 a) a = id2 g ;

  (* distortion cost on 1-morphisms *)
  cost : forall {A B : Obj}, Hom A B -> R ;

  cost_id :
    forall A : Obj, cost (id1 A) = 0 ;

  cost_comp :
    forall (A B C : Obj) (f : Hom A B) (g : Hom B C),
      cost (comp1 f g) <= cost f + cost g ;

  (* invariance under 2-isomorphism *)
  cost_two_invariant :
    forall (A B : Obj) (f g : Hom A B), Two f g -> cost f = cost g ;

  (* height on objects *)
  height : Obj -> R ;

  (* transport inequality *)
  height_trans :
    forall (A B : Obj) (f : Hom A B),
      height B <= height A + cost f
}.
\end{verbatim}
\end{scriptsize}

From this interface, the familiar consequences of the transport inequality follow directly. The axiom \texttt{cost\_two\_invariant} ensures that the scalar cost depends only upon a $1$-morphism up to coherent change of presentation, and the invertible $2$-cells carry no independent numerical weight.

\begin{scriptsize}
\begin{verbatim}
Section DerivedLemmas.

Context (C : QuantTwoOneCat).

Lemma height_no_increase_if_cost_zero :
  forall (A B : Obj C) (f : Hom C A B),
    cost (C:=C) f = 0 ->
    height C B <= height C A.
Proof.
  intros A B f Hc.
  pose proof (height_trans C A B f) as H.
  lra.
Qed.

Lemma height_increase_forces_cost :
  forall (A B : Obj C) (f : Hom C A B),
    height C A <= height C B ->
    height C B - height C A <= cost (C:=C) f.
Proof.
  intros A B f Hle.
  pose proof (height_trans C A B f) as H.
  lra.
Qed.

Lemma cost_equal_under_two_cell :
  forall (A B : Obj C) (f g : Hom C A B),
    Two C f g ->
    cost (C:=C) f = cost (C:=C) g.
Proof.
  intros A B f g a.
  apply (cost_two_invariant C A B f g a).
Qed.

End DerivedLemmas.
\end{verbatim}
\end{scriptsize}

\begin{remark}[Coherence is present, and it is numerically silent]\label{rem:CoqCoherenceSilent}
The lemma \texttt{cost\_equal\_under\_two\_cell} is the toy formal analogue of the requirement that the cost factors through $\Ho(\mathfrak{A})$. The $2$-cells remain present and invertible, and yet they do not perturb the scalar accounting. This silence is load-bearing rather than incidental. It is the very property on which the collapse of the descent to finite codes, described in Remark~\ref{rem:TwoMorphismErasure}, depends, and its consequence for the management of the $2$-cells in the formal interface is taken up in Remark~\ref{rem:CoqPseudoNat}.
\end{remark}

\subsection{A toy elliptic-package fragment}\label{subsec:CoqEllPkg}

We now sketch an elliptic-package fragment in which the $2$-morphisms encode harmless changes of presentation of transports. One may keep the geometric content abstract. The cost \texttt{mu} plays the role of the model defect $\mu$ of Definition~\ref{def:ModelDefect}, and the height \texttt{hDelta} plays the role of $h_\Delta$ of Definition~\ref{def:DiscriminantHeight}.

\begin{scriptsize}
\begin{verbatim}
(*
  Toy elliptic packages as a (2,1)-category.

  Objects: elliptic packages (left abstract)
  1-morphisms: transports between packages
  2-morphisms: coherent identifications between transports, all invertible
*)

Parameter EllPkg  : Type.
Parameter EllHom  : EllPkg -> EllPkg -> Type.
Parameter EllTwo  : forall {E F : EllPkg}, EllHom E F -> EllHom E F -> Type.

Parameter Ell_id1  : forall E : EllPkg, EllHom E E.
Parameter Ell_comp1 : forall {E F G : EllPkg},
  EllHom E F -> EllHom F G -> EllHom E G.

Axiom Ell_comp1_assoc :
  forall (E F G H : EllPkg)
         (f : EllHom E F) (g : EllHom F G) (h : EllHom G H),
    Ell_comp1 (Ell_comp1 f g) h = Ell_comp1 f (Ell_comp1 g h).

Axiom Ell_comp1_id_l :
  forall (E F : EllPkg) (f : EllHom E F),
    Ell_comp1 (Ell_id1 E) f = f.

Axiom Ell_comp1_id_r :
  forall (E F : EllPkg) (f : EllHom E F),
    Ell_comp1 f (Ell_id1 F) = f.

Parameter Ell_id2 : forall {E F : EllPkg} (f : EllHom E F), EllTwo f f.
Parameter Ell_vcomp2 :
  forall {E F : EllPkg} {f g h : EllHom E F},
    EllTwo f g -> EllTwo g h -> EllTwo f h.

Parameter Ell_inv2 :
  forall {E F : EllPkg} {f g : EllHom E F}, EllTwo f g -> EllTwo g f.

Axiom Ell_inv2_left :
  forall {E F : EllPkg} {f g : EllHom E F} (a : EllTwo f g),
    Ell_vcomp2 a (Ell_inv2 a) = Ell_id2 f.

Axiom Ell_inv2_right :
  forall {E F : EllPkg} {f g : EllHom E F} (a : EllTwo f g),
    Ell_vcomp2 (Ell_inv2 a) a = Ell_id2 g.

(* Model defect mu(f) *)
Parameter mu : forall {E F : EllPkg}, EllHom E F -> R.

Axiom mu_id :
  forall E : EllPkg, mu (Ell_id1 E) = 0.

Axiom mu_comp :
  forall (E F G : EllPkg) (f : EllHom E F) (g : EllHom F G),
    mu (Ell_comp1 f g) <= mu f + mu g.

Axiom mu_two_invariant :
  forall (E F : EllPkg) (f g : EllHom E F), EllTwo f g -> mu f = mu g.

(* Discriminant height *)
Parameter hDelta : EllPkg -> R.

(* Transport inequality, mirroring Lemma~\ref{lem:TransportDiscriminant} *)
Axiom hDelta_trans :
  forall (E F : EllPkg) (f : EllHom E F),
    hDelta F <= hDelta E + mu f.
\end{verbatim}
\end{scriptsize}

Once one has \texttt{mu\_two\_invariant}, the familiar consequences of the transport inequality may be written without collapsing $2$-cells into equalities.

\begin{scriptsize}
\begin{verbatim}
Lemma hDelta_no_increase_if_mu_zero :
  forall (E F : EllPkg) (f : EllHom E F),
    mu f = 0 ->
    hDelta F <= hDelta E.
Proof.
  intros E F f Hmu.
  pose proof (hDelta_trans E F f) as H.
  lra.
Qed.

Lemma hDelta_increase_forces_mu :
  forall (E F : EllPkg) (f : EllHom E F),
    hDelta E <= hDelta F ->
    hDelta F - hDelta E <= mu f.
Proof.
  intros E F f Hle.
  pose proof (hDelta_trans E F f) as H.
  lra.
Qed.
\end{verbatim}
\end{scriptsize}

\subsection{Normalization as a $2$-endofunctor with a pseudonatural comparison}\label{subsec:CoqPseudoNat}

In the body of the paper, normalization is treated as a disciplined choice of a reference point in each extensional fiber, together with a distinguished transport. In a $(2,1)$-category it is pleasant to write this as a $2$-endofunctor equipped with a pseudonatural transformation. The toy interface below keeps this shape, while leaving the coherence laws as axioms. The purpose is to keep the $2$-cells explicit, so that the formal picture matches the ambient language of Universe~A.

\begin{scriptsize}
\begin{verbatim}
(*
  A toy 2-endofunctor on the elliptic-package (2,1)-category.

  We keep only what is needed for pseudonaturality.
*)

Record EllEndo2 := {
  Fobj : EllPkg -> EllPkg ;
  Fmor : forall {E F : EllPkg}, EllHom E F -> EllHom (Fobj E) (Fobj F) ;

  F_id1 :
    forall E : EllPkg, Fmor (Ell_id1 E) = Ell_id1 (Fobj E) ;

  F_comp1 :
    forall (E F G : EllPkg) (f : EllHom E F) (g : EllHom F G),
      Fmor (Ell_comp1 f g) = Ell_comp1 (Fmor f) (Fmor g) ;

  Ftwo :
    forall (E F : EllPkg) (f g : EllHom E F),
      EllTwo f g -> EllTwo (Fmor f) (Fmor g)
}.

(*
  Pseudonatural transformation between 2-endofunctors.

  For each 1-morphism f, we store an explicit invertible 2-cell witnessing the
  pseudonaturality square.
*)

Record EllPseudoNat (F G : EllEndo2) := {
  eta_obj :
    forall E : EllPkg, EllHom (Fobj F E) (Fobj G E) ;

  eta_nat :
    forall (E F0 : EllPkg) (f : EllHom E F0),
      EllTwo
        (Ell_comp1 (eta_obj E) (Fmor F f))
        (Ell_comp1 (Fmor G f) (eta_obj F0)) ;

  eta_nat_inv :
    forall (E F0 : EllPkg) (f : EllHom E F0),
      EllTwo
        (Ell_comp1 (Fmor G f) (eta_obj F0))
        (Ell_comp1 (eta_obj E) (Fmor F f)) ;

  eta_nat_left :
    forall (E F0 : EllPkg) (f : EllHom E F0),
      Ell_vcomp2 (eta_nat E F0 f) (eta_nat_inv E F0 f)
        =
      Ell_id2 (Ell_comp1 (eta_obj E) (Fmor F f)) ;

  eta_nat_right :
    forall (E F0 : EllPkg) (f : EllHom E F0),
      Ell_vcomp2 (eta_nat_inv E F0 f) (eta_nat E F0 f)
        =
      Ell_id2 (Ell_comp1 (Fmor G f) (eta_obj F0))
}.
\end{verbatim}
\end{scriptsize}

\begin{remark}[Pseudonaturality as explicit invertible $2$-cells]\label{rem:CoqPseudoNat}
The fields \texttt{eta\_nat} and \texttt{eta\_nat\_inv} are the literal $(2,1)$-categorical content of pseudonaturality in the sense of Definition~\ref{def:PseudoNatural}. One keeps an explicit invertible $2$-morphism rather than an equality of composites.

One might ask why the interface keeps these $2$-cells together with their round-trip laws \texttt{eta\_nat\_left} and \texttt{eta\_nat\_right}, rather than showing only that the relevant squares commute. The collapse performed in the descent to finite codes, described in Remark~\ref{rem:TwoMorphismErasure}, is sound only once the cost has been shown invariant under such $2$-cells. Keeping the cells explicit is what permits that invariance to be checked by the machine through \texttt{cost\_equal\_under\_two\_cell}, rather than presumed. If the $2$-cells were collapsed at the outset, then the independence of the scalar from its presentation would be assumed in place of the very property on which the descent rests. The interface thus holds the step that establishes coherence apart from the step that reads off the scalar, and it certifies that the first does not disturb the second. This is the formal counterpart of the reading given in Remark~\ref{rem:PseudoNaturalNotBicategorical}, where the $2$-cells carry weight at the level of derivation and fall silent at the level of the final inequality.
\end{remark}

\begin{remark}[The limit of the strict prototype]\label{rem:StrictPrototypeLimit}
The strict prototype is faithful here for one reason. It transports scalars. The field \texttt{Fmor} and its laws \texttt{F\_id1} and \texttt{F\_comp1} ask the $2$-endofunctor to preserve composition strictly, and the only output the development reads from a composite is its cost, upon which an associativity isomorphism acts as the identity. A strict associator and a coherent one therefore agree on everything the numerical accounting consults, and the pseudonaturality cell \texttt{eta\_nat} keeps its full force, since it is a $2$-cell between parallel $1$-morphisms and not a matter of associativity at all. The strictness of Remark~\ref{rem:strictness} and the un-collapsed $2$-cells of Definition~\ref{def:PseudoNatural} are thus consistent, the first concerning the composition laws and the second the cells between parallel transports.

The trade-off shows its edge once the transport carries geometric data in place of a scalar. A base change of a Weierstrass model, a pullback of local sections, and the analytic uniformization of an elliptic curve compose only up to a canonical isomorphism rather than strictly, so a $2$-functor that moved such data could not satisfy \texttt{F\_comp1} as a strict equation. A faithful formalization at that level would replace the strict $2$-functors with pseudofunctors, would carry associator and unitor cells, and would require the pseudonaturality squares of Definition~\ref{def:PseudoNatural} to cohere with those cells through the bicategorical axioms. The present interface is adequate precisely because the scalars it transports -- the cost and the height -- reach Universe~P as fast Cauchy names rather than as the geometric objects whose comparison produced them, so that the coherence the strict prototype omits is coherence that the scalar never feels. A homotopy-minded reader who wishes to transport the data themselves should read the strict choice as the boundary of the prototype rather than as a claim about the geometry.
\end{remark}

\subsection{A toy intensional Szpiro interface}\label{subsec:CoqSzpiro}

Finally, one may state the schematic form of an intensional Szpiro-type inequality in this toy $(2,1)$-categorical environment. The purpose is to show that the defect term is attached to a chosen transport, while the coherence $2$-cells do not contribute to the scalar bound.

\begin{scriptsize}
\begin{verbatim}
(*
  Conductor complexity and Szpiro-style inequality (toy version).

  We keep the statement schematic. It mirrors the shape of
  Proposition~\ref{prop:NormalizedToIntensionalSzpiro}.
*)

Parameter n : EllPkg -> R.

Record EllNormalization2 := {
  NF2 : EllEndo2 ;
  nu  : forall E : EllPkg, EllHom (Fobj NF2 E) E ;

  (* one may require idempotence as a separate axiom if desired *)
}.

Definition defect (N : EllNormalization2) (E : EllPkg) : R :=
  mu (nu N E).

Parameter C_eps : R -> R.

Conjecture Intensional_Szpiro :
  forall (eps : R) (N : EllNormalization2) (E : EllPkg),
    0 < eps ->
    hDelta E <= (6 + eps) * n E + defect N E + C_eps eps.
\end{verbatim}
\end{scriptsize}

\begin{remark}[A small point of fidelity]\label{rem:CoqFidelity}
In this toy formulation, the $(2,1)$-categorical structure is literal, and the cost is invariant under $2$-isomorphism by \texttt{mu\_two\_invariant}. Thus a change of coherence data may alter the $2$-cells that witness a comparison, and yet it cannot alter the scalar defect. The cost remains attached to the transport itself.
\end{remark}

This toy subsection is limited in scope, and it offers a check rather than a claim of formalized arithmetic geometry. The grammar of Universe~A can be written, even in a small proof assistant, without erasing the $(2,1)$-categorical level. One keeps coherence explicit, and one keeps the scalarization steps separate from the coherence steps.



\section{An intensional statement of the Birch and Swinnerton-Dyer
conjecture}\label{sec:BSD}
%==================================================================

The intensional Szpiro inequality of Section~\ref{sec:IntensionalSzpiro} took the form of a transport bound, where a height was held below a conductor term together with a defect. The conjecture of Birch and Swinnerton-Dyer is of another shape. It concerns the lattice of rational points of an elliptic curve, and it relates the growth of a sequence of local point counts to the rank of that lattice, with a refined part that prescribes a leading constant. We state it within the quantitative framework of the preceding sections.

Throughout, we return to the form in which Birch and Swinnerton-Dyer first phrased their observation~\cite{Birch1965}, before the modularity theorem, so that the statement rests upon finite point counts and upon a lattice of solutions, and not upon the analytic continuation of an $L$-function. This is what permits the statement to descend, as the Szpiro inequality descended above, to a pre-set-theoretic universe.

\subsection{The Mordell--Weil package}\label{subsec:MWPackage}

\begin{definition}[Mordell--Weil package]\label{def:MordellWeilPackage}
Let $K$ be a number field. A \emph{Mordell--Weil package} over $K$ is an elliptic package $\mathcal{E} = (E, \mathcal{W}, \{\sigma_v\}_{v\in S}, \lambda)$ over $K$ in the sense of Definition~\ref{def:EllipticPackage}, enriched with the following data:
\begin{enumerate}
\item[(i)] the group of rational points $E(K)$, which by the Mordell--Weil theorem~\cite{SilvermanAEC} is finitely generated, so that one has a presentation
  \[
    E(K) \;\cong\; \Z^{r} \oplus T
  \]
with $T$ a finite torsion subgroup and $r$ the rank, given by a finite list of generators $P_1, \dots, P_r$ of the free part together with a finite presentation of $T$.
\item[(ii)] the canonical N\'eron--Tate height pairing
  \[
    \langle \cdot, \cdot \rangle \colon E(K) \times E(K) \longrightarrow \R
  \]
which descends to a positive definite form on the free part $E(K) \otimes \R \cong \R^{r}$.
\end{enumerate}
A change of generating set of the free part is a transport in the sense of the labeled signature of Definition~\ref{def:Signatures}, an element of $\mathrm{GL}_r(\Z)$ acting on the chosen basis. The height of an individual generator depends upon that choice, and so it carries a transport cost.
\end{definition}

\begin{definition}[The regulator as a transport invariant]\label{def:RegulatorAsTransport}
With the notation of Definition~\ref{def:MordellWeilPackage}, the \emph{regulator} of $\mathcal{E}$ is the Gram determinant of the canonical height pairing in a basis $P_1, \dots, P_r$ of the free part,
\[
  \Reg(\mathcal{E}) \;=\; \det\bigl( \langle P_i, P_j \rangle \bigr)_{1 \le i,j \le r}.
\]
It is the squared covolume of the Mordell--Weil lattice under the height metric. The change-of-basis transports of Definition~\ref{def:MordellWeilPackage} have determinant $\pm 1$, and so the regulator is invariant under them. Within the quantitative category of Definition~\ref{def:TransportDistance} it is therefore a zero-cost invariant of the lattice, while the heights of the individual generators are basis-dependent and bear the cost.
\end{definition}

\begin{remark}[The volume reading of the regulator]\label{rem:RegulatorVolumeReading}
The regulator is the volume of the cell of the Mordell--Weil lattice. One may read it as the volume of the transport between the lattice presented by a chosen basis and the lattice presented by the height form itself. The invariance noted above is the statement that this volume does not depend upon the presentation, even though the coordinates of the generators do. In this the regulator behaves as the short-map heights of Proposition~\ref{prop:HeightShort} behave, a quantity attached to the object and not to the manner of its display.
\end{remark}

\subsection{Local point counts as data of Universe~P}\label{subsec:PointCounts}

\begin{definition}[Point-count datum]\label{def:PointCountDatum}
Let $\mathcal{E}$ be an elliptic package over $K$ with code $d \in \mathsf{PkgCode}_K$ in the sense of Definition~\ref{def:EllPkgCode}. For each finite place $\mathfrak{p}$ of $K$ of good reduction, with residue field $\kappa(\mathfrak{p})$ of cardinality $\mathrm{N}\mathfrak{p}$, set
\[
  N_{\mathfrak{p}}(\mathcal{E}) \;=\; \# E\bigl(\kappa(\mathfrak{p})\bigr),
\]
the number of points of the reduced curve over the residue field. Each $N_{\mathfrak{p}}(\mathcal{E})$ is obtained from $d$ by counting the solutions of the reduced Weierstrass congruence together with the point at infinity, a finite and decidable procedure. The assignment $\mathfrak{p} \mapsto N_{\mathfrak{p}}(\mathcal{E})$ is therefore an object of Universe~P in the sense of Definition~\ref{def:UniverseP}, a code-indexed family of integers of the kind described in Definition~\ref{def:LogIntDatum}.
\end{definition}

\begin{definition}[Partial Euler transport]\label{def:PartialEulerTransport}
With the notation of Definition~\ref{def:PointCountDatum}, the \emph{partial Euler transport} of $\mathcal{E}$ at level $X$ is the finite product of positive rationals
\[
  \Pi_{\mathcal{E}}(X)
  \;=\;
  \prod_{\substack{\mathfrak{p}\ \text{good}\\ \mathrm{N}\mathfrak{p} \le X}}
  \frac{N_{\mathfrak{p}}(\mathcal{E})}{\mathrm{N}\mathfrak{p}}.
\]
For each $X$ this is an element of Universe~P, and the family $X \mapsto \Pi_{\mathcal{E}}(X)$ is a Universe-P sequence of rationals.
\end{definition}

\begin{remark}[The local factors and the absent series]\label{rem:EulerFactorMotivation}
Each factor $N_{\mathfrak{p}}/\mathrm{N}\mathfrak{p}$ is the reciprocal of the local Euler factor of the Hasse--Weil series at the central point, so that $\Pi_{\mathcal{E}}(X)$ is the reciprocal of a truncated Euler product. We keep this as motivation, and we do not invoke the analytic continuation of that series. The growth of $\Pi_{\mathcal{E}}(X)$ is an arithmetical matter, settled by the integers $N_{\mathfrak{p}}$ alone.
\end{remark}

\subsection{The rank form}\label{subsec:BSDRank}

\begin{conjecture}[Intensional Birch and Swinnerton-Dyer, rank
form]\label{conjecture:IntensionalBSDRank} Let $\mathcal{E}$ be a Mordell--Weil package over $K$ in the sense of Definition~\ref{def:MordellWeilPackage}, with lattice of rank $r$. Then there exists a positive real constant $C_{\mathcal{E}}$ such that
\[
  \Pi_{\mathcal{E}}(X) \;\sim\; C_{\mathcal{E}} \, (\log X)^{r}
  \qquad (X \to \infty).
\]
Equivalently, the rank of the Mordell--Weil lattice is the growth exponent of the partial Euler transport,
\[
  r \;=\; \lim_{X \to \infty} \frac{\log \Pi_{\mathcal{E}}(X)}{\log \log X}.
\]
\end{conjecture}

This is the asymptotic form in which Birch and Swinnerton-Dyer first phrased their conjecture~\cite{Birch1965}. Read within the framework, it equates two quantities of distinct provenance, the growth exponent of a sequence of local point counts on the one side and the rank of an arithmetic lattice on the other. Both are objects of Universe~P. The left member is read from the family $\mathfrak{p} \mapsto N_{\mathfrak{p}}$, and the right member is read from the lattice of rational points, and the conjecture asserts their agreement.

\begin{remark}[On the delicacy of the asymptotic]\label{rem:RankSubtlety}
The asymptotic with a finite non-zero constant is a delicate assertion. Goldfeld has shown that, once it is granted, then it carries the Riemann hypothesis for the associated series, and so it lies deeper than the equality of the two ranks alone~\cite{Goldfeld1982}. We state it in this strong asymptotic form for two reasons. It is the form that is free of analytic continuation, and it is the form in which the growth exponent, and with it the rank, appears as a limit of arithmetical quantities.
\end{remark}

\subsection{The refined form}\label{subsec:BSDRefined}

The refined conjecture prescribes the leading constant $C_{\mathcal{E}}$. We assemble its several ingredients, each as a quantity attached to the package.

\begin{definition}[Local package costs]\label{def:LocalPackageCosts}
Let $\mathcal{E}$ be a Mordell--Weil package over $K$. Its \emph{local package costs} are the finite invariants the package carries: the Tamagawa number $c_v(\mathcal{E})$ at each place $v$ of bad reduction, that is the index of the identity component in the special fiber of the N\'eron model, and the order $\# T$ of the torsion subgroup of Definition~\ref{def:MordellWeilPackage}. Each is obtained from the package code by a finite procedure -- Tate's algorithm for the $c_v$ and a torsion search bounded by the reduction data for $\# T$ -- in the manner of Tate's algorithm~\cite{SilvermanAEC, Cremona1997}.
\end{definition}

\begin{definition}[The period as a comparison cost]\label{def:PeriodAsTransport}
The \emph{period} $\Omega_{\mathcal{E}}$ of a Mordell--Weil package over $K$ is the cost attached to the comparison between the algebraic presentation of $E$ by its Weierstrass model and its analytic presentation as a complex torus $\C / \Lambda$. Explicitly, at each archimedean place $v$ it is the integral of the N\'eron differential over the local points, the covolume of the period lattice $\Lambda_v$, and $\Omega_{\mathcal{E}}$ is the product of these local contributions. As a covolume of a comparison between two presentations, the algebraic and the analytic, it is an instance of a comparison problem in the sense of Definition~\ref{def:ComparisonProblem}. Its value is a computable real number, approached by the arithmetic-geometric mean, and so it is an object of Universe~P, a Cauchy sequence of rationals rather than an appeal to set existence.
\end{definition}

\begin{remark}[The period as an evaluation in the cost quantale]\label{rem:PeriodComparisonCategory}
It is worth indicating along which path the period of Definition~\ref{def:PeriodAsTransport} is paid, so that it reads as an evaluation in the enriched language developed above rather than as an isolated real number. Fix an archimedean place $v$ of $K$. Over the curve $E$ at $v$ two presentations sit side by side. The first is the algebraic presentation $\mathrm{Alg}_v(E)$, the Weierstrass model with its N\'eron differential $\omega$, a datum of the kind carried by the package of Definition~\ref{def:EllipticPackage}. The second is the analytic presentation $\mathrm{Anal}_v(E)$, the complex torus $\C/\Lambda_v$ with its translation-invariant differential $dz$, where $\Lambda_v$ is the period lattice. Both lie over the same object of $\mathcal{C}_{\mathrm{ext}}$, and the uniformization at $v$ furnishes a vertical transport
\[
  \omega_v \colon \mathrm{Alg}_v(E) \longrightarrow \mathrm{Anal}_v(E)
\]
in the sense of Definition~\ref{def:ComparisonProblem}, the passage that reads the algebraic differential through the flat coordinate of the torus. Its cost is the logarithmic covolume by which the two normalizations of the differential disagree, oriented so that a differential of unit covolume is costless, and the local period $\Omega_v$ is the exponential of that cost. The archimedean places assemble by the tensor of the cost quantale of Definition~\ref{def:CostQuantale}, since the additive law $a+b$ on $[0,\infty]$ carries the multiplicative law $\Omega_v\,\Omega_{v'}$ of the periods, so that
\[
  \log \Omega_{\mathcal{E}} = \sum_{v \mid \infty} c(\omega_v).
\]
In the transport distance of Definition~\ref{def:TransportDistance} the period is therefore the $\mathcal{V}$-distance
\[
  d_{\mathcal{Q}}\bigl(\mathrm{Alg}_v(E),\mathrm{Anal}_v(E)\bigr) = c(\omega_v)
\]
between the two presentations, an evaluation of the hom-object of that $\mathcal{V}$-category of the same nature as the costs that the height functor of Definition~\ref{def:QuantitativeCategory}\textup{(iii)} controls. As Remark~\ref{rem:LawvereMetric} permits, the reverse passage $\mathrm{Anal}_v(E)\to\mathrm{Alg}_v(E)$ carries its own cost, and the two need not agree -- the asymmetry is the arithmetic content of the comparison rather than a defect of the account.
\end{remark}

\begin{remark}[The period across the two universes]\label{rem:PeriodAcrossUniverses}
The two readings of the period belong to two universes, and it is the cost map alone that passes between them. In Universe~A the transport $\omega_v$ of Remark~\ref{rem:PeriodComparisonCategory} is a $1$-morphism in the full sense, a comparison between the algebraic and the analytic presentation of $E$ at $v$, and its geometric content -- the flat coordinate of the torus, the integral of the N\'eron differential -- is what marks it as a comparison rather than a coincidence of two numbers. That content is never transcribed into Universe~P. What descends is the scalar cost $c(\omega_v)=\log\Omega_v$, and it descends as a fast Cauchy name of rationals in the sense of Definition~\ref{def:UniverseP}\textup{(iv)}, the arithmetic-geometric mean furnishing the modulus of convergence in the manner of Remark~\ref{rem:BSDErrorMeasures}. The certified field code of Definition~\ref{def:CertifiedFieldCode} and the package code of Definition~\ref{def:EllPkgCode} carry the period only under this name, as a coded real among the other coded reals, rather than as an analytic object.

The reading of the period as a numeric constraint variable in Universe~P is therefore the correct one, and it is the whole of what Universe~P is asked to hold. The morphism status of the period is the business of Universe~A, where it performs one task that the bare number cannot. By exhibiting $\Omega_v$ as the covolume of a comparison between presentations, it places the period within the scope of the descent obstruction of Theorem~\ref{thm:DescentObstruction}, so that the period cannot be collapsed into the regulator of Definition~\ref{def:RegulatorAsTransport} by any change of presentation. The analytic transport thus earns the distinctness of the period at the categorical level, while the arithmetic that the descent verifies sees only the fast Cauchy name. The bridge between the two levels is the same scalarization that attaches a cost to a $1$-morphism throughout the work, applied here to a transport whose geometry happens to be analytic.
\end{remark}

\begin{definition}[The Shafarevich defect]\label{def:ShafarevichDefect}
The Tate--Shafarevich group $\Sha(E/K)$ classifies the torsors under $E$ that possess a rational point over every completion $K_v$ and yet possess none over $K$. In the language of the framework it is the accumulated defect of Definition~\ref{def:AccumulatedDefect} of the descent that would globalize a family of everywhere-local solutions, a measure of the failure of local data to transport to a global point. The refined conjecture presupposes that this defect is finite, an assertion that is open in general, and under that hypothesis we write $\# \Sha(E/K)$ for its order.
\end{definition}

\begin{conjecture}[Intensional Birch and Swinnerton-Dyer, refined
form]\label{conjecture:IntensionalBSDRefined} Let $\mathcal{E}$ be a Mordell--Weil package over $K$ with finite Shafarevich defect and rank $r$. Then the leading constant $C_{\mathcal{E}}$ of Conjecture~\ref{conjecture:IntensionalBSDRank} is the total transport cost
\[
  C_{\mathcal{E}}
  \;=\;
  \kappa_r \cdot
  \frac{\# \Sha(E/K) \cdot \Omega_{\mathcal{E}} \cdot \Reg(\mathcal{E})
        \cdot \prod_{v} c_v(\mathcal{E})}
       {\bigl( \# T \bigr)^{2}}.
\]
Here $\kappa_r$ is an explicit elementary factor depending only upon the rank $r$ and the archimedean places, given by Goldfeld~\cite{Goldfeld1982}.
\end{conjecture}

\begin{remark}[Agreement with the received refined conjecture]\label{rem:RefinedClassicalForm}
The quantity
\[
  \frac{\# \Sha(E/K) \cdot \Omega_{\mathcal{E}} \cdot \Reg(\mathcal{E})
        \cdot \prod_{v} c_v(\mathcal{E})}{\bigl( \# T \bigr)^{2}}
\]
is the leading Taylor coefficient of the Hasse--Weil series at the central point in the received refined conjecture of Tate~\cite{Tate1966}. The factor $\kappa_r$ is the elementary normalization that relates that coefficient to the constant of the partial product, and it does not disturb the descent of the statement to Universe~P, since it is itself a computable real. Read in this way, the refined conjecture asserts that a constant computed from local point counts equals a total transport cost assembled from a regulator volume, a period, the local package costs, and the Shafarevich defect.
\end{remark}

\begin{remark}[The transport reading of the constant]\label{rem:TotalTransportReading}
Every member of the numerator and denominator is a transport quantity. The regulator is the covolume of the lattice transport of Definition~\ref{def:RegulatorAsTransport}. The period is the covolume of the algebraic-analytic comparison of Definition~\ref{def:PeriodAsTransport}. The Tamagawa numbers and the torsion order are the local package costs of Definition~\ref{def:LocalPackageCosts}. The Shafarevich order is the accumulated descent defect of Definition~\ref{def:ShafarevichDefect}. The refined conjecture states that the growth constant of the local counts is their product, up to the elementary factor $\kappa_r$.
\end{remark}

\subsection{Descent of the statement to Universe~P}\label{subsec:BSDDescent}

\begin{proposition}[The conjecture descends to Universe~P]\label{prop:BSDDescentToP}
Fix a number field $K$ presented by a certified field code in the sense of Definition~\ref{def:CertifiedFieldCode}, and a package code $d \in \mathsf{PkgCode}_K$ in the sense of Definition~\ref{def:EllPkgCode}. Then the rank form and the refined form, Conjectures~\ref{conjecture:IntensionalBSDRank} and~\ref{conjecture:IntensionalBSDRefined}, are each expressible as a sentence in the language of second-order arithmetic, and the assertion of each is formalizable in $\mathsf{ACA}_0$. No quantification over a totality of all elliptic curves enters, and no appeal to the analytic continuation of an $L$-function is made.
\end{proposition}

\begin{proof}
The point counts $N_{\mathfrak{p}}(d)$ are primitive recursive in $d$, being counts of solutions of a congruence, so that $\mathfrak{p} \mapsto N_{\mathfrak{p}}(d)$ and the partial Euler transport $\Pi_d(X)$ are code-indexed sequences of rationals. The constants $C_d$, $\Omega_d$, and $\kappa_r$ are computable reals, presented as $\mathsf{ACA}_0$-definable Cauchy sequences. The regulator is the determinant of an integer-indexed matrix of canonical heights, themselves computable reals. The rank $r$, the torsion order $\# T$, the Tamagawa numbers $c_v$, and the order $\# \Sha$, granted finite, are integers extracted from $d$ by arithmetical predicates. The asymptotic relation and the limit defining the growth exponent are then arithmetical statements about these sequences, of bounded quantifier complexity over the reals so coded. As in Proposition~\ref{prop:APDescentToP}, the resulting sentence speaks only of finite codes and of sequences indexed by them, and it remains within $\mathsf{ACA}_0$. The single ingredient that would carry the statement beyond arithmetic, the analytic continuation of the Hasse--Weil series that the modularity theorem provides, is absent from the asymptotic form by the choice of formulation.
\end{proof}

\begin{remark}[Error measures of the real constants and the strength of the descent]\label{rem:BSDErrorMeasures}
The descent of Proposition~\ref{prop:BSDDescentToP} asks a little more of the real constants than their bare computability, and it is worth saying what. The normalizing factor $\kappa_r$ of Remark~\ref{rem:RefinedClassicalForm}, the period $\Omega_{\mathcal{E}}$, and the leading constant $C_{\mathcal{E}}$ are each presented as a Cauchy sequence of rationals, and the assertion that $\Pi_{\mathcal{E}}(X)\sim C_{\mathcal{E}}(\log X)^{r}$ carries with it an error measure, the rate at which the partial Euler transport approaches its asymptote. Two real-analytic inputs govern that rate. The first is the equidistribution of the normalized local traces $a_{\mathfrak{p}}/2\sqrt{\mathrm{N}\mathfrak{p}}$ against the Sato--Tate measure, which controls the fluctuation of the individual factors $N_{\mathfrak{p}}/\mathrm{N}\mathfrak{p}$ of Definition~\ref{def:PartialEulerTransport}. The second is an asymptotic estimate, analogous to Mertens' second theorem, for the weighted sum of those traces, which controls the accumulation of the factors into the product.

What the descent requires of these inputs is only that their consequence for the partial product be phrased as a bound between coded reals carrying an explicit modulus of convergence. A proof of the inputs within $\mathsf{ACA}_0$ is neither claimed nor needed, the equidistribution lying deeper. Once the error measure is presented as such a modulus, the limit that defines the growth exponent of Conjecture~\ref{conjecture:IntensionalBSDRank}, and the limit that defines each real constant, are the limits of bounded monotone sequences of rationals, and $\mathsf{ACA}_0$ proves the existence of exactly these limits. The integral of the Sato--Tate measure and the Mertens constant then enter as computable reals of the same standing as $\kappa_r$, defined by arithmetical Cauchy names. In this manner the error measures of the real integration formulas are handled inside the universe by the one principle that separates $\mathsf{ACA}_0$ from its weaker companions, the convergence of bounded monotone arithmetical sequences, and the descent of Proposition~\ref{prop:BSDDescentToP} closes without an appeal to a stronger system. The standing of the modulus itself, as an existential parameter of the conjecture rather than a conditionally proven rate, is taken up in Remark~\ref{rem:ModulusExistentialParameter}.
\end{remark}

\begin{remark}[Why the modularity theorem is set aside]\label{rem:WhyAvoidModularity}
The modularity theorem identifies the Hasse--Weil series of an elliptic curve over $\Q$ with the $L$-series of a modular form, and with that identification come the analytic continuation and the functional equation. That identification is the route by which the central value and its derivatives acquire their meaning. It is also the route by which the statement of the conjecture would draw upon an analytic object whose construction reaches past arithmetic. By returning to the asymptotic form of Birch and Swinnerton-Dyer, one keeps the statement on the arithmetical side of that divide. The leading constant is then approached through a delicate limit rather than read from a value of an analytic series, and in exchange the conjecture may be stated, in full, within Universe~P, without recourse to the set-theoretic ambient of Universe~H.
\end{remark}

\begin{remark}[The obstruction keeps the ingredients
apart]\label{rem:BSDLinkToFramework} The same discipline that gave the Szpiro inequality a reading as a transport bound in Proposition~\ref{prop:NormalizedToIntensionalSzpiro} gives the refined constant a reading as a transport cost. The descent obstruction of Theorem~\ref{thm:DescentObstruction} is what holds the several invariants apart. If the regulator, the period, and the Shafarevich defect could be collapsed into one another by a change of presentation, then each would carry a vanishing relative cost, and the obstruction forbids that. They remain distinct measurements rather than artifacts of a chosen vocabulary, and the refined conjecture is the statement that their product governs the growth of the local counts.
\end{remark}
%==================================================================
\section{The Tate--Shafarevich group}\label{sec:TateShafarevich}
%==================================================================

The destination of this section is a single statement, that the Tate--Shafarevich group has finite exponent, or what is the same under a classical pairing recalled below, that its divisible part is trivial. Two of the three ingredients this asks for are unconditional. The finite-level shadows $\Sha(E/K)[n]$ are finite for every $n$, by ordinary finite descent, and every class outside the divisible part carries a positive, computable transport cost, by the non-degeneracy of the Cassels--Tate pairing. What is left, once these two facts are in hand, is a single named hypothesis, that every class of $\Sha(E/K)$ -- the divisible ones included -- carries a positive cost bounded in total by the package. We call this the Shafarevich cost bound, state it below, and close the section by showing exactly what it still asks for.

\subsection{Finite descent and the Shafarevich cost bound}
\label{subsec:ShaFinitenessCostBound}

This subsection does not prove the finiteness of $\Sha(E/K)$ in full generality. It proves the finite-level part of descent inside Universe~P, and then it isolates the exact cost-bound hypothesis under which under which an infinite Tate--Shafarevich group becomes impossible. The proof is therefore a reduction rather than a completion. Its value is that the remaining obstruction is no longer hidden in an analytic continuation, nor in an ambient set universe. It is placed in a finite transport inequality between local descent data and the Mordell--Weil package.

\begin{remark}[What is expressed in $\mathsf{ACA}_{0}$ and what is proved there]
\label{rem:ExpressibilityVersusProvability}
It is worth drawing, at the outset of the descent, a line that the rest of the section keeps. Two questions are easily run together, and they are not the same. The first is whether the finiteness of $\Sha(E/K)$ can be \emph{expressed} as a sentence over finite codes in the language of $\mathsf{ACA}_{0}$. The second is whether that sentence can be \emph{proved} in $\mathsf{ACA}_{0}$. The first holds, and the constructions of Universe~P are arranged so that it holds: the finite-level shadows, the Selmer packets, the descent costs, and the orders of the classes are all finite codes or coded reals, as Proposition~\ref{prop:ACAFormalizable} and Proposition~\ref{prop:FiniteLevelShaACA0} attest. The second does not hold in general, and the reduction does not pretend to it. What the reduction shows is that the proof needs, over and above $\mathsf{ACA}_{0}$, only a short list of classical inputs, each free of modularity and of analytic continuation, and each far weaker than the comprehension of an ambient set universe.

Two such inputs enter here. The finite generation of $E(K)$ and the termination of Tate's algorithm furnish the package, as in Definition~\ref{def:MordellWeilPackage} and Definition~\ref{def:LocalPackageCosts}. The non-degeneracy of the Cassels--Tate pairing modulo the divisible part, recalled in Lemma~\ref{lem:CasselsTateNondegeneracy}, furnishes the separation. When a step asks more than $\mathsf{ACA}_{0}$ proves, the honest ambient is the comprehension scheme $\Pi^{1}_{1}\text{-}\mathsf{CA}_{0}$, in which the Galois cohomology of a number field and the duality that the pairing rests upon are available. That scheme is itself a fragment of second-order arithmetic, far below $\mathsf{ZFC}$, so the program stays independent of the set universe of Universe~H throughout. We mark each importation where it falls, and we keep separate the sentence, which lives in $\mathsf{ACA}_{0}$, from the proof, which calls upon these few recalled principles.

A reader from reverse mathematics may wish the upper bound sharpened, and one point of the importation does admit a calibration of long standing. The formation of the maximal divisible subgroup is the point at which $\Pi^{1}_{1}$ comprehension enters, since the theorem that every countable abelian group is the direct sum of a divisible group and a reduced group is equivalent to $\Pi^{1}_{1}\text{-}\mathsf{CA}_{0}$ over $\mathsf{RCA}_{0}$~\cite{FriedmanSimpsonSmith1983,simpson2009subsystems}. The radical of the Cassels--Tate pairing being $\Sha(E/K)_{\mathrm{div}}$, the non-degeneracy of Lemma~\ref{lem:CasselsTateNondegeneracy} presupposes that subgroup as a completed set, and $\Pi^{1}_{1}\text{-}\mathsf{CA}_{0}$ is the strength at which such a subgroup is available with no help from special structure. The finite-level material stands apart. The exactness of the $n$-descent sequence at a fixed level, the Hilbert-symbol evaluations of the local invariants, and the comparison of a torsor with its split presentation concern finite Galois modules over finitely many completions, they are arithmetical in the codes, and Lemma~\ref{lem:SelmerPacketFiniteACA0} together with Proposition~\ref{prop:FiniteLevelShaACA0} carries them out in $\mathsf{ACA}_{0}$. What exceeds $\mathsf{ACA}_{0}$ is the global half of the duality, the passage from the finite levels to the cohomology of the full Galois group of $K$ and the identification of the radical of the pairing with the divisible part. For that half we claim only the upper bound $\Pi^{1}_{1}\text{-}\mathsf{CA}_{0}$, and whether it can be lowered, or whether the duality reverses to the comprehension it invokes, is to our knowledge open.
\end{remark}

\subsubsection{Finite asymptotic rank energy}
\label{subsubsec:FiniteAsymptoticRankEnergy}

\begin{definition}[Finite rank-energy window]
\label{def:FiniteRankEnergyWindow}
Let $\mathcal{E}$ be a Mordell--Weil package over $K$, and let $\Pi_{\mathcal{E}}(X)$ be the partial Euler transport of Definition~\ref{def:PartialEulerTransport}. For $X>e$, define the finite rank-energy quotient
\[
  \rho_{\mathcal{E}}(X)
  \coloneqq
  \frac{\log \Pi_{\mathcal{E}}(X)}{\log\log X}.
\]
A \emph{rank-energy modulus} for $\mathcal{E}$ is a map $M_{\mathcal{E}} \colon \Q_{>0}\to \mathbb{N}$, presented in Universe~P, such that for every rational $\epsilon>0$ and every integer $X\ge M_{\mathcal{E}}(\epsilon)$ one has
\[
  \bigl|\rho_{\mathcal{E}}(X)-r\bigr| < \epsilon.
\]
Here $r$ is the rank of the lattice in Definition~\ref{def:MordellWeilPackage}, and all logarithms are read through fast Cauchy names of rationals as in Definition~\ref{def:UniverseP}.
\end{definition}

\begin{proposition}[Rank conservation as bounded asymptotic energy]
\label{prop:RankConservationFiniteAsymptotics}
Assume Conjecture~\ref{conjecture:IntensionalBSDRank} in its modulus-bearing form. Then the equality of the asymptotic exponent with the Mordell--Weil rank is expressed in $\mathsf{ACA}_{0}$ by the existence of a rank-energy modulus $M_{\mathcal{E}}$ in the sense of Definition~\ref{def:FiniteRankEnergyWindow}. Conversely, the existence of such a modulus gives the exponent equality
\[
  r
  =
  \lim_{X\to\infty}
  \frac{\log \Pi_{\mathcal{E}}(X)}{\log\log X}
\]
inside the coded real presentation of Universe~P.
\end{proposition}

\begin{proof}
The assertion is an unwinding of the definition of a limit for coded real sequences. The sequence $X\mapsto \Pi_{\mathcal{E}}(X)$ is obtained from the finite point-count datum of Definition~\ref{def:PointCountDatum}. Thus $X\mapsto \rho_{\mathcal{E}}(X)$ is a sequence of coded reals arithmetical in the package code. In $\mathsf{ACA}_{0}$, convergence of such a sequence to the integer $r$ is given by a modulus assigning to each rational tolerance a finite threshold. This is precisely Definition~\ref{def:FiniteRankEnergyWindow}.
\end{proof}

\begin{remark}[The modulus is a parameter of the statement]
\label{rem:ModulusExistentialParameter}
The phrase \emph{modulus-bearing form} in Proposition~\ref{prop:RankConservationFiniteAsymptotics} deserves a precise gloss, since two readings of the modulus are possible and only one is intended. Classically, the existence of the limit in Conjecture~\ref{conjecture:IntensionalBSDRank} and the existence of a modulus for it are interchangeable. Inside Universe~P they are not interchangeable for free. The modulus $M_{\mathcal{E}}$ is a second-order object, and the modulus-bearing form of the conjecture is the sentence asserting, existentially, that some such object exists alongside the limit it witnesses. The quantification is purely existential and the modulus is a parameter of the statement. No specific rate of convergence is assumed anywhere in this paper. In particular, neither the rate that would accompany the Riemann hypothesis for the associated series, which the bare asymptotic carries with it by the theorem of Goldfeld recalled in Remark~\ref{rem:RankSubtlety}, nor an effective form of Sato--Tate equidistribution, enters as an input. The two real-analytic inputs named in Remark~\ref{rem:BSDErrorMeasures} are named as the provenance from which a modulus would classically be manufactured, and they are never invoked as premises. The distinction is what keeps the descent of Proposition~\ref{prop:BSDDescentToP} unconditional as a statement about statements. Were a conditionally proven rate assumed in place of the existential parameter, then the condition would carry over into Universe~P, and the paper claims no such thing.
\end{remark}

\begin{remark}[What rank conservation does not see]
\label{rem:RankConservationDoesNotSeeSha}
Proposition~\ref{prop:RankConservationFiniteAsymptotics} captures the finite asymptotic content of the rank form. It compares the logarithmic growth of the point-count product with the rank of the algebraic lattice. It does not by itself count locally soluble torsors without global points. Those torsors live in the cokernel of a descent map, while the rank-energy quotient reads only the growth exponent of the partial Euler transport. A further coupling principle is therefore needed before rank conservation can impose finiteness upon Definition~\ref{def:ShafarevichDefect}.
\end{remark}

\subsubsection{Finite descent packets}
\label{subsubsec:FiniteDescentPackets}

\begin{definition}[Certified descent torsor code]
\label{def:CertifiedDescentTorsorCode}
Fix an integer $n\ge 2$. A \emph{certified $n$-descent torsor code} for $\mathcal{E}$ is a finite tuple
\[
  \xi=(C_{\xi},\pi_{\xi},a_{\xi},S_{\xi},L_{\xi})
\]
with the following data:
\begin{enumerate}
\item[(i)] a projective genus-one model $C_{\xi}$ over $K$ given by equations with coefficients in the certified field code of Definition~\ref{def:CertifiedFieldCode},
\item[(ii)] a finite morphism $\pi_{\xi}\colon C_{\xi}\to E$ presenting $C_{\xi}$ as an $n$-covering of $E$,
\item[(iii)] an action code $a_{\xi}$ witnessing the simply transitive action of $E[n]$ on the geometric fibers of $\pi_{\xi}$,
\item[(iv)] a finite set of places $S_{\xi}$ containing the bad places of $\mathcal{E}$ and the places over $n$,
\item[(v)] finite local solubility certificates $L_{\xi}$ for the places in $S_{\xi}$, together with local lifting certificates derived from Hensel's lemma for the places outside $S_{\xi}$.
\end{enumerate}
Two such codes are declared equivalent when they present the same class in $H^{1}(K,E[n])$. The equivalence relation is arithmetical in the two codes.
\end{definition}

\begin{definition}[Intensional Selmer packet]
\label{def:IntensionalSelmerPacket}
The \emph{intensional $n$-Selmer packet} of $\mathcal{E}$ is the finite coded subset
\[
  \Sel^{\mathfrak{P}}_{n}(\mathcal{E})
  \subseteq
  \bigl\{\text{certified $n$-descent torsor codes for }\mathcal{E}\bigr\}\big/{\sim}
\]
consisting of those equivalence classes that are everywhere locally soluble by the certificates of Definition~\ref{def:CertifiedDescentTorsorCode}. The split covering coming from a rational point gives the distinguished zero class.
\end{definition}

\begin{lemma}[Finite Selmer packets in \texorpdfstring{$\mathsf{ACA}_{0}$}{ACA0}]
\label{lem:SelmerPacketFiniteACA0}
For every fixed integer $n\ge 2$ and every Mordell--Weil package $\mathcal{E}$ over $K$, the packet $\Sel^{\mathfrak{P}}_{n}(\mathcal{E})$ is finite in Universe~P.
\end{lemma}

\begin{proof}
The standard finite descent construction reduces the possible $n$-covering classes to a finite module of $S$-unit and ideal-class data, where $S$ contains the bad places of $\mathcal{E}$ and the places above $n$. In the present language this finite module is represented by a list of integer tuples bounded by the field code, the Weierstrass coefficients, and $n$. Local solubility is then checked by the finite certificates included in Definition~\ref{def:CertifiedDescentTorsorCode}. The formation of the list, the local tests, and the quotient by the arithmetical equivalence relation are finite procedures. Proposition~\ref{prop:ACAFormalizable} therefore places the construction inside $\mathsf{ACA}_{0}$.
\end{proof}

\begin{proposition}[Finite-level Shafarevich finiteness]
\label{prop:FiniteLevelShaACA0}
For every fixed integer $n\ge 2$, the $n$-torsion subgroup $\Sha(E/K)[n]$ is finite in Universe~P. More precisely, there is an injection of coded groups
\[
  \Sha(E/K)[n]
  \hookrightarrow
  \Sel^{\mathfrak{P}}_{n}(\mathcal{E}).
\]
Consequently, when there exists an integer $m\ge 2$ with $\Sha(E/K)=\Sha(E/K)[m]$, then $\Sha(E/K)$ is finite in Universe~P.
\end{proposition}

\begin{proof}
A class in $\Sha(E/K)[n]$ is represented by a torsor under $E$ that becomes split after multiplication by $n$. Such a torsor gives an $n$-covering that is locally soluble at every completion of $K$, hence it determines an element of $\Sel^{\mathfrak{P}}_{n}(\mathcal{E})$. The classes arising from global rational points form the image of $E(K)/nE(K)$, and passing to the quotient gives the displayed injection. Lemma~\ref{lem:SelmerPacketFiniteACA0} gives finiteness. The last assertion follows by taking $n=m$.
\end{proof}

\begin{remark}[The remaining gap]
\label{rem:FiniteLevelNotFullSha}
Proposition~\ref{prop:FiniteLevelShaACA0} is the part of the proof that the present finite machinery supplies without a new arithmetic principle. It proves the finiteness of each finite-level shadow of $\Sha$. It leaves open the possibility that non-zero classes occur at unbounded orders. The full finiteness problem is therefore the problem of ruling out an unbounded ascent through the Selmer packets.
\end{remark}

\subsubsection{The descent cost bound}
\label{subsubsec:DescentCostBound}

\begin{definition}[Certified Shafarevich cost]
\label{def:CertifiedShafarevichCost}
Let $\xi$ be a certified descent torsor code representing a locally soluble torsor under $E$. The \emph{certified Shafarevich cost} of $\xi$ is
\[
  \varsigma_{\mathcal{E}}(\xi)
  \coloneqq
  d_{\mathcal{Q}_{\mathrm{desc}}}
  \bigl(\mathrm{Loc}(\xi),\mathrm{Glob}(\xi)\bigr)
\]
where $\mathcal{Q}_{\mathrm{desc}}$ is the quantitative category whose objects are finite local descent atlases and global torsor presentations, and whose cost is the transport cost induced from Definition~\ref{def:ModelDefect} and Definition~\ref{def:TransportDistance}. The object $\mathrm{Loc}(\xi)$ is the finite atlas of local solution certificates. The object $\mathrm{Glob}(\xi)$ is the split global torsor presentation. The cost is represented by a fast Cauchy name of rationals in the sense of Definition~\ref{def:UniverseP}. The zero class has cost $0$. A positive lower bound for the cost of every non-zero class is not part of the definition.
\end{definition}

\begin{definition}[Finite Shafarevich load]
\label{def:FiniteShafarevichLoad}
For a finite list $\Xi=(\xi_{1},\ldots,\xi_{m})$ of pairwise inequivalent non-zero locally soluble torsor codes, define its \emph{Shafarevich load} by
\[
  \mathfrak{D}_{\mathcal{E}}(\Xi)
  \coloneqq
  \sum_{i=1}^{m}\varsigma_{\mathcal{E}}(\xi_i).
\]
This sum is a coded real of Universe~P.
\end{definition}

\begin{hypothesis}[Shafarevich speed limit]
\label{hyp:ShafarevichCostBound}
There exist coded real numbers $\eta_{\mathcal{E}}>0$ and $B_{\mathcal{E}}\ge 0$ with the following two properties:
\begin{enumerate}
\item[(i)] for every non-zero class $\xi\in\Sha(E/K)$ one has
  \[
    \varsigma_{\mathcal{E}}(\xi)\ge \eta_{\mathcal{E}}
  \]
\item[(ii)] for every finite list $\Xi$ of pairwise inequivalent non-zero classes in $\Sha(E/K)$ one has
  \[
    \mathfrak{D}_{\mathcal{E}}(\Xi)
    \le
    B_{\mathcal{E}}.
  \]
\end{enumerate}
\end{hypothesis}

\begin{proposition}[Finiteness from the cost bound]
\label{prop:CostBoundFinitenessSha}
Assume Hypothesis~\ref{hyp:ShafarevichCostBound}. Then $\Sha(E/K)$ is finite in Universe~P, and
\[
  \#\Sha(E/K)
  \le
  1+\bigl\lfloor B_{\mathcal{E}}/\eta_{\mathcal{E}}\bigr\rfloor.
\]
\end{proposition}

\begin{proof}
Let $\Xi=(\xi_{1},\ldots,\xi_{m})$ be a finite list of pairwise inequivalent non-zero classes in $\Sha(E/K)$. By Hypothesis~\ref{hyp:ShafarevichCostBound},
\[
  m\eta_{\mathcal{E}}
  \le
  \sum_{i=1}^{m}\varsigma_{\mathcal{E}}(\xi_i)
  =
  \mathfrak{D}_{\mathcal{E}}(\Xi)
  \le
  B_{\mathcal{E}}.
\]
Thus $m\le B_{\mathcal{E}}/\eta_{\mathcal{E}}$. Hence no list of pairwise inequivalent non-zero classes can have length greater than $\lfloor B_{\mathcal{E}}/\eta_{\mathcal{E}}\rfloor$. The quotient presentation of the torsor codes has an arithmetical equivalence relation. By arithmetical comprehension, one may form the predicate saying that a code is the least representative of its class. Starting from the least non-zero representative and repeating the least-search construction at most $\lfloor B_{\mathcal{E}}/\eta_{\mathcal{E}}\rfloor$ times, one obtains all non-zero classes. Were another representative left over, then adjoining its least code would produce a longer pairwise inequivalent list, contradicting the bound above. Adding the zero class gives the displayed estimate.
\end{proof}
\subsection{Separation from the Cassels--Tate pairing}
\label{subsec:TowardShafarevichCostBound}

Clause (i) of Hypothesis~\ref{hyp:ShafarevichCostBound} asks that every non-zero class of $\Sha(E/K)$ leave a visible finite defect in the descent category. This subsection proves that clause for every class outside the divisible part, drawing it from a classical pairing rather than assuming it.

\subsubsection{Integral shadows of descent cost}
\label{subsubsec:IntegralDescentShadow}

\begin{definition}[Integral descent refinement]
\label{def:IntegralDescentRefinement}
Let $\mathcal{E}$ be a Mordell--Weil package over $K$. An \emph{integral descent refinement} of the certified Shafarevich cost of Definition~\ref{def:CertifiedShafarevichCost} consists of the following data. For each certified locally soluble torsor code $\xi$, one has a finite-support map
\[
  m_{\xi}\colon \{\text{finite places of }K\}\longrightarrow \mathbb{N}
\]
arithmetical in the code of $\xi$, together with the integral descent shadow
\[
  \varsigma^{\sharp}_{\mathcal{E}}(\xi)
  \coloneqq
  \frac{1}{[K:\Q]}
  \sum_{v<\infty} m_{\xi}(v)\log \mathrm{N}v.
\]
The refinement is required to satisfy $\varsigma_{\mathcal{E}}(\xi)\ge\varsigma^{\sharp}_{\mathcal{E}}(\xi)$ for every certified locally soluble torsor code $\xi$.
\end{definition}

The integer $m_{\xi}(v)$ measures how many non-unit finite corrections are needed at $v$ in order to compare the local descent atlas of $\xi$ with the split global presentation. It is a valuation-level shadow of Definition~\ref{def:ModelDefect}. The archimedean comparisons may still be present in $\varsigma_{\mathcal{E}}(\xi)$, yet they do not enter the atom estimate below.

\begin{definition}[Descent separation]
\label{def:DescentSeparation}
An integral descent refinement satisfies \emph{descent separation} for $\mathcal{E}$ when every non-zero class $\xi\in\Sha(E/K)$ has a representative for which
\[
  \sum_{v<\infty} m_{\xi}(v) \ge 1.
\]
Equivalently, every non-zero Shafarevich class has a representative with $\varsigma^{\sharp}_{\mathcal{E}}(\xi)>0$.
\end{definition}

\begin{lemma}[The atom of a separated descent class]
\label{lem:SeparatedDescentAtom}
Assume Definition~\ref{def:IntegralDescentRefinement}, and suppose that descent separation holds in the sense of Definition~\ref{def:DescentSeparation}. Then every non-zero class $\xi\in\Sha(E/K)$ satisfies
\[
  \varsigma_{\mathcal{E}}(\xi)
  \ge
  \frac{\log 2}{[K:\Q]}.
\]
\end{lemma}

\begin{proof}
Let $\xi$ be non-zero. By Definition~\ref{def:DescentSeparation}, there is at least one finite place $v$ with $m_{\xi}(v)>0$. Since $\mathrm{N}v\ge 2$, one obtains
\[
  \varsigma^{\sharp}_{\mathcal{E}}(\xi)
  =
  \frac{1}{[K:\Q]}
  \sum_{w<\infty} m_{\xi}(w)\log \mathrm{N}w
  \ge
  \frac{\log 2}{[K:\Q]}.
\]
Definition~\ref{def:IntegralDescentRefinement} gives $\varsigma_{\mathcal{E}}(\xi)\ge\varsigma^{\sharp}_{\mathcal{E}}(\xi)$, and the claim follows.
\end{proof}

\begin{remark}[What the atom proves]
\label{rem:AtomProvesFirstClause}
Lemma~\ref{lem:SeparatedDescentAtom} proves clause \textup{(i)} of Hypothesis~\ref{hyp:ShafarevichCostBound} with $\eta_{\mathcal{E}}=\log 2/[K:\Q]$. It does so without a completed set of torsors. It uses only finite place codes, finite-support valuation profiles, and rational Cauchy names for logarithms in the sense of Definition~\ref{def:UniverseP}.
\end{remark}

\subsubsection{Separation from the Cassels--Tate pairing}
\label{subsubsec:CasselsTateSeparation}

The descent separation of Definition~\ref{def:DescentSeparation} was stated as a property one asks of a refinement, rather than as a property one holds. There is, however, a pairing of long standing upon the Tate--Shafarevich group that yields it, for every class outside the divisible part, with no appeal to modularity and none to the continuation of an $L$-function. We recall the pairing in the finite language of Universe~P, and we then read its non-degeneracy as the separation that the cost bound asks for.

\begin{definition}[The Cassels--Tate pairing on torsor codes]
\label{def:CasselsTatePairingCode}
Let $\xi$ and $\eta$ be certified locally soluble torsor codes for $\mathcal{E}$ in the sense of Definition~\ref{def:CertifiedDescentTorsorCode}. For each place $v$ of $K$, the local component
\[
  \langle \xi,\eta\rangle_{v}
  \in
  \tfrac{1}{n}\Z/\Z
\]
is the local invariant of the cup product of the two local cocycle codes at $v$, computed from the solubility certificates $L_{\xi}$ and $L_{\eta}$ by a finite Hilbert-symbol evaluation. The component is zero at every place outside the finite set $S_{\xi}\cup S_{\eta}$. The \emph{Cassels--Tate pairing} of $\xi$ and $\eta$ is the finite sum
\[
  \langle \xi,\eta\rangle_{\mathrm{CT}}
  =
  \sum_{v} \langle \xi,\eta\rangle_{v}
  \in
  \Q/\Z.
\]
The sum has finite support, so it is a coded rational of Universe~P, arithmetical in the two codes. The pairing is bilinear, and on an elliptic curve it is alternating, so that one has $\langle\xi,\xi\rangle_{\mathrm{CT}}=0$.
\end{definition}

\begin{lemma}[Non-degeneracy modulo the divisible part]
\label{lem:CasselsTateNondegeneracy}
The radical of the Cassels--Tate pairing of Definition~\ref{def:CasselsTatePairingCode} is the maximal divisible subgroup $\Sha(E/K)_{\mathrm{div}}$. That is, a class $\xi$ pairs to zero with every class of $\Sha(E/K)$ if and only if $\xi$ lies in $\Sha(E/K)_{\mathrm{div}}$. The induced pairing upon the quotient $\Sha(E/K)/\Sha(E/K)_{\mathrm{div}}$ is therefore non-degenerate.
\end{lemma}

\begin{proof}
This is the duality theorem of Cassels and Tate for the Tate--Shafarevich group, in the form given for an elliptic curve by Cassels and in general by Milne~\cite{Cassels1965, Tate1963, MilneADT}, with the alternating refinement of Poonen and Stoll~\cite{PoonenStoll1999}. The construction rests upon local and global duality in the Galois cohomology of $K$, and it does not pass through the analytic theory of the curve. We import the statement as a recalled fact, in the manner in which the finite generation of $E(K)$ and the termination of Tate's algorithm are imported in Definition~\ref{def:MordellWeilPackage} and Definition~\ref{def:LocalPackageCosts}. The strength that this importation asks of the universe is taken up in Remark~\ref{rem:ExpressibilityVersusProvability}.
\end{proof}

\begin{lemma}[A non-zero local term leaves a finite correction]
\label{lem:NonzeroLocalTermForcesCorrection}
Let $\xi$ be a certified locally soluble torsor code, and let $v$ be a finite place at which $\langle\xi,\eta\rangle_{v}\neq 0$ for some certified class $\eta$. Then the integral descent refinement of Definition~\ref{def:IntegralDescentRefinement} has $m_{\xi}(v)\ge 1$ at that place.
\end{lemma}

\begin{proof}
A non-zero local invariant at $v$ is the obstruction to comparing the local descent atlas of $\xi$ at $v$ with the split presentation by a unit transport. If one had $m_{\xi}(v)=0$, then by Definition~\ref{def:IntegralDescentRefinement} the local atlas would carry no non-unit correction at $v$, the local cocycle code would be cohomologically trivial there, and one would have $\langle\xi,\eta\rangle_{v}=0$ for every $\eta$. The non-vanishing of the local term therefore forces a non-unit finite correction, that is $m_{\xi}(v)\ge 1$.
\end{proof}

\begin{proposition}[Separation of the non-divisible classes]
\label{prop:SeparationFromCasselsTate}
There is a positive coded real $\eta_{\mathcal{E}}$ of Universe~P such that every class $\xi\in\Sha(E/K)$ outside $\Sha(E/K)_{\mathrm{div}}$ has
\[
  \varsigma_{\mathcal{E}}(\xi)
  \ge
  \eta_{\mathcal{E}}.
\]
For a class separated at a finite place one may take the finite atom $\log 2/[K:\Q]$, and for the residual classes one may take the least cost $\eta^{\infty}_{\mathcal{E}}>0$ of a non-trivial archimedean comparison of the package, so that $\eta_{\mathcal{E}}=\min\{\log 2/[K:\Q],\,\eta^{\infty}_{\mathcal{E}}\}$.
\end{proposition}

\begin{proof}
Let $\xi\notin\Sha(E/K)_{\mathrm{div}}$. By Lemma~\ref{lem:CasselsTateNondegeneracy} there is a certified class $\eta$ with $\langle\xi,\eta\rangle_{\mathrm{CT}}\neq 0$. The pairing is the finite sum of its local components of Definition~\ref{def:CasselsTatePairingCode}, so at least one local term $\langle\xi,\eta\rangle_{v}$ is non-zero. If that place $v$ is finite, then Lemma~\ref{lem:NonzeroLocalTermForcesCorrection} gives $m_{\xi}(v)\ge 1$, so $\sum_{w<\infty}m_{\xi}(w)\ge 1$, which is descent separation in the sense of Definition~\ref{def:DescentSeparation}, and Lemma~\ref{lem:SeparatedDescentAtom} then gives $\varsigma_{\mathcal{E}}(\xi)\ge\log 2/[K:\Q]$. If instead the only non-zero local terms sit at archimedean places, then a non-trivial real local invariant is the obstruction to a unit archimedean comparison, so the archimedean part of the certified cost of Definition~\ref{def:CertifiedShafarevichCost} is at least the package constant $\eta^{\infty}_{\mathcal{E}}>0$. In either case the displayed bound holds.
\end{proof}

\begin{remark}[Separation is now held, not asked for]
\label{rem:SeparationPossessed}
With Proposition~\ref{prop:SeparationFromCasselsTate}, the descent separation of Definition~\ref{def:DescentSeparation} no longer belongs to the properties one assumes, for every class outside the divisible part. Clause~\textup{(i)} of Hypothesis~\ref{hyp:ShafarevichCostBound} then holds upon the non-divisible quotient, and it holds upon the whole of $\Sha(E/K)$ exactly when the divisible part is trivial. The divisible classes pair to zero with everything, so the Cassels--Tate pairing is silent upon them, and their separation cannot be drawn from this source. One reads the residual archimedean case of the proposition the same way. A class whose Cassels--Tate non-triviality is carried at the archimedean places alone takes values in the $2$-torsion of $\Q/\Z$, so its image in $\Sha(E/K)/\Sha(E/K)_{\mathrm{div}}$ is killed by $2$. Such images form a subgroup of $(\Sha(E/K)/\Sha(E/K)_{\mathrm{div}})[2]$, which is finite by Proposition~\ref{prop:FiniteLevelShaACA0}, so they add at most a package-bounded constant to any load and do not bear upon finiteness. The single statement that remains is therefore the triviality of the divisible part, which is bounded exponent for $\Sha(E/K)$.
\end{remark}

\subsubsection{The divisible part as the sole obstruction}
\label{subsubsec:DivisiblePartObstruction}

The reading of Remark~\ref{rem:SeparationPossessed} places the whole weight of finiteness upon one quantity. We name it here, and we state in what sense the rest of the section is a search for a bound upon it.

\begin{definition}[The divisible part of the Shafarevich group]
\label{def:DivisibleSubgroupSha}
The \emph{divisible part} of $\Sha(E/K)$ is the maximal divisible subgroup
\[
  \Sha(E/K)_{\mathrm{div}}
  =
  \bigcap_{n\ge 1} n\,\Sha(E/K)
\]
which is the intersection of the images of multiplication by $n$ over all positive integers $n$. A group has finite exponent when there is an integer $n_{0}\ge 2$ that kills every element, that is when $\Sha(E/K)=\Sha(E/K)[n_{0}]$.
\end{definition}

\begin{proposition}[The divisible part is the sole obstruction]
\label{prop:DivisiblePartSoleObstruction}
In Universe~P the following hold for $\Sha(E/K)$.
\begin{enumerate}
\item[(i)] The group $\Sha(E/K)$ is finite if and only if it has finite exponent.
\item[(ii)] The group $\Sha(E/K)$ has finite exponent if and only if its divisible part $\Sha(E/K)_{\mathrm{div}}$ is trivial and the quotient $\Sha(E/K)/\Sha(E/K)_{\mathrm{div}}$ has finite exponent.
\end{enumerate}
Under the Cassels--Tate pairing of Lemma~\ref{lem:CasselsTateNondegeneracy} the quotient $\Sha(E/K)/\Sha(E/K)_{\mathrm{div}}$ carries a non-degenerate alternating pairing, so that, once finite, it has square order.
\end{proposition}

\begin{proof}
For \textup{(i)}, if $\Sha(E/K)$ is finite, then its exponent is finite. Conversely, if $\Sha(E/K)=\Sha(E/K)[n_{0}]$ for some $n_{0}\ge 2$, then Proposition~\ref{prop:FiniteLevelShaACA0} gives finiteness in Universe~P, since $\Sha(E/K)[n_{0}]$ injects into the finite packet $\Sel^{\mathfrak{P}}_{n_{0}}(\mathcal{E})$. For \textup{(ii)}, a group of finite exponent has no non-zero divisible element, since a divisible element is divisible by every $n$ while a bounded exponent kills it, so the divisible part is trivial and the quotient has the same finite exponent. Conversely, if the divisible part is trivial, then the group equals its quotient by that part, and a finite exponent of the quotient is a finite exponent of the group. The pairing statement is Lemma~\ref{lem:CasselsTateNondegeneracy} together with the alternating property of Definition~\ref{def:CasselsTatePairingCode}, a non-degenerate alternating form upon a finite abelian group forcing square order.
\end{proof}

\begin{remark}[What is left to search for]
\label{rem:TowerSearchesBoundedExponent}
Proposition~\ref{prop:DivisiblePartSoleObstruction} turns the finiteness of $\Sha(E/K)$ into a single arithmetical demand, that the group have finite exponent, or what is the same under the pairing, that its divisible part be trivial. The finite-level finiteness of Proposition~\ref{prop:FiniteLevelShaACA0} and the separation of Proposition~\ref{prop:SeparationFromCasselsTate} are then in hand without further hypothesis, and they settle the demand entirely for every class outside $\Sha(E/K)_{\mathrm{div}}$. What is left is the divisible part itself, upon which the pairing is silent by construction. Hypothesis~\ref{hyp:ShafarevichCostBound} names the most direct route to a positive answer, a single uniform cost bound covering the divisible classes along with the rest, and Proposition~\ref{prop:CostBoundFinitenessSha} shows that this one hypothesis closes the proof outright. The remainder of this section says what can, and cannot, yet be said about it.
\end{remark}

\begin{proposition}[Equivalence of the cost bound with finiteness]
\label{prop:CostBoundEquivalentFiniteness}
Grant the classical inputs of this section, namely the non-degeneracy of Lemma~\ref{lem:CasselsTateNondegeneracy} together with the integral descent refinement of Definition~\ref{def:IntegralDescentRefinement} and the separation it yields in Proposition~\ref{prop:SeparationFromCasselsTate}, and read the cost of a class through its least certified representative, as in the proof of Proposition~\ref{prop:CostBoundFinitenessSha}. Then Hypothesis~\ref{hyp:ShafarevichCostBound} holds if and only if $\Sha(E/K)$ is finite in Universe~P.
\end{proposition}

\begin{proof}
One direction is Proposition~\ref{prop:CostBoundFinitenessSha}. For the converse, suppose that $\Sha(E/K)$ is finite. A finite abelian group has a finite exponent $m$, and a divisible group $D$ satisfies $D=mD$, so the divisible part, being at once finite and divisible, satisfies $\Sha(E/K)_{\mathrm{div}}=m\,\Sha(E/K)_{\mathrm{div}}=0$. Every non-zero class therefore lies outside the divisible part, and Proposition~\ref{prop:SeparationFromCasselsTate} supplies a positive coded real $\eta_{\mathcal{E}}$ with $\varsigma_{\mathcal{E}}(\xi)\ge\eta_{\mathcal{E}}$ upon each of them, which is clause~\textup{(i)}. For clause~\textup{(ii)}, the non-zero classes are enumerated by a finite list of certified codes, each taken least in its class by the arithmetical least-search of the proof of Proposition~\ref{prop:CostBoundFinitenessSha}, and one sets $B_{\mathcal{E}}$ to be the sum of their certified costs, a finite sum of coded reals and hence a coded real of Universe~P. Every finite list $\Xi$ of pairwise inequivalent non-zero classes is then a sublist of that full list, and every summand of a load is non-negative, so $\mathfrak{D}_{\mathcal{E}}(\Xi)\le B_{\mathcal{E}}$, which is clause~\textup{(ii)}.
\end{proof}

\begin{remark}[The cost bound is a reformulation rather than a weakening]
\label{rem:CostBoundReformulation}
Proposition~\ref{prop:CostBoundEquivalentFiniteness} should temper any impression that Hypothesis~\ref{hyp:ShafarevichCostBound} is a milder demand than the finiteness it implies. Modulo the classical inputs named there the two statements coincide, and the hypothesis is the finiteness conjecture rewritten in the transport vocabulary of Universe~P. The rewriting neither strengthens nor weakens what is asked. It places the conjecture in the shape of a finite quantitative inequality, and it exposes clause~\textup{(i)} as the part the Cassels--Tate pairing settles outside the divisible part. The equivalence reads the cost of a class through its least certified representative. Under a reading in which the representative varies freely, clause~\textup{(ii)} asks in addition a bound uniform over representatives, a strengthening upon which no argument of this paper draws.
\end{remark}

\begin{remark}[How the cost bound meets the divisible part]
\label{rem:CostBoundMeetsDivisibility}
It is worth being exact about the mechanism by which Hypothesis~\ref{hyp:ShafarevichCostBound} excludes a non-trivial divisible part, since two mechanisms could be imagined and only one is at work. Suppose $\Sha(E/K)_{\mathrm{div}}\neq 0$. The group $\Sha(E/K)$ is torsion, a non-zero divisible torsion abelian group contains a copy of the Pr\"ufer group $\Q_{\ell}/\Z_{\ell}$ for some prime $\ell$, and such a copy supplies finite lists of pairwise inequivalent non-zero classes of every length. Clause~\textup{(i)} prices each of them at $\eta_{\mathcal{E}}$ or more, so the loads of ever longer lists surpass every bound, against clause~\textup{(ii)}. The exclusion is therefore an exclusion by counting. The hypothesis forbids an infinitude of non-zero classes, divisibility entails that infinitude, and the divisible part falls with it.

A second mechanism plays no part here. The cost $\varsigma_{\mathcal{E}}$ of Definition~\ref{def:CertifiedShafarevichCost} compares one torsor code with the split presentation, and its axioms do not couple the cost of a class to the costs of its $n$-th roots. A chain $\xi=n\,\xi_{n}$, available for every $n$ when $\xi$ is divisible, therefore exerts no downward pressure upon $\varsigma_{\mathcal{E}}(\xi)$. The metric is neutral upon the existence of divisible classes, a neutrality of a piece with Remark~\ref{rem:SeparationPossessed}, where the one pairing that separates classes was seen to be silent precisely upon the divisible ones. The exclusion is effected by the finitude the hypothesis demands, with the cost bound acting as a constraint imposed upon the metric rather than read off from it.
\end{remark}
%------------------------------------------------------------------
\subsection{What remains}
\label{subsec:BSDWhatRemains}
%------------------------------------------------------------------

The two results above delimit the reduction precisely. Proposition~\ref{prop:SeparationFromCasselsTate} separates every class outside the divisible part, and Proposition~\ref{prop:DivisiblePartSoleObstruction} shows that this separation is exactly the content finiteness asks for. The divisible part is therefore the whole of what is left, and it resists the local methods that serve elsewhere in this paper.

A class of $\Sha(E/K)$ is locally soluble at every completion of $K$. Once a local point is chosen at a place $v$ the torsor becomes isomorphic to $E$ over $K_v$, so the Kodaira type, the conductor exponent, and the Tamagawa number, each computed from the local model by Tate's algorithm and Ogg's formula, agree across distinct classes. A predicate built from local reduction data alone therefore cannot separate a non-zero class from zero. What distinguishes them is global and cohomological, which is where Proposition~\ref{prop:SeparationFromCasselsTate} stops.

The reduction thus isolates Hypothesis~\ref{hyp:ShafarevichCostBound} as its single remaining object, a uniform positive cost on every non-zero class together with a uniform bound on the total cost of any finite family of them. It holds on the non-divisible quotient, and by Proposition~\ref{prop:CostBoundEquivalentFiniteness}, modulo the classical inputs of this section, it is equivalent to the finiteness of $\Sha(E/K)$.


\bibliographystyle{smfalpha}
\bibliography{intensional_semantics}

\end{document}